% Sphères — Mathématiques 1ère E — Cours signé OnBuch+
% Programme officiel MINESEC. Compiler avec : tectonic N20-spheres.tex
\newcommand{\DOCMATIERE}{Mathématiques}
\newcommand{\DOCNIVEAU}{1ère E}
\newcommand{\DOCMODULE}{Géométrie dans l'espace}
\newcommand{\DOCLECON}{N20}
\newcommand{\DOCTITRE}{Sphères}
% OnBuch+ — préambule « cours signés » ENRICHI (compilé avec tectonic / XeLaTeX)
% Système visuel « Pop » d'OnBuch+ : fond crème (jamais blanc), bordures encre
% épaisses, ombres « dures » décalées, titres Archivo Black en capitales.
% Version MATHÉMATIQUES : graphes (pgfplots/tikz), boîtes pédagogiques
% (définition, propriété, méthode, attention, le savais-tu, exercice résolu,
% exercice, TP, bilan), page de garde et signature OnBuch+ ancrée en bas.
%
% Macros attendues dans main.tex AVANT \input{preamble} :
%   \DOCMATIERE  \DOCNIVEAU  \DOCMODULE  \DOCLECON (numéro)  \DOCTITRE
\documentclass[11pt]{article}

\usepackage{fontspec}
\usepackage[french]{babel}
\usepackage[a4paper,margin=2cm,top=3.4cm,bottom=2.8cm,headheight=1.4cm,footskip=1.3cm]{geometry}
\usepackage{amsmath,amssymb,amsfonts}
\usepackage{mathtools} % \xmapsto, \coloneqq... souvent utilisés par les modèles
\usepackage{cancel} % simplifications barrées
% \abs{z} (module, valeur absolue) et \norm{u} (norme), utilisés par les modèles
\providecommand{\abs}{}\let\abs\relax\DeclarePairedDelimiter{\abs}{\lvert}{\rvert}
\providecommand{\norm}{}\let\norm\relax\DeclarePairedDelimiter{\norm}{\lVert}{\rVert}
\usepackage[table]{xcolor} % [table] : fournit aussi \rowcolors (alternance de lignes)
\usepackage{pagecolor}
\usepackage{tikz}
\usetikzlibrary{arrows.meta,positioning,calc,shapes.geometric,shapes.misc,shapes.arrows,decorations.pathmorphing,decorations.pathreplacing,decorations.markings,patterns,shadows,mindmap,trees,fit,backgrounds,matrix,intersections,angles,quotes}
\usepackage{pgfplots}
\usepackage{pgf-pie} % diagrammes circulaires \pie (statistiques)
\pgfplotsset{compat=1.18}
\usepgfplotslibrary{fillbetween,groupplots} % aires entre courbes (calcul intégral)
\usepackage{enumitem}
\usepackage{fancyhdr}
% [most] charge toutes les bibliothèques tcolorbox (listings, minted, raster,
% documentation...) et gonfle inutilement la mémoire TeX sur un document long
% (~300 boîtes) : on ne charge que ce qui est réellement utilisé.
\usepackage[skins,breakable]{tcolorbox}
\usepackage{graphicx}
\usepackage{colortbl}
\usepackage{booktabs}
\usepackage{tabularx}
\usepackage{diagbox} % case d'en-tête diagonale des tableaux à double entrée
% Colonnes extensibles centrée / gauche / droite (C, L, R), souvent utilisées par les modèles
\newcolumntype{C}{>{\centering\arraybackslash}X}
\newcolumntype{L}{>{\raggedright\arraybackslash}X}
\newcolumntype{R}{>{\raggedleft\arraybackslash}X}
\usepackage{array}
\usepackage{multicol}
\usepackage{siunitx}
% parse-numbers=false : accepte aussi \SI{2,5 \times 10^{-3}}{...}
\sisetup{locale=FR,output-decimal-marker={,},group-minimum-digits=4,parse-numbers=false}
\DeclareSIUnit\ohm{\ensuremath{\symup{\Omega}}} % U+2126 absent de la police
\usepackage{qrcode}
% \degree : commande fréquemment utilisée par les modèles pour un angle ou une
% température en dehors de \SI{}{} (non fournie par les paquets chargés).
\providecommand{\degree}{^{\circ}}
% \qed (fin de démonstration), utilisé par les modèles sans amsthm
\providecommand{\qed}{\hfill$\square$}
% \barre : double barre des valeurs interdites dans un tableau de variations (tkz-tab)
\providecommand{\barre}{\ensuremath{\|}}
% environnement proof (amsthm non chargé) utilisé par les modèles
\ifdefined\proof\else\newenvironment{proof}[1][Démonstration]{\par\noindent\textit{#1.}\ }{\qed\par}\fi
\usepackage{lastpage}
\usepackage{hyperref}
\hypersetup{hidelinks}

% ============================================================
% Palette « Pop » OnBuch+ — mêmes hex que la classe P (plus_ui.dart)
% ============================================================
\definecolor{poporange}{HTML}{F59321}
\definecolor{popdark}{HTML}{E07A0C}
\definecolor{popcream}{HTML}{FFF6E8}
\definecolor{popink}{HTML}{1C1714}
\definecolor{poppurple}{HTML}{7A5AE0}
\definecolor{popgreen}{HTML}{2E9E6A}
\definecolor{poppink}{HTML}{E0457B}
\definecolor{popblue}{HTML}{2F7DE1}
\definecolor{popgold}{HTML}{C98A12}
\definecolor{popmuted}{HTML}{8A7F76}
\definecolor{popline}{HTML}{EADFD2}
\definecolor{popgray}{HTML}{8A7F76}
\colorlet{popgrey}{popgray}
% Alias tolérants (le modèle nomme parfois les couleurs différemment)
\colorlet{popwhite}{white}
\colorlet{popblack}{popink}
\colorlet{popred}{poppink}
\colorlet{popyellow}{popgold}
% Teintes claires pour fonds de tableaux / zones
\colorlet{popblueL}{popblue!10!white}
\colorlet{poporangeL}{poporange!14!white}
\colorlet{popgreenL}{popgreen!12!white}
\colorlet{poppurpleL}{poppurple!12!white}
\colorlet{poppinkL}{poppink!12!white}
\colorlet{popgoldL}{popgold!14!white}

\pagecolor{popcream}

% ============================================================
% Typographie
% ============================================================
\defaultfontfeatures{Ligatures=TeX}
\setmainfont{PlusJakartaSans-Medium.ttf}[Path=../../../fonts/,BoldFont=PlusJakartaSans-Bold.ttf]
% Maths en sans-serif (Fira Math) avec chiffres et lettres droites de Plus
% Jakarta, pour que les formules se fondent dans le texte.
\usepackage[math-style=ISO,bold-style=ISO]{unicode-math}
\setmathfont{FiraMath-Regular.otf}[Path=../../../fonts/]
\setmathfont{PlusJakartaSans-Medium.ttf}[Path=../../../fonts/,range={up/num,up/{Latin,latin},\mathup}]
\usepackage{icomma} % virgule décimale sans espace en mode math
% Caractères mathématiques Unicode absents des polices de texte (Plus Jakarta,
% Archivo Black) : rendus par leur équivalent mathématique, en texte comme en
% formule — sinon ils s'affichent en carré vide (ex. « Arithmétique dans ℤ »).
\usepackage{newunicodechar}
\newunicodechar{ℝ}{\ensuremath{\mathbb{R}}}
\newunicodechar{ℤ}{\ensuremath{\mathbb{Z}}}
\newunicodechar{ℂ}{\ensuremath{\mathbb{C}}}
\newunicodechar{ℕ}{\ensuremath{\mathbb{N}}}
\newunicodechar{ℚ}{\ensuremath{\mathbb{Q}}}
\newunicodechar{𝔻}{\ensuremath{\mathbb{D}}}
\newunicodechar{α}{\ensuremath{\alpha}}
\newunicodechar{β}{\ensuremath{\beta}}
\newunicodechar{γ}{\ensuremath{\gamma}}
\newunicodechar{δ}{\ensuremath{\delta}}
\newunicodechar{ε}{\ensuremath{\varepsilon}}
\newunicodechar{θ}{\ensuremath{\theta}}
\newunicodechar{λ}{\ensuremath{\lambda}}
\newunicodechar{μ}{\ensuremath{\mu}}
\newunicodechar{σ}{\ensuremath{\sigma}}
\newunicodechar{τ}{\ensuremath{\tau}}
\newunicodechar{φ}{\ensuremath{\varphi}}
\newunicodechar{ψ}{\ensuremath{\psi}}
\newunicodechar{ω}{\ensuremath{\omega}}
\newunicodechar{Σ}{\ensuremath{\Sigma}}
% majuscules produites par \MakeUppercase dans les titres (α-aminés -> Α)
\newunicodechar{Α}{\ensuremath{\alpha}}
\newunicodechar{Β}{\ensuremath{\beta}}
\newunicodechar{Γ}{\ensuremath{\Gamma}}
\newunicodechar{Δ}{\ensuremath{\Delta}}
\newunicodechar{Ε}{\ensuremath{\varepsilon}}
\newunicodechar{Θ}{\ensuremath{\Theta}}
\newunicodechar{Λ}{\ensuremath{\Lambda}}
\newunicodechar{Μ}{\ensuremath{\mu}}
\newunicodechar{Τ}{\ensuremath{\tau}}
\newunicodechar{Φ}{\ensuremath{\Phi}}
\newunicodechar{Ψ}{\ensuremath{\Psi}}
\newunicodechar{Ω}{\ensuremath{\Omega}}
\newunicodechar{↦}{\ensuremath{\mapsto}}
\newunicodechar{∈}{\ensuremath{\in}}
\newunicodechar{∉}{\ensuremath{\notin}}
\newunicodechar{∧}{\ensuremath{\wedge}}
\newunicodechar{⇔}{\ensuremath{\Leftrightarrow}}
\newunicodechar{⟺}{\ensuremath{\Longleftrightarrow}}
\newunicodechar{⇒}{\ensuremath{\Rightarrow}}
\newunicodechar{⟹}{\ensuremath{\Longrightarrow}}
\newunicodechar{∘}{\ensuremath{\circ}}
\newunicodechar{∪}{\ensuremath{\cup}}
\newunicodechar{∩}{\ensuremath{\cap}}
\newunicodechar{≡}{\ensuremath{\equiv}}
\newunicodechar{⊂}{\ensuremath{\subset}}
\newunicodechar{⊥}{\ensuremath{\perp}}
\newunicodechar{∃}{\ensuremath{\exists}}
\newunicodechar{∀}{\ensuremath{\forall}}
\newunicodechar{∛}{\ensuremath{\sqrt[3]{\,}}}
\newunicodechar{∜}{\ensuremath{\sqrt[4]{\,}}}
\newunicodechar{⃗}{\ensuremath{{}^{\to}}}
\newunicodechar{ⁿ}{\ifmmode^{n}\else\textsuperscript{\ensuremath{n}}\fi}
\newunicodechar{ˣ}{\ifmmode^{x}\else\textsuperscript{\ensuremath{x}}\fi}
\newunicodechar{ᵇ}{\ifmmode^{b}\else\textsuperscript{\ensuremath{b}}\fi}
\newunicodechar{ʳ}{\ifmmode^{r}\else\textsuperscript{\ensuremath{r}}\fi}
\newunicodechar{ᵘ}{\ifmmode^{u}\else\textsuperscript{\ensuremath{u}}\fi}
\newunicodechar{ʸ}{\ifmmode^{y}\else\textsuperscript{\ensuremath{y}}\fi}
\newunicodechar{ᵃ}{\ifmmode^{a}\else\textsuperscript{\ensuremath{a}}\fi}
\newunicodechar{ᵉ}{\ifmmode^{e}\else\textsuperscript{\ensuremath{e}}\fi}
\newunicodechar{ᵏ}{\ifmmode^{k}\else\textsuperscript{\ensuremath{k}}\fi}
\newunicodechar{ᵖ}{\ifmmode^{p}\else\textsuperscript{\ensuremath{p}}\fi}
\newunicodechar{ᴺ}{\ifmmode^{N}\else\textsuperscript{\ensuremath{N}}\fi}
\newunicodechar{ᶜ}{\ifmmode^{c}\else\textsuperscript{\ensuremath{c}}\fi}
\newunicodechar{ʰ}{\ifmmode^{h}\else\textsuperscript{\ensuremath{h}}\fi}
\newunicodechar{ᵠ}{\ifmmode^{q}\else\textsuperscript{\ensuremath{q}}\fi}
\newunicodechar{ₙ}{\ifmmode_{n}\else\textsubscript{\ensuremath{n}}\fi}
\newunicodechar{ₐ}{\ifmmode_{a}\else\textsubscript{\ensuremath{a}}\fi}
\newunicodechar{ᵢ}{\ifmmode_{i}\else\textsubscript{\ensuremath{i}}\fi}
\newunicodechar{ᵣ}{\ifmmode_{r}\else\textsubscript{\ensuremath{r}}\fi}
\newunicodechar{ₖ}{\ifmmode_{k}\else\textsubscript{\ensuremath{k}}\fi}
\newunicodechar{ᵦ}{\ifmmode_{\beta}\else\textsubscript{\ensuremath{\beta}}\fi}
\newunicodechar{ₓ}{\ifmmode_{x}\else\textsubscript{\ensuremath{x}}\fi}
\newfontfamily\popdisplay{ArchivoBlack-Regular.ttf}[Path=../../../fonts/]
\newfontfamily\poplogo{SpaceGrotesk-Bold.ttf}[Path=../../../fonts/]
\newfontfamily\popsemi{PlusJakartaSans-SemiBold.ttf}[Path=../../../fonts/]
\newfontfamily\popxbold{PlusJakartaSans-ExtraBold.ttf}[Path=../../../fonts/]
\newfontfamily\popxboldls{PlusJakartaSans-ExtraBold.ttf}[Path=../../../fonts/,LetterSpace=3.0]

\setlist[enumerate]{leftmargin=1.7em,itemsep=5pt,topsep=4pt}
\setlist[itemize]{leftmargin=1.7em,itemsep=4pt,topsep=4pt,label={\color{poporange}\textbullet}}
\setlength{\parindent}{0pt}
\setlength{\parskip}{5pt}
\renewcommand{\arraystretch}{1.25}

% pgfplots : style Pop par défaut
\pgfplotsset{
  every axis/.append style={axis line style={popink,line width=1pt},tick style={popink},
    grid style={popline,line width=0.5pt},label style={font=\small},tick label style={font=\footnotesize}},
  popcurve/.style={poporange,line width=1.8pt,no markers,smooth},
}
% Même style disponible dans un \draw TikZ simple (hors pgfplots)
\tikzset{popcurve/.style={poporange,line width=1.8pt}}

% ============================================================
% Pill / squiggle
% ============================================================
\newcommand{\poppill}[3]{%
  \tikz[baseline=-0.55ex]\node[draw=popink,line width=1.2pt,rounded corners=2.6pt,%
    fill=#2,inner xsep=6.5pt,inner ysep=3pt]{%
    \popxboldls\fontsize{7}{7}\selectfont\color{#3}\MakeUppercase{#1}};%
}
\newcommand{\popsquiggle}[1]{%
  \tikz[baseline=-0.4ex]{
    \draw[#1,line width=2.6pt,line cap=round]
      plot [smooth] coordinates {(0,0.09) (0.34,0.01) (0.68,0.21) (1.02,0.01) (1.36,0.10)};
  }%
}
% Mot-clé surligné (marqueur orange sous le texte)
\newcommand{\cle}[1]{{\popxbold\color{popdark}#1}}

% ============================================================
% En-tête / pied
% ============================================================
\pagestyle{fancy}
\fancyhf{}
\renewcommand{\headrulewidth}{0pt}
\renewcommand{\footrulewidth}{0pt}
\lhead{\raisebox{-0.15ex}{\poplogo\fontsize{13}{13}\selectfont\color{popink}OnBuch\color{poporange}+}}
\rhead{\poppill{\DOCMATIERE\ \textbullet\ \DOCNIVEAU\ \textbullet\ Leçon \DOCLECON}{white}{popink}}
\lfoot{\popxbold\fontsize{7.6}{7.6}\selectfont\color{popmuted}\MakeUppercase{Cours signé OnBuch+}\ \textbullet\ {\popsemi\DOCTITRE}}
\rfoot{\tikz[baseline=-0.6ex]\node[rounded corners=8.5pt,fill=popink,minimum height=17pt,minimum width=17pt,inner xsep=5pt,inner sep=0]{\popxbold\fontsize{8}{8}\selectfont\color{popcream}\thepage\,/\,\pageref*{LastPage}};}

% ============================================================
% Titres
% ============================================================
\newcommand{\chaptitre}[2]{%
  \begin{tcolorbox}[enhanced,colback=poporange,colframe=popink,boxrule=2.6pt,arc=11pt,
      drop shadow=popink,left=15pt,right=15pt,top=12pt,bottom=11pt,before skip=3pt,after skip=18pt]
    \poppill{#2}{popink}{popcream}\par\vspace{9pt}
    {\raggedright\hyphenpenalty=10000\exhyphenpenalty=10000\popdisplay\fontsize{22}{24}\selectfont\color{popcream}\MakeUppercase{#1}\par}
  \end{tcolorbox}
}
% Section numérotée : pastille orange avec le numéro + titre Archivo
\newcounter{popsec}
\newcommand{\coursec}[1]{%
  \stepcounter{popsec}\setcounter{popsub}{0}%
  \par\vspace{16pt}\noindent
  \begin{minipage}{\linewidth}
  \tikz[baseline=-0.7ex]\node[draw=popink,line width=1.6pt,fill=poporange,rounded corners=4pt,minimum size=22pt,inner sep=0]{\popdisplay\fontsize{12}{12}\selectfont\color{popcream}\thepopsec};%
  \hspace{8pt}{\popdisplay\fontsize{15}{17}\selectfont\color{popink}\MakeUppercase{#1}}\par
  \vspace{4pt}\noindent{\color{poporange}\rule{36pt}{3.6pt}}
  \end{minipage}\par\nopagebreak\vspace{8pt}
}
\newcounter{popsub}
\newcommand{\courssub}[1]{%
  \stepcounter{popsub}%
  \par\vspace{9pt}\noindent{\popxbold\fontsize{11.5}{13}\selectfont\color{popink}{\color{poporange}\thepopsec.\thepopsub}\hspace{6pt}#1}\par\nopagebreak\vspace{4pt}
}

% ============================================================
% Boîtes pédagogiques — même recette : fond blanc, bordure épaisse colorée,
% ombre dure encre, badge plein. Titre optionnel [#1].
% ============================================================
\tcbset{popbase/.style={breakable,enhanced,colback=white,boxrule=2.3pt,arc=9pt,
  shadow={3pt}{-3pt}{0pt}{popink},left=14pt,right=14pt,top=11pt,bottom=11pt,before skip=11pt,after skip=15pt,
  fontupper=\popsemi\fontsize{10.6}{14.5}\selectfont\color{popink},
  fontlower=\popsemi\fontsize{10.6}{14.5}\selectfont\color{popink}}}
% Coche dessinée (le glyphe ✓ n'existe pas dans les polices du document)
\newcommand{\popcheck}[1]{\tikz[baseline=-0.5ex]\draw[#1,line width=1.6pt,line cap=round,line join=round](-0.13,0.0)--(-0.04,-0.09)--(0.13,0.1);}
\newcommand{\popboxhead}[3]{\poppill{#1}{#2}{white}\ifx\relax#3\relax\else\hspace{6pt}{\popxbold\fontsize{10.6}{12}\selectfont\color{popink}#3}\fi\par\vspace{7pt}}

\newtcolorbox{aretenir}[1][]{popbase,colframe=popgreen,before upper={\popboxhead{À retenir}{popgreen}{#1}}}
\newtcolorbox{exemplebox}[1][]{popbase,colframe=poporange,before upper={\popboxhead{Exemple}{poporange}{#1}}}
\newtcolorbox{definition}[1][]{popbase,colframe=popblue,before upper={\popboxhead{Définition}{popblue}{#1}}}
\newtcolorbox{propriete}[1][]{popbase,colframe=poppurple,before upper={\popboxhead{Propriété}{poppurple}{#1}}}
\newtcolorbox{methode}[1][]{popbase,colframe=popgold,before upper={\popboxhead{Méthode}{popgold}{#1}}}
\newtcolorbox{attention}[1][]{popbase,colframe=poppink,before upper={\popboxhead{Attention}{poppink}{#1}}}
\newtcolorbox{experience}[1][]{popbase,colframe=popblue,colback=popblueL,before upper={\popboxhead{Expérience}{popblue}{#1}}}
% « Le savais-tu ? » : carte violette pleine (contexte, culture, Cameroun)
\newtcolorbox{savaistu}[1][]{popbase,colback=poppurple,colframe=popink,
  fontupper=\popsemi\fontsize{10.4}{14.2}\selectfont\color{white},
  before upper={\poppill{Le savais-tu\,?}{white}{poppurple}\ifx\relax#1\relax\else\hspace{6pt}{\popxbold\color{white}#1}\fi\par\vspace{7pt}}}
% Exercice résolu : énoncé en haut, \tcblower puis la solution sur fond crème
\newcounter{popexo}\newcounter{popexr}
\newtcolorbox{exoresolu}[1][]{popbase,colframe=poporange,colbacklower=popcream,
  segmentation style={popink,line width=1.2pt,dashed},
  before upper={\stepcounter{popexr}\popboxhead{Exercice résolu \thepopexr}{poporange}{#1}},
  before lower={\poppill{Solution}{popink}{popcream}\par\vspace{6pt}}}
% Exercice (non résolu en place) — la correction est dans la partie Corrigés
\newtcolorbox{exercice}[1][]{popbase,colframe=popink,
  before upper={\stepcounter{popexo}\popboxhead{Exercice \thepopexo}{popink}{#1}}}
% Corrigé
\newtcolorbox{corrige}[1][]{popbase,colframe=popgreen,colback=popgreenL,
  before upper={\popboxhead{Corrigé}{popgreen}{#1}}}
% Objectifs / prérequis / bilan
\newtcolorbox{objectifs}{popbase,colframe=popink,colback=poporangeL,
  before upper={\popboxhead{Objectifs de la leçon}{popink}{}}}
\newtcolorbox{prerequis}{popbase,colframe=popink,colback=popblueL,
  before upper={\popboxhead{Prérequis}{popblue}{}}}
\newtcolorbox{bilan}{popbase,colframe=popink,colback=white,boxrule=2.8pt,
  before upper={\popboxhead{Fiche bilan}{poporange}{L'essentiel à connaître pour l'examen}}}
% Figure encadrée avec légende
\newtcolorbox{popfigure}[1][]{enhanced,breakable=false,colback=white,colframe=popline,boxrule=1.4pt,arc=8pt,
  left=10pt,right=10pt,top=10pt,bottom=8pt,before skip=10pt,after skip=12pt,halign=center,
  lower separated=false,halign lower=center,
  fontlower=\popsemi\fontsize{9}{11.5}\selectfont\color{popmuted},
  #1}
\newcommand{\legende}[1]{\par\vspace{6pt}{\popsemi\fontsize{9}{11.5}\selectfont\color{popmuted}#1}}

% ============================================================
% Signature OnBuch+
% ============================================================
\newcommand{\poplogoinline}[1]{{\poplogo\fontsize{#1}{#1}\selectfont\color{popink}OnBuch\color{poporange}+}}
\newcommand{\signatureonbuch}{%
  \begin{tcolorbox}[enhanced,colback=popink,colframe=popink,boxrule=2.2pt,arc=10pt,
      drop shadow=poporange,left=14pt,right=14pt,top=10pt,bottom=10pt,before skip=0pt,after skip=0pt]
    \tikz[baseline=-0.7ex]\node[circle,fill=popgreen,draw=popcream,line width=1.3pt,minimum size=22pt,inner sep=0]{\popcheck{white}};%
    \hspace{9pt}\begin{minipage}[c]{0.62\linewidth}
      {\popxbold\fontsize{12}{13}\selectfont\color{popcream}Signé\ \textbullet\ {\poplogo OnBuch\color{poporange}+}}\par\vspace{2pt}
      {\raggedright\popsemi\fontsize{8}{10.5}\selectfont\color{popline}\MakeUppercase{Équipe pédagogique OnBuch+}\\\MakeUppercase{Conforme au programme officiel MINESEC}\par}
    \end{minipage}\hfill
    {\poplogo\fontsize{22}{22}\selectfont\color{popcream}OnBuch\color{poporange}+}
  \end{tcolorbox}%
}

% ============================================================
% Page de garde
% #1 titre · #2 matière · #3 niveau · #4 module · #5 n° leçon · #6 durée ·
% #7 description / objectifs (texte ou itemize)
% ============================================================
\newcommand{\pagegarde}[7]{%
  \thispagestyle{empty}
  \newgeometry{a4paper,margin=1.7cm,top=1.6cm,bottom=1.4cm}%
  \begin{tikzpicture}[remember picture,overlay]
    \fill[poporange!18] ([xshift=-3.2cm,yshift=-3.6cm]current page.north east) circle (4.6cm);
    \fill[poppurple!14] ([xshift=2.4cm,yshift=7.6cm]current page.south west) circle (3.8cm);
  \end{tikzpicture}%
  {\poplogo\fontsize{30}{30}\selectfont\color{popink}OnBuch\color{poporange}+}\hfill
  \raisebox{0.6ex}{\poppill{Cours signé \textbullet\ Premium}{popink}{popcream}}\par
  \vspace{1pt}\popsquiggle{poppurple}\par
  \vspace{8pt}
  \begin{tcolorbox}[enhanced,colback=poporange,colframe=popink,boxrule=2.8pt,arc=13pt,
      drop shadow=popink,left=18pt,right=18pt,top=12pt,bottom=13pt,after skip=10pt]
    \poppill{#2 \textbullet\ #3}{popink}{popcream}\hspace{5pt}\poppill{#4}{white}{popink}\par\vspace{8pt}
    {\popdisplay\fontsize{14}{14}\selectfont\color{popink}LEÇON #5}\par\vspace{4pt}
    {\raggedright\hyphenpenalty=10000\exhyphenpenalty=10000\popdisplay\fontsize{25}{27}\selectfont\color{popcream}\MakeUppercase{#1}\par}
    \vspace{8pt}
    {\popxbold\fontsize{9.5}{9.5}\selectfont\color{popink}\MakeUppercase{Durée conseillée : #6}}
  \end{tcolorbox}
  \begin{tcolorbox}[enhanced,colback=white,colframe=popink,boxrule=2.2pt,arc=10pt,
      drop shadow=popink,left=16pt,right=16pt,top=10pt,bottom=10pt,after skip=8pt]
    \poppill{Dans cette leçon}{poppurple}{white}\par\vspace{5pt}
    \popsemi\fontsize{9.3}{12.6}\selectfont\color{popink}#7
  \end{tcolorbox}
  \noindent\begin{minipage}[t]{0.487\linewidth}
    \begin{tcolorbox}[enhanced,colback=popgreen,colframe=popink,boxrule=2.2pt,arc=10pt,
        drop shadow=popink,left=11pt,right=11pt,top=10pt,bottom=10pt,height=3.1cm,valign=center]
      \tcbox[colback=white,colframe=popink,boxrule=1.3pt,arc=5pt,boxsep=0pt,nobeforeafter,
        left=3pt,right=3pt,top=3pt,bottom=3pt,valign=center]{\qrcode[height=1.9cm]{https://play.google.com/store/apps/details?id=com.luvvix.onbuch&hl=fr}}%
      \hspace{8pt}\begin{minipage}[c]{0.5\linewidth}
        {\popdisplay\fontsize{11}{12.5}\selectfont\color{white}\MakeUppercase{Télécharge OnBuch}}\par\vspace{4pt}
        {\raggedright\popsemi\fontsize{8.4}{11}\selectfont\color{white}Gratuit \textbullet\ Google Play\\Tuteur IA, annales, concours, résultats\par}
      \end{minipage}
    \end{tcolorbox}
  \end{minipage}\hfill
  \begin{minipage}[t]{0.487\linewidth}
    \begin{tcolorbox}[enhanced,colback=poppurple,colframe=popink,boxrule=2.2pt,arc=10pt,
        drop shadow=popink,left=11pt,right=11pt,top=10pt,bottom=10pt,height=3.1cm,valign=center]
      \tikz[baseline=-0.7ex]\node[circle,fill=white,draw=popink,line width=1.3pt,minimum size=1.6cm,inner sep=0]{\popdisplay\fontsize{24}{24}\selectfont\color{poppurple}+};%
      \hspace{8pt}\begin{minipage}[c]{0.6\linewidth}
        {\popdisplay\fontsize{11}{12.5}\selectfont\color{white}\MakeUppercase{Passe à OnBuch+}}\par\vspace{4pt}
        {\raggedright\popsemi\fontsize{8.4}{11}\selectfont\color{white}Tous les cours signés, Tuteur IA illimité, fiches PDF — onglet OnBuch+ de l'app\par}
      \end{minipage}
    \end{tcolorbox}
  \end{minipage}
  \vspace{8pt}
  \popcontenu
  \vfill
  \signatureonbuch
  \restoregeometry
  \newpage
}

% Rangée « Ce cours contient » (page de garde)
\newcommand{\poptile}[3]{% #1 couleur #2 gros texte #3 légende
  \begin{tcolorbox}[enhanced,colback=#1,colframe=popink,boxrule=2pt,arc=9pt,shadow={2.5pt}{-2.5pt}{0pt}{popink},
      left=6pt,right=6pt,top=9pt,bottom=9pt,halign=center,valign=center,height=1.75cm,width=\linewidth,nobeforeafter]
    {\popdisplay\fontsize{12.5}{13}\selectfont\color{white}\MakeUppercase{#2}}\par\vspace{3pt}
    {\centering\popsemi\fontsize{7.8}{9.5}\selectfont\color{white}#3\par}
  \end{tcolorbox}}
\newcommand{\popcontenu}{%
  \par\noindent{\popxboldls\fontsize{7.5}{7.5}\selectfont\color{popmuted}CE COURS SIGNÉ CONTIENT}\par\vspace{7pt}
  \noindent\begin{minipage}[t]{0.235\linewidth}\poptile{popblue}{Cours}{complet, illustré, pas à pas}\end{minipage}\hfill
  \begin{minipage}[t]{0.235\linewidth}\poptile{poporange}{Méthodes}{et exercices résolus}\end{minipage}\hfill
  \begin{minipage}[t]{0.235\linewidth}\poptile{poppink}{Exercices}{type examen + corrigés}\end{minipage}\hfill
  \begin{minipage}[t]{0.235\linewidth}\poptile{popgreen}{Fiche bilan}{l'essentiel à retenir}\end{minipage}\par}

% Sommaire
\newtcolorbox{sommaire}{enhanced,colback=white,colframe=popink,boxrule=2.2pt,arc=10pt,
  drop shadow=popink,left=18pt,right=18pt,top=13pt,bottom=13pt,before skip=6pt,after skip=20pt,
  before upper={\poppill{Sommaire}{poppurple}{white}\par\vspace{9pt}\popsemi\fontsize{10.6}{14.5}\selectfont\color{popink}}}

% Signature de fin de cours — poussée tout en bas de la dernière page
\newcommand{\findecours}{%
  \par\vfill\nopagebreak
  \begin{center}\popsquiggle{poporange}\end{center}\vspace{4pt}
  \signatureonbuch
}
\AtBeginDocument{\def\llbracket{\mathopen{[\mkern-3.5mu[}}\def\rrbracket{\mathclose{]\mkern-3.5mu]}}} % intervalles d'entiers (glyphe ⟦ absent de la police)
% Glyphes absents de Fira Math : redéfinis à partir de glyphes disponibles
\AtBeginDocument{%
  \def\setminus{\mathbin{\backslash}}\let\smallsetminus\setminus
  \def\vdots{\mathord{\vbox{\baselineskip3.2pt\lineskiplimit0pt\kern4pt\hbox{.}\hbox{.}\hbox{.}}}}%
  \def\ddots{\mathinner{\mkern1mu\raise6.5pt\hbox{.}\mkern2mu\raise3.6pt\hbox{.}\mkern2mu\raise0.7pt\hbox{.}\mkern1mu}}%
  \def\triangle{\mathord{\scalebox{0.9}{$\symup{\Delta}$}}}%
  \def\nabla{\mathord{\raisebox{\depth}{\scalebox{1}[-1]{$\symup{\Delta}$}}}}%
  \def\checkmark{\popcheck{popdark}}%
  \def\star{\mathbin{\ast}}\def\bigstar{\mathord{\ast}}%
  \def\diamond{\mathbin{\scalebox{0.6}{$\Diamond$}}}%
  \def\bigcup{\mathop{\mathchoice{\raisebox{-0.3em}{\scalebox{1.7}{$\cup$}}}{\scalebox{1.25}{$\cup$}}{\cup}{\cup}}\displaylimits}%
  \def\bigcap{\mathop{\mathchoice{\raisebox{-0.3em}{\scalebox{1.7}{$\cap$}}}{\scalebox{1.25}{$\cap$}}{\cap}{\cap}}\displaylimits}%
  \def\lesseqqgtr{\mathrel{\lessgtr}}\def\bigoplus{\mathop{\oplus}\displaylimits}%
}
\providecommand{\ssi}{si et seulement si} % abréviation fréquente
\AtBeginDocument{\providecommand{\wideparen}{\overparen}} % arc de cercle

%% FIGFIX-BEGIN v5 — correctifs globaux des figures (docs/onbuch-generation/tools/figfix)
\usepackage{adjustbox}\usepackage{etoolbox}\tcbuselibrary{raster}
% toute figure TikZ/pgfplots plus large que la ligne est réduite à la largeur disponible
\BeforeBeginEnvironment{tikzpicture}{\begin{adjustbox}{max width=\linewidth}}
\AfterEndEnvironment{tikzpicture}{\end{adjustbox}}
% légendes pgfplots lisibles par-dessus les courbes
\pgfplotsset{every axis/.append style={legend style={fill=white,fill opacity=0.92,text opacity=1,draw=black!15,inner sep=3pt}}}
% étiquettes d'arêtes (syntaxe edge["..."]) avec fond blanc
\tikzset{every edge quotes/.append style={fill=white,inner sep=1.2pt}}
%% FIGFIX-END
\begin{document}

\pagegarde{Sphères}{Mathématiques}{1ère E}{Géométrie dans l'espace}{N20}{8 h}{Cette leçon étudie la sphère dans l'espace muni d'un repère orthonormé : son équation cartésienne, la reconnaissance des ensembles d'équation x²+y²+z²+ax+by+cz+d=0, puis les positions relatives d'une sphère avec une droite ou un plan (plan tangent, cercle d'intersection). Elle fournit les outils analytiques indispensables pour modéliser des situations concrètes de portée, de contact et de section.
\vspace{4pt}
\begin{multicols}{2}\small
\begin{itemize}
\item À la fin de cette leçon, je sais déterminer l'équation cartésienne d'une sphère connaissant son centre et son rayon
\item À la fin de cette leçon, je sais déterminer l'équation d'une sphère connaissant un diamètre ou passant par des points donnés
\item À la fin de cette leçon, je sais reconnaître la nature de l'ensemble d'équation x²+y²+z²+ax+by+cz+d=0 (sphère, point, ensemble vide) et donner ses éléments caractéristiques
\item À la fin de cette leçon, je sais calculer la distance d'un point à un plan et à une droite dans l'espace
\item À la fin de cette leçon, je sais déterminer analytiquement l'intersection d'une sphère et d'une droite et discuter selon le nombre de points
\item À la fin de cette leçon, je sais déterminer l'intersection d'une sphère et d'un plan (vide, point de tangence, cercle) et calculer le centre et le rayon du cercle d'intersection
\end{itemize}\end{multicols}}

\begin{sommaire}
\begin{multicols}{2}
\begin{enumerate}
\item Sphère : définition et équation cartésienne
\item Ensemble d'équation x²+y²+z²+ax+by+cz+d = 0
\item Position relative d'une sphère et d'une droite
\item Position relative d'une sphère et d'un plan ; plan tangent
\item Cercle d'intersection d'une sphère et d'un plan
\item Synthèse et applications
\item Activité d'intégration
\item Méthodes et exercices résolus
\item Exercices
\item Corrigés des exercices
\item Fiche bilan
\end{enumerate}
\end{multicols}
\end{sommaire}

\begin{objectifs}
\begin{itemize}
\item À la fin de cette leçon, je sais déterminer l'équation cartésienne d'une sphère connaissant son centre et son rayon
\item À la fin de cette leçon, je sais déterminer l'équation d'une sphère connaissant un diamètre ou passant par des points donnés
\item À la fin de cette leçon, je sais reconnaître la nature de l'ensemble d'équation x²+y²+z²+ax+by+cz+d=0 (sphère, point, ensemble vide) et donner ses éléments caractéristiques
\item À la fin de cette leçon, je sais calculer la distance d'un point à un plan et à une droite dans l'espace
\item À la fin de cette leçon, je sais déterminer analytiquement l'intersection d'une sphère et d'une droite et discuter selon le nombre de points
\item À la fin de cette leçon, je sais déterminer l'intersection d'une sphère et d'un plan (vide, point de tangence, cercle) et calculer le centre et le rayon du cercle d'intersection
\item À la fin de cette leçon, je sais déterminer une équation du plan tangent à une sphère en un point donné
\item À la fin de cette leçon, je sais modéliser et résoudre un problème concret (portée d'un émetteur, section d'une sphère) à l'aide des équations de sphères
\end{itemize}
\end{objectifs}

\begin{prerequis}
\begin{itemize}
\item Repère orthonormé de l'espace et coordonnées d'un point, d'un vecteur
\item Distance entre deux points et milieu d'un segment dans l'espace
\item Équation cartésienne d'un plan, vecteur normal, représentation paramétrique d'une droite
\item Distance d'un point à un plan : d(A,P) = |ax₀+by₀+cz₀+d| / √(a²+b²+c²)
\item Résolution d'équations du second degré et discriminant
\end{itemize}
\end{prerequis}

\newpage

\coursec{Sphère : définition et équation cartésienne}

\emph{Situation d'accroche.} Au marché Mokolo de Yaoundé, on croise partout la même forme : l'antenne parabolique du vendeur de télévisions, le ballon de football que les enfants font rouler entre les étals, l'orange que l'on épluche. Tous ces objets partagent une géométrie commune : celle de la \cle{sphère}. Un technicien de Camtel doit installer un émetteur dont la portée est une sphère de $5$ km de rayon : comment savoir si un quartier donné est couvert ? Un ballon posé sur le sol touche le sol en un seul point : pourquoi ? Un puits creusé dans une colline sphérique coupe la surface selon un cercle : quel est son rayon ?

\medskip
\textbf{Problématique.} Comment décrire une sphère de l'espace par une équation, et comment utiliser cette équation pour décider si un point est sur la sphère, à l'intérieur ou à l'extérieur ? C'est l'objet de cette première section : nous allons définir rigoureusement la sphère, établir son \cle{équation cartésienne} dans un repère orthonormé, puis découvrir une caractérisation remarquable de la sphère de diamètre $[AB]$ à l'aide du produit scalaire.

\courssub{Définition géométrique de la sphère}

\textbf{Intuition.} Dans le plan, le cercle de centre $\Omega$ et de rayon $R$ est l'ensemble des points situés à la distance $R$ de $\Omega$. Dans l'espace, on ne se contente plus de tourner dans un plan : on « tourne dans toutes les directions ». Le cercle devient alors une \emph{surface} : la sphère. Imagine une corde de longueur $R$ fixée en un point $\Omega$ : en tendant la corde dans toutes les directions possibles de l'espace, son extrémité libre décrit exactement la sphère de centre $\Omega$ et de rayon $R$.

\begin{definition}[Sphère]
Soit $\Omega$ un point de l'espace et $R$ un réel strictement positif. On appelle \cle{sphère} de centre $\Omega$ et de rayon $R$, notée $\mathcal{S}(\Omega, R)$, l'ensemble des points $M$ de l'espace tels que :
\[ \Omega M = R. \]
Autrement dit : $\mathcal{S}(\Omega, R) = \{ M \in \mathcal{E} \;/\; \Omega M = R \}$.
\end{definition}

\begin{attention}[Sphère ou boule ?]
La sphère est une \textbf{surface} : elle ne contient que les points situés \emph{exactement} à la distance $R$ de $\Omega$. Un ballon de football (l'enveloppe) est une sphère ; une orange pleine est une \emph{boule}. Ne dis jamais « le point est dans la sphère » : un point est \emph{sur} la sphère, ou bien à l'\emph{intérieur}, ou bien à l'\emph{extérieur}.
\end{attention}

\begin{popfigure}
\begin{center}
\begin{tikzpicture}[scale=1.15]
  % Sphère : disque + grand cercle + équateur en pointillés
  \shade[ball color=popblueL, opacity=0.55] (0,0) circle (2);
  \draw[popblue, thick] (0,0) circle (2);
  \draw[popblue, dashed] (-2,0) arc (180:360:2 and 0.6);
  \draw[popblue] (2,0) arc (0:180:2 and 0.6);
  % Centre et rayon
  \coordinate (O) at (0,0);
  \coordinate (M) at (1.414,1.414);
  \fill[popdark] (O) circle (1.6pt) node[below left=2pt] {$\Omega$};
  \draw[poporange, very thick] (O) -- (M) node[midway, below right, popdark] {$R$};
  \fill[poporange] (M) circle (1.8pt) node[above right=1pt, popdark] {$M$};
\end{tikzpicture}
\end{center}
\legende{Sphère $\mathcal{S}(\Omega, R)$ : pour tout point $M$ de la surface, $\Omega M = R$.}
\end{popfigure}

\begin{definition}[Boule ouverte, boule fermée]
Soit $\Omega$ un point de l'espace et $R > 0$.
\begin{itemize}
  \item La \cle{boule ouverte} de centre $\Omega$ et de rayon $R$ est l'ensemble des points $M$ tels que $\Omega M < R$ (l'intérieur strict de la sphère).
  \item La \cle{boule fermée} de centre $\Omega$ et de rayon $R$ est l'ensemble des points $M$ tels que $\Omega M \leq R$ (la sphère \emph{plus} son intérieur).
\end{itemize}
Un point $M$ tel que $\Omega M > R$ est dit \emph{extérieur} à la sphère.
\end{definition}

\begin{exemplebox}[Couverture de l'émetteur Camtel]
L'émetteur du technicien est modélisé par une sphère de centre $\Omega$ (la position de l'antenne) et de rayon $R = 5$ km. Un quartier situé au point $M$ est couvert si et seulement si $\Omega M \leq 5$ km, c'est-à-dire si $M$ appartient à la \textbf{boule fermée} de centre $\Omega$ et de rayon $5$. La sphère elle-même est la « limite de portée » du signal.
\end{exemplebox}

\courssub{Équation cartésienne d'une sphère}

\textbf{Intuition.} Plaçons-nous dans un repère orthonormé $(O\,;\vec{i},\vec{j},\vec{k})$ de l'espace. Dire que $M(x\,;y\,;z)$ appartient à la sphère de centre $\Omega(a\,;b\,;c)$ et de rayon $R$, c'est dire que la distance $\Omega M$ vaut $R$. Or nous connaissons la formule de la distance entre deux points : il suffit de la traduire !

\begin{propriete}[Équation cartésienne d'une sphère — démonstration exigible]
Dans un repère orthonormé $(O\,;\vec{i},\vec{j},\vec{k})$, la sphère $\mathcal{S}(\Omega, R)$ de centre $\Omega(a\,;b\,;c)$ et de rayon $R > 0$ admet pour équation cartésienne :
\[ (x-a)^2 + (y-b)^2 + (z-c)^2 = R^2. \]
Autrement dit : $M(x\,;y\,;z) \in \mathcal{S}(\Omega, R) \iff (x-a)^2+(y-b)^2+(z-c)^2 = R^2$.
\end{propriete}

\begin{experience}[Démonstration guidée de l'équation cartésienne]
Suivons le raisonnement pas à pas.
\begin{enumerate}
  \item \textbf{Traduction de l'appartenance.} Par définition, $M(x\,;y\,;z) \in \mathcal{S}(\Omega, R) \iff \Omega M = R$.
  \item \textbf{Formule de la distance.} Dans un repère orthonormé, avec $\Omega(a\,;b\,;c)$ et $M(x\,;y\,;z)$ :
  \[ \Omega M = \sqrt{(x-a)^2 + (y-b)^2 + (z-c)^2}. \]
  \item \textbf{Équivalence.} On a donc :
  \[ M \in \mathcal{S}(\Omega, R) \iff \sqrt{(x-a)^2+(y-b)^2+(z-c)^2} = R. \]
  \item \textbf{Élévation au carré.} Les deux membres sont positifs ($R>0$), donc l'égalité est préservée en élevant au carré :
  \[ M \in \mathcal{S}(\Omega, R) \iff (x-a)^2+(y-b)^2+(z-c)^2 = R^2. \]
\end{enumerate}
L'équivalence est démontrée. $\blacksquare$
\end{experience}

\begin{aretenir}[Les deux formes de l'équation d'une sphère]
\textbf{Forme réduite} (centre et rayon visibles) :
\[ (x-a)^2 + (y-b)^2 + (z-c)^2 = R^2. \]
\textbf{Forme développée}, obtenue en développant les carrés :
\[ x^2+y^2+z^2 - 2ax - 2by - 2cz + \left(a^2+b^2+c^2-R^2\right) = 0. \]
On passe de l'une à l'autre par développement (réduite $\to$ développée) ou par \emph{formes canoniques} (développée $\to$ réduite), technique que nous détaillerons dans la section suivante.
\end{aretenir}

\begin{exemplebox}[Équation d'une sphère de centre et rayon donnés]
Déterminons une équation cartésienne de la sphère $\mathcal{S}$ de centre $\Omega(1\,;-2\,;3)$ et de rayon $R = 4$.

La forme réduite donne directement :
\[ (x-1)^2 + (y-(-2))^2 + (z-3)^2 = 4^2 \quad\text{soit}\quad (x-1)^2+(y+2)^2+(z-3)^2 = 16. \]
Développons pour obtenir la forme développée :
\begin{align*}
(x-1)^2+(y+2)^2+(z-3)^2 &= 16\\
x^2-2x+1 + y^2+4y+4 + z^2-6z+9 &= 16\\
x^2+y^2+z^2 -2x+4y-6z -2 &= 0.
\end{align*}
Vérification avec la formule : $-2a=-2$, $-2b=4$, $-2c=-6$ et $a^2+b^2+c^2-R^2 = 1+4+9-16 = -2$. C'est cohérent.
\end{exemplebox}

\begin{attention}[Les deux pièges classiques]
\begin{enumerate}
  \item \textbf{Ne pas confondre $R$ et $R^2$.} Le membre de droite est $R^2$, pas $R$ ! Pour la sphère de rayon $4$, on écrit $=16$, jamais $=4$.
  \item \textbf{Attention aux signes des coordonnées du centre.} Dans $(x-a)^2$, le signe devant $a$ est un \emph{moins}. Ainsi $(y+2)^2$ correspond à $b=-2$ et non $b=2$. Pour lire le centre sur une équation, change le signe de ce qui accompagne $x$, $y$, $z$ dans les parenthèses.
\end{enumerate}
\end{attention}

\begin{exoresolu}[Sphère définie par son centre et un point]
\textbf{Énoncé.} Déterminer une équation cartésienne de la sphère $\mathcal{S}$ de centre $\Omega(2\,;-1\,;0)$ passant par le point $A(2\,;3\,;4)$.
\tcblower
\textbf{Solution détaillée.}

Le centre est connu, mais pas le rayon. Or le rayon est précisément la distance du centre à n'importe quel point de la sphère, donc $R = \Omega A$.

\textbf{Étape 1 : calcul du rayon.}
\begin{align*}
R = \Omega A &= \sqrt{(2-2)^2+(3-(-1))^2+(4-0)^2}\\
&= \sqrt{0+16+16} = \sqrt{32} = 4\sqrt{2}.
\end{align*}
Donc $R^2 = 32$ (inutile de garder le radical : seul $R^2$ intervient dans l'équation).

\textbf{Étape 2 : écriture de l'équation.}
\[ (x-2)^2+(y+1)^2+z^2 = 32. \]

\textbf{Étape 3 (facultative) : forme développée.}
\[ x^2-4x+4+y^2+2y+1+z^2 = 32 \iff x^2+y^2+z^2-4x+2y-27=0. \]

\textbf{Vérification.} Le point $A$ doit satisfaire l'équation : $(2-2)^2+(3+1)^2+4^2 = 0+16+16=32$. C'est bien le cas.
\end{exoresolu}

\courssub{Sphère de diamètre $[AB]$ : caractérisation par le produit scalaire}

\textbf{Intuition.} Dans le plan, tu connais le théorème de l'angle inscrit : un triangle $AMB$ est rectangle en $M$ si et seulement si $M$ appartient au cercle de diamètre $[AB]$. Ce résultat se généralise à l'espace : le cercle devient une sphère, et la condition « angle droit » s'exprime élégamment par le produit scalaire $\overrightarrow{MA}\cdot\overrightarrow{MB}=0$.

\begin{propriete}[{Caractérisation de la sphère de diamètre $[AB]$}]
Soient $A$ et $B$ deux points distincts de l'espace. Un point $M$ appartient à la sphère de diamètre $[AB]$ si et seulement si :
\[ \overrightarrow{MA}\cdot\overrightarrow{MB} = 0. \]
\end{propriete}

\begin{experience}[{Démonstration guidée via le milieu de $[AB]$}]
Notons $I$ le milieu de $[AB]$. La sphère de diamètre $[AB]$ a pour centre $I$ et pour rayon $\dfrac{AB}{2} = IA$.
\begin{enumerate}
  \item \textbf{Décomposition de Chasles.} Introduisons le point $I$ dans chaque vecteur :
  \[ \overrightarrow{MA} = \overrightarrow{MI}+\overrightarrow{IA} \qquad\text{et}\qquad \overrightarrow{MB} = \overrightarrow{MI}+\overrightarrow{IB}. \]
  \item \textbf{Utilisation du milieu.} Comme $I$ est le milieu de $[AB]$, $\overrightarrow{IB} = -\overrightarrow{IA}$. Donc :
  \[ \overrightarrow{MA}\cdot\overrightarrow{MB} = \left(\overrightarrow{MI}+\overrightarrow{IA}\right)\cdot\left(\overrightarrow{MI}-\overrightarrow{IA}\right). \]
  \item \textbf{Identité remarquable vectorielle.} En développant (le produit scalaire est distributif et commutatif) :
  \[ \overrightarrow{MA}\cdot\overrightarrow{MB} = \overrightarrow{MI}^{\,2} - \overrightarrow{IA}^{\,2} = MI^2 - IA^2. \]
  \item \textbf{Conclusion.} Par conséquent :
  \[ \overrightarrow{MA}\cdot\overrightarrow{MB}=0 \iff MI^2 = IA^2 \iff MI = IA = \frac{AB}{2}, \]
  ce qui signifie exactement que $M$ appartient à la sphère de centre $I$ et de rayon $\dfrac{AB}{2}$, c'est-à-dire la sphère de diamètre $[AB]$. $\blacksquare$
\end{enumerate}
\end{experience}

\begin{popfigure}
\begin{center}
\begin{tikzpicture}[scale=1.05]
  \shade[ball color=poppurpleL, opacity=0.5] (0,0) circle (2.2);
  \draw[poppurple, thick] (0,0) circle (2.2);
  \draw[poppurple, dashed] (-2.2,0) arc (180:360:2.2 and 0.65);
  \draw[poppurple] (2.2,0) arc (0:180:2.2 and 0.65);
  % Diamètre AB
  \coordinate (A) at (-1.9,-1.1);
  \coordinate (B) at (1.9,1.1);
  \coordinate (I) at (0,0);
  \coordinate (M) at (0.6,2.116);
  \draw[popdark, thick] (A) -- (B);
  \draw[poporange, thick] (M) -- (A);
  \draw[poporange, thick] (M) -- (B);
  % Angle droit en M
  \draw[popgreen, thick] ($(M)!0.28cm!(A)$) -- ++($(M)!0.28cm!(B)-(M)$) -- ++($(M)!0.28cm!(A)-(M)$) -- cycle;
  \fill[popdark] (A) circle (1.6pt) node[below left] {$A$};
  \fill[popdark] (B) circle (1.6pt) node[above right] {$B$};
  \fill[popblue] (I) circle (1.6pt) node[below right=1pt] {$I$};
  \fill[poporange] (M) circle (1.8pt) node[above=1pt, popdark] {$M$};
\end{tikzpicture}
\end{center}
\legende{Sphère de diamètre $[AB]$, de centre $I$ milieu de $[AB]$ : pour tout point $M$ de la sphère, $\overrightarrow{MA}\cdot\overrightarrow{MB}=0$ (angle droit en $M$).}
\end{popfigure}

\begin{methode}[{Équation d'une sphère de diamètre $[AB]$}]
Pour déterminer l'équation cartésienne de la sphère de diamètre $[AB]$ :
\begin{enumerate}
  \item Calculer les coordonnées du \textbf{centre} $I$, milieu de $[AB]$ :
  $I\left(\dfrac{x_A+x_B}{2}\,;\dfrac{y_A+y_B}{2}\,;\dfrac{z_A+z_B}{2}\right)$.
  \item Calculer le \textbf{rayon} $R = \dfrac{AB}{2}$ (ou directement $R^2 = \dfrac{AB^2}{4}$).
  \item Écrire l'équation $(x-x_I)^2+(y-y_I)^2+(z-z_I)^2 = R^2$.
\end{enumerate}
\emph{Variante rapide} : on peut aussi écrire directement $\overrightarrow{MA}\cdot\overrightarrow{MB}=0$ avec $M(x\,;y\,;z)$, ce qui donne immédiatement la forme développée.
\end{methode}

\begin{exoresolu}[{Sphère de diamètre $[AB]$}]
\textbf{Énoncé.} Soient $A(1\,;0\,;2)$ et $B(3\,;2\,;0)$. Déterminer une équation cartésienne de la sphère $\mathcal{S}$ de diamètre $[AB]$.
\tcblower
\textbf{Solution détaillée.}

\textbf{Méthode 1 : centre et rayon.}

Le centre est $I$, milieu de $[AB]$ :
\[ I\left(\frac{1+3}{2}\,;\frac{0+2}{2}\,;\frac{2+0}{2}\right) = I(2\,;1\,;1). \]
Calculons $AB^2 = (3-1)^2+(2-0)^2+(0-2)^2 = 4+4+4 = 12$, donc $AB = 2\sqrt{3}$ et $R^2 = \dfrac{AB^2}{4} = 3$.

L'équation de $\mathcal{S}$ est donc :
\[ (x-2)^2+(y-1)^2+(z-1)^2 = 3. \]

\textbf{Méthode 2 (vérification) : produit scalaire.} Avec $M(x\,;y\,;z)$ :
\[ \overrightarrow{MA} = \begin{pmatrix}1-x\\-y\\2-z\end{pmatrix}, \qquad \overrightarrow{MB} = \begin{pmatrix}3-x\\2-y\\-z\end{pmatrix}. \]
\begin{align*}
\overrightarrow{MA}\cdot\overrightarrow{MB} &= (1-x)(3-x)+(-y)(2-y)+(2-z)(-z)\\
&= x^2-4x+3 + y^2-2y + z^2-2z.
\end{align*}
La condition $\overrightarrow{MA}\cdot\overrightarrow{MB}=0$ donne $x^2+y^2+z^2-4x-2y-2z+3=0$. Or en développant la forme réduite : $x^2-4x+4+y^2-2y+1+z^2-2z+1=3$, soit $x^2+y^2+z^2-4x-2y-2z+3=0$. Les deux méthodes concordent parfaitement.
\end{exoresolu}

\courssub{Intérieur et extérieur d'une sphère : inéquations}

La traduction analytique de la boule est immédiate à partir de l'équation de la sphère : il suffit de remplacer le signe $=$ par une inégalité.

\begin{aretenir}[Position d'un point par rapport à une sphère]
Soit $\mathcal{S}$ la sphère d'équation $(x-a)^2+(y-b)^2+(z-c)^2=R^2$ et $M(x\,;y\,;z)$ un point. On pose $d^2 = (x-a)^2+(y-b)^2+(z-c)^2$ (c'est $\Omega M^2$). Alors :
\begin{itemize}
  \item $d^2 = R^2$ : $M$ est \textbf{sur} la sphère ;
  \item $d^2 < R^2$ : $M$ est à l'\textbf{intérieur} (boule ouverte) ;
  \item $d^2 \leq R^2$ : $M$ est dans la \textbf{boule fermée} ;
  \item $d^2 > R^2$ : $M$ est à l'\textbf{extérieur}.
\end{itemize}
\end{aretenir}

\begin{exemplebox}[Le quartier est-il couvert ?]
Reprenons l'émetteur Camtel : sphère de centre $\Omega(0\,;0\,;0)$ et de rayon $5$ (unité : le km), d'équation $x^2+y^2+z^2=25$. Le quartier de Nkolbisson est modélisé par le point $M(3\,;4\,;1)$. Calculons :
\[ \Omega M^2 = 3^2+4^2+1^2 = 9+16+1 = 26 > 25. \]
Donc $\Omega M = \sqrt{26} \approx 5{,}1$ km $> 5$ km : le quartier est \textbf{hors de portée}, de justesse ! Le technicien devra rehausser l'antenne ou installer un relais.
\end{exemplebox}

\begin{savaistu}[Pourquoi la Terre est-elle (presque) une sphère ?]
Sous l'effet de sa propre gravitation, tout astre suffisamment massif adopte une forme sphérique : la sphère est la surface qui minimise l'énergie pour un volume donné — c'est aussi pourquoi les bulles de savon sont sphériques. La Terre est en réalité légèrement aplatie aux pôles (on parle d'\emph{ellipsoïde}), mais pour les calculs de portée des antennes du Cameroun, le modèle sphérique est largement suffisant. Le rayon moyen de la Terre est d'environ $\SI{6371}{km}$ : la sphère d'équation $x^2+y^2+z^2 = 6371^2$ centrée au centre de la Terre est un excellent modèle de sa surface.
\end{savaistu}

\begin{exercice}[Pour s'entraîner]
\begin{enumerate}
  \item Déterminer une équation cartésienne de la sphère de centre $\Omega(0\,;1\,;-3)$ et de rayon $R=\sqrt{5}$, sous forme réduite puis développée.
  \item La sphère $\mathcal{S}$ a pour centre $\Omega(1\,;1\,;1)$ et passe par l'origine $O$ du repère. Donner son équation.
  \item Soient $A(2\,;1\,;3)$ et $B(0\,;3\,;1)$. Déterminer l'équation de la sphère de diamètre $[AB]$ par les deux méthodes du cours.
  \item Le point $P(2\,;-1\,;2)$ appartient-il à la boule fermée de centre $\Omega(1\,;0\,;1)$ et de rayon $2$ ?
\end{enumerate}
\end{exercice}

\begin{corrige}[Exercice 1]
\begin{enumerate}
  \item Forme réduite : $x^2+(y-1)^2+(z+3)^2=5$. Forme développée : $x^2+y^2-2y+1+z^2+6z+9=5$, soit $x^2+y^2+z^2-2y+6z+5=0$.
  \item $R^2=\Omega O^2=1+1+1=3$, d'où $(x-1)^2+(y-1)^2+(z-1)^2=3$.
  \item Milieu $I(1\,;2\,;2)$ ; $AB^2=4+4+4=12$ donc $R^2=3$ : $(x-1)^2+(y-2)^2+(z-2)^2=3$. Par le produit scalaire : $(2-x)(-x)+(1-y)(3-y)+(3-z)(1-z)=0$ donne $x^2-2x+y^2-4y+3+z^2-4z+3=0$, soit $x^2+y^2+z^2-2x-4y-4z+6=0$, ce qui correspond au développement de la forme réduite.
  \item $\Omega P^2=(2-1)^2+(-1-0)^2+(2-1)^2=1+1+1=3 \leq 4=R^2$ : oui, $P$ est dans la boule fermée (et même à l'intérieur strict).
\end{enumerate}
\end{corrige}

\medskip
\textbf{En résumé.} La sphère $\mathcal{S}(\Omega,R)$ est l'ensemble des points $M$ tels que $\Omega M = R$. Dans un repère orthonormé, avec $\Omega(a\,;b\,;c)$, elle admet l'équation $(x-a)^2+(y-b)^2+(z-c)^2=R^2$, qui se développe en $x^2+y^2+z^2-2ax-2by-2cz+(a^2+b^2+c^2-R^2)=0$. La sphère de diamètre $[AB]$ est caractérisée par $\overrightarrow{MA}\cdot\overrightarrow{MB}=0$, et son équation s'obtient via le milieu de $[AB]$ et le rayon $\frac{AB}{2}$. Dans la section suivante, nous ferons le chemin inverse : partant d'une équation $x^2+y^2+z^2+ax+by+cz+d=0$, nous déterminerons la nature géométrique de l'ensemble qu'elle représente.

\coursec{Ensemble d'équation $x^2+y^2+z^2+ax+by+cz+d = 0$}

Dans la section précédente, nous avons vu qu'une sphère de centre $\Omega(x_0\,;\,y_0\,;\,z_0)$ et de rayon $R$ admet, dans un repère orthonormé $(O\,;\vec{\imath},\vec{\jmath},\vec{k})$ de l'espace, l'équation cartésienne
\[
(x-x_0)^2+(y-y_0)^2+(z-z_0)^2 = R^2,
\]
c'est-à-dire, après développement,
\[
x^2+y^2+z^2-2x_0\,x-2y_0\,y-2z_0\,z+x_0^2+y_0^2+z_0^2-R^2=0.
\]
Toute sphère possède donc une équation de la forme $x^2+y^2+z^2+ax+by+cz+d=0$. Se pose alors la question \emph{réciproque}, et c'est toute la problématique de cette section : si l'on vous donne une équation de cette forme, l'ensemble des points $M(x\,;\,y\,;\,z)$ dont les coordonnées la vérifient est-il \emph{toujours} une sphère ? La réponse, comme nous allons le démontrer, est nuancée : ce peut être une sphère, un seul point, ou... aucun point du tout. Savoir trancher entre ces trois cas est une compétence classique du Probatoire.

\courssub{De quelles équations parlons-nous exactement ?}

\begin{definition}[Équation du type étudié]
On considère, dans un repère orthonormé de l'espace, l'ensemble $(E)$ des points $M(x\,;\,y\,;\,z)$ dont les coordonnées vérifient une équation de la forme
\[
x^2+y^2+z^2+ax+by+cz+d=0,
\]
où $a$, $b$, $c$ et $d$ sont quatre réels donnés.
\end{definition}

Deux conditions structurelles caractérisent ces équations, et il faut les vérifier \emph{avant} tout calcul :
\begin{itemize}
\item les coefficients devant $x^2$, $y^2$ et $z^2$ sont \cle{égaux} (ici tous égaux à $1$ ; s'ils valent tous $2$, on divise toute l'équation par $2$ pour s'y ramener) ;
\item il n'y a \cle{aucun terme rectangle}, c'est-à-dire pas de terme en $xy$, en $yz$ ni en $zx$.
\end{itemize}

\begin{attention}[Vérifier la forme de l'équation avant de conclure]
L'équation $x^2+2y^2+z^2-4x+1=0$ n'est \textbf{pas} du type étudié (coefficients de $x^2$, $y^2$, $z^2$ différents) : son ensemble n'est en général pas une sphère (c'est un ellipsoïde, hors programme). De même, $x^2+y^2+z^2+2xy-1=0$ contient un terme en $xy$ : les résultats de cette section ne s'appliquent pas. Première étape réflexe, toujours : \textbf{regarder les coefficients des carrés et l'absence de termes rectangles}.
\end{attention}

\courssub{L'idée directrice : compléter les carrés}

L'intuition est simple. Dans l'équation développée d'une sphère, les termes $x^2-2x_0x$ proviennent du développement de $(x-x_0)^2$. Pour « remonter » de la forme développée à la forme $(x-x_0)^2+(y-y_0)^2+(z-z_0)^2=R^2$, il suffit donc de \emph{reconstituer les carrés}, variable par variable. C'est la technique dite des \cle{formes canoniques}, que tu as déjà rencontrée en Seconde pour le trinôme $ax^2+bx+c$.

Rappelons l'identité fondamentale : pour tous réels $x$ et $a$,
\[
x^2+ax = \left(x+\frac{a}{2}\right)^2-\frac{a^2}{4}.
\]
Autrement dit, pour « compléter le carré » en $x$, on prend la \textbf{moitié du coefficient de $x$}, on l'élève au carré, et on compense en soustrayant ce carré. On procédera de même, indépendamment, pour $y$ et pour $z$.

\courssub{Le théorème fondamental et sa démonstration}

\begin{propriete}[Théorème — nature de l'ensemble $(E)$]
Dans un repère orthonormé de l'espace, soit $(E)$ l'ensemble des points $M(x\,;\,y\,;\,z)$ tels que
\[
x^2+y^2+z^2+ax+by+cz+d=0.
\]
On pose
\[
k=\frac{a^2+b^2+c^2}{4}-d \qquad \text{et} \qquad \Omega\left(-\frac{a}{2}\,;\,-\frac{b}{2}\,;\,-\frac{c}{2}\right).
\]
Alors trois cas, et trois cas seulement, se présentent :
\begin{itemize}
\item si $k>0$ : $(E)$ est la \cle{sphère} de centre $\Omega$ et de rayon $R=\sqrt{k}$ ;
\item si $k=0$ : $(E)$ est réduit au \cle{singleton} $\{\Omega\}$ ;
\item si $k<0$ : $(E)$ est \cle{l'ensemble vide} $\varnothing$.
\end{itemize}
\emph{Cette démonstration est exigible au Probatoire : elle ne demande que la technique des formes canoniques.}
\end{propriete}

\begin{experience}[Démonstration guidée du théorème]
Partons de l'équation et transformons-la par équivalences successives.

\textbf{Étape 1 : compléter le carré en $x$.}
\[
x^2+ax=\left(x+\frac{a}{2}\right)^2-\frac{a^2}{4}.
\]

\textbf{Étape 2 : compléter le carré en $y$.}
\[
y^2+by=\left(y+\frac{b}{2}\right)^2-\frac{b^2}{4}.
\]

\textbf{Étape 3 : compléter le carré en $z$.}
\[
z^2+cz=\left(z+\frac{c}{2}\right)^2-\frac{c^2}{4}.
\]

\textbf{Étape 4 : regrouper.} En reportant dans l'équation initiale :
\begin{align*}
x^2+y^2+z^2+ax+by+cz+d=0
&\iff \left(x+\tfrac{a}{2}\right)^2-\tfrac{a^2}{4}+\left(y+\tfrac{b}{2}\right)^2-\tfrac{b^2}{4}+\left(z+\tfrac{c}{2}\right)^2-\tfrac{c^2}{4}+d=0\\
&\iff \left(x+\tfrac{a}{2}\right)^2+\left(y+\tfrac{b}{2}\right)^2+\left(z+\tfrac{c}{2}\right)^2=\underbrace{\tfrac{a^2+b^2+c^2}{4}-d}_{=\,k}.
\end{align*}

\textbf{Étape 5 : discussion selon le signe de $k$.} Notons $\Omega\left(-\frac{a}{2}\,;\,-\frac{b}{2}\,;\,-\frac{c}{2}\right)$. Le membre de gauche n'est autre que $\Omega M^2$. L'équation équivaut donc à
\[
\Omega M^2 = k.
\]
\begin{itemize}
\item Si $k>0$ : $\Omega M = \sqrt{k}$, donc $(E)$ est la sphère de centre $\Omega$ et de rayon $\sqrt{k}$.
\item Si $k=0$ : $\Omega M = 0$, donc $M=\Omega$ : l'ensemble $(E)$ est le singleton $\{\Omega\}$.
\item Si $k<0$ : une somme de carrés de réels est toujours positive ou nulle, donc ne peut jamais être égale à un nombre strictement négatif : aucun point ne convient, $(E)=\varnothing$.
\end{itemize}
Le théorème est démontré. $\blacksquare$
\end{experience}

\begin{aretenir}[À retenir absolument]
Pour l'ensemble d'équation $x^2+y^2+z^2+ax+by+cz+d=0$ :
\[
\text{Centre (potentiel) : } \Omega\left(-\frac{a}{2}\,;\,-\frac{b}{2}\,;\,-\frac{c}{2}\right), \qquad k=\frac{a^2+b^2+c^2}{4}-d.
\]
La nature de $(E)$ dépend du signe de $k$ : $k>0$ sphère de rayon $\sqrt{k}$ ; $k=0$ point $\Omega$ ; $k<0$ ensemble vide. On ne conclut \textbf{jamais} « sphère » sans avoir calculé $k$ et vérifié qu'il est strictement positif.
\end{aretenir}

\begin{popfigure}
\begin{center}
\begin{tikzpicture}[
  font=\small,
  decision/.style={rectangle, rounded corners, draw=popblue, fill=popblueL, text width=4.6cm, align=center, minimum height=1.1cm},
  result/.style={rectangle, rounded corners, draw=popdark, text width=3.6cm, align=center, minimum height=1.3cm},
  good/.style={result, draw=popgreen, fill=popgreenL},
  mid/.style={result, draw=popgold, fill=popgoldL},
  bad/.style={result, draw=poppink, fill=poppinkL},
  fleche/.style={-{Stealth[length=2.5mm]}, thick, popmuted}
]
\node[decision, align=center] (calc) at (0,0) {Équation $x^2+y^2+z^2+ax+by+cz+d=0$\\[2pt] On calcule $k=\dfrac{a^2+b^2+c^2}{4}-d$};
\node[good, align=center] (pos) at (-5.4,-3.2) {$k>0$\\[2pt] \textbf{Sphère} de centre $\Omega\left(-\frac{a}{2};-\frac{b}{2};-\frac{c}{2}\right)$ et de rayon $\sqrt{k}$};
\node[mid, align=center] (nul) at (0,-3.2) {$k=0$\\[2pt] \textbf{Singleton} $\{\Omega\}$ (un seul point)};
\node[bad, align=center] (neg) at (5.4,-3.2) {$k<0$\\[2pt] \textbf{Ensemble vide} $\varnothing$ (aucun point)};
\draw[fleche] (calc.south) -- (pos.north);
\draw[fleche] (calc.south) -- (nul.north);
\draw[fleche] (calc.south) -- (neg.north);
\end{tikzpicture}
\end{center}
\legende{Arbre de décision : la nature de l'ensemble $(E)$ se lit sur le signe de $k$.}
\end{popfigure}

\courssub{Méthode pratique : identifier centre et rayon}

\begin{methode}[Déterminer la nature de $(E)$]
Pour déterminer la nature et les éléments caractéristiques de l'ensemble d'équation $x^2+y^2+z^2+ax+by+cz+d=0$ :
\begin{enumerate}
\item \textbf{Vérifier la forme} : coefficients de $x^2$, $y^2$, $z^2$ égaux (sinon, diviser pour s'y ramener) et aucun terme en $xy$, $yz$, $zx$.
\item \textbf{Compléter les carrés} successivement : écrire $x^2+ax=\left(x+\frac{a}{2}\right)^2-\frac{a^2}{4}$, et de même pour $y$ et $z$.
\item \textbf{Réécrire l'équation} sous la forme $(x-x_0)^2+(y-y_0)^2+(z-z_0)^2=k$.
\item \textbf{Conclure selon le signe de $k$} : sphère (donner centre et rayon $\sqrt{k}$), singleton, ou ensemble vide.
\end{enumerate}
En pratique, on peut aussi appliquer directement les formules $\Omega\left(-\frac{a}{2}\,;\,-\frac{b}{2}\,;\,-\frac{c}{2}\right)$ et $k=\frac{a^2+b^2+c^2}{4}-d$, mais la rédaction par formes canoniques est plus sûre et toujours acceptée.
\end{methode}

\begin{attention}[L'erreur classique : le centre sans la division par 2]
Beaucoup d'élèves, voyant l'équation $x^2+y^2+z^2-2x+4y-6z+10=0$, annoncent un centre $(-2\,;\,4\,;\,-6)$ ou $(2\,;\,-4\,;\,6)$. C'est faux ! Les coordonnées du centre sont les \textbf{opposés des moitiés} des coefficients de $x$, $y$, $z$ :
\[
\Omega\left(-\frac{a}{2}\,;\,-\frac{b}{2}\,;\,-\frac{c}{2}\right).
\]
Ici $a=-2$, $b=4$, $c=-6$, donc $\Omega(1\,;\,-2\,;\,3)$. Retiens l'astuce : dans $(x-x_0)^2$, le développement donne $-2x_0x$ ; le coefficient de $x$ vaut donc $-2x_0$, d'où le facteur $\frac{1}{2}$ et le changement de signe.
\end{attention}

\courssub{Trois exemples chiffrés, un pour chaque cas}

\begin{exemplebox}[Cas $k>0$ : une sphère]
Montrons que l'ensemble $(E)$ d'équation $x^2+y^2+z^2-2x+4y-6z+10=0$ est la sphère de centre $\Omega(1\,;\,-2\,;\,3)$ et de rayon $2$.

On complète les carrés :
\begin{align*}
x^2-2x &= (x-1)^2-1,\\
y^2+4y &= (y+2)^2-4,\\
z^2-6z &= (z-3)^2-9.
\end{align*}
L'équation équivaut donc à
\begin{align*}
(x-1)^2-1+(y+2)^2-4+(z-3)^2-9+10 &= 0\\
(x-1)^2+(y+2)^2+(z-3)^2 &= 4.
\end{align*}
On reconnaît $\Omega M^2=4$ avec $\Omega(1\,;\,-2\,;\,3)$. Comme $k=4>0$, $(E)$ est la sphère de centre $\Omega(1\,;\,-2\,;\,3)$ et de rayon $R=\sqrt{4}=2$.

Vérification par la formule : $k=\dfrac{(-2)^2+4^2+(-6)^2}{4}-10=\dfrac{4+16+36}{4}-10=14-10=4>0$. $\checkmark$
\end{exemplebox}

\begin{exemplebox}[Cas $k=0$ : un singleton]
Soit $(E)$ l'ensemble d'équation $x^2+y^2+z^2+2x-2y+4z+6=0$. On complète :
\[
x^2+2x=(x+1)^2-1, \qquad y^2-2y=(y-1)^2-1, \qquad z^2+4z=(z+2)^2-4.
\]
L'équation équivaut à
\[
(x+1)^2+(y-1)^2+(z+2)^2 = 1+1+4-6 = 0.
\]
Une somme de carrés de réels est nulle si et seulement si chaque carré est nul : $x=-1$, $y=1$, $z=-2$. L'ensemble $(E)$ est donc réduit au seul point $\Omega(-1\,;\,1\,;\,-2)$. On dit parfois, par abus de langage, que c'est une « sphère de rayon nul ».
\end{exemplebox}

\begin{exemplebox}[Cas $k<0$ : l'ensemble vide]
Soit $(E)$ l'ensemble d'équation $x^2+y^2+z^2-4x+2y+8=0$ (ici $c=0$ : il n'y a pas de terme en $z$). On complète :
\[
x^2-4x=(x-2)^2-4, \qquad y^2+2y=(y+1)^2-1, \qquad z^2=(z-0)^2.
\]
L'équation équivaut à
\[
(x-2)^2+(y+1)^2+z^2 = 4+1-8 = -3.
\]
Or une somme de carrés de réels est toujours positive ou nulle : elle ne peut valoir $-3$. Donc $(E)=\varnothing$ : aucun point de l'espace ne vérifie cette équation. Vérification : $k=\dfrac{16+4+0}{4}-8=5-8=-3<0$. $\checkmark$
\end{exemplebox}

\begin{exoresolu}[Bilan : reconnaître et caractériser]
Énoncé. Dans un repère orthonormé de l'espace, on considère l'ensemble $(\mathcal{S})$ des points $M(x\,;\,y\,;\,z)$ tels que
\[
x^2+y^2+z^2-6x+2y-4z+13=0.
\]
Déterminer la nature de $(\mathcal{S})$ et préciser ses éléments caractéristiques.
\tcblower
Solution détaillée.

\textbf{Forme de l'équation.} Les coefficients de $x^2$, $y^2$, $z^2$ sont tous égaux à $1$ et il n'y a aucun terme en $xy$, $yz$ ou $zx$ : on est bien dans le cadre du théorème.

\textbf{Formes canoniques.}
\begin{align*}
x^2-6x &= (x-3)^2-9,\\
y^2+2y &= (y+1)^2-1,\\
z^2-4z &= (z-2)^2-4.
\end{align*}

\textbf{Réécriture.} L'équation équivaut à
\[
(x-3)^2+(y+1)^2+(z-2)^2 = 9+1+4-13 = 1.
\]

\textbf{Conclusion.} Comme $k=1>0$, $(\mathcal{S})$ est la sphère de centre $\Omega(3\,;\,-1\,;\,2)$ et de rayon $R=\sqrt{1}=1$.

Contrôle rapide par la formule : $k=\dfrac{(-6)^2+2^2+(-4)^2}{4}-13=\dfrac{36+4+16}{4}-13=14-13=1$. $\checkmark$
\end{exoresolu}

\begin{savaistu}[Pourquoi cette discussion est-elle si utile ?]
La même idée — compléter les carrés pour reconnaître une géométrie cachée derrière une équation — fonctionne dans tout le programme : en Seconde pour la parabole $y=ax^2+bx+c$, en Première pour le cercle dans le plan ($x^2+y^2+ax+by+c=0$), et ici pour la sphère. C'est aussi le principe utilisé par les logiciels de géolocalisation : lorsqu'un récepteur GPS (par exemple celui d'un téléphone sur le réseau MTN ou Orange à Yaoundé) reçoit les signaux de plusieurs satellites, chaque signal se traduit par une équation de sphère centrée sur un satellite, et le récepteur doit reconnaître et exploiter ces équations pour calculer sa position. Comprendre la discussion selon le signe de $k$, c'est comprendre pourquoi certaines configurations ne donnent aucune position possible !
\end{savaistu}

\begin{exercice}[À toi de jouer]
Dans chaque cas, déterminer la nature de l'ensemble des points $M(x\,;\,y\,;\,z)$ et préciser ses éléments caractéristiques :
\begin{enumerate}
\item $x^2+y^2+z^2+4x-6y+2z+5=0$ ;
\item $x^2+y^2+z^2-2x+10y+26=0$ ;
\item $x^2+y^2+z^2-8x+4z+20=0$.
\end{enumerate}
\end{exercice}

\begin{corrige}[Exercice « À toi de jouer »]
\textbf{1.} On complète : $x^2+4x=(x+2)^2-4$, $y^2-6y=(y-3)^2-9$, $z^2+2z=(z+1)^2-1$. L'équation équivaut à
\[
(x+2)^2+(y-3)^2+(z+1)^2 = 4+9+1-5 = 9.
\]
Comme $k=9>0$, c'est la sphère de centre $\Omega(-2\,;\,3\,;\,-1)$ et de rayon $R=3$.

\textbf{2.} On complète : $x^2-2x=(x-1)^2-1$, $y^2+10y=(y+5)^2-25$, $z^2=z^2$. L'équation équivaut à
\[
(x-1)^2+(y+5)^2+z^2 = 1+25-26 = 0.
\]
Comme $k=0$, l'ensemble est réduit au singleton $\{\Omega\}$ avec $\Omega(1\,;\,-5\,;\,0)$.

\textbf{3.} On complète : $x^2-8x=(x-4)^2-16$, $z^2+4z=(z+2)^2-4$. L'équation équivaut à
\[
(x-4)^2+y^2+(z+2)^2 = 16+4-20 = 0.
\]
Comme $k=0$, l'ensemble est réduit au singleton $\{\Omega\}$ avec $\Omega(4\,;\,0\,;\,-2)$. (Attention : ici $b=0$ et $c=4$ ; ne pas confondre les rôles des variables.)
\end{corrige}

En résumé, tu sais désormais passer de l'équation développée $x^2+y^2+z^2+ax+by+cz+d=0$ à la forme géométrique $(x-x_0)^2+(y-y_0)^2+(z-z_0)^2=k$, et conclure sans hésiter sur la nature de l'ensemble. Cette compétence sera notre outil principal dans les sections suivantes, lorsqu'il s'agira d'étudier l'intersection d'une sphère avec une droite ou avec un plan : la plupart des problèmes se ramèneront, après calcul, à reconnaître une équation de ce type.

\coursec{Position relative d'une sphère et d'une droite}

Dans la section précédente, nous avons appris à reconnaître, à partir d'une équation de la forme $x^2+y^2+z^2+ax+by+cz+d=0$, si l'ensemble des points correspondants est une sphère, un point ou l'ensemble vide. Nous savons donc désormais manipuler une sphère $\mathcal{S}(\Omega,R)$ par son équation. Il est temps de poser une question naturelle de géométrie : si l'on lance une droite dans l'espace, que peut-il se passer avec cette sphère ?

Pense à une balle de football suspendue dans l'air et à un fil tendu (une corde à linge, par exemple) : soit le fil ne touche pas la balle, soit il l'effleure en un seul point, soit il la traverse en entrant d'un côté et en ressortant de l'autre. Il n'y a pas d'autre possibilité. Cette intuition est exactement le contenu mathématique de cette section.

\courssub{Rappel : représentation paramétrique d'une droite de l'espace}

\begin{definition}[Représentation paramétrique d'une droite]
L'espace est muni d'un repère orthonormé $(O\,;\vec{\imath},\vec{\jmath},\vec{k})$. Soit $(D)$ la droite passant par le point $A(x_A\,;y_A\,;z_A)$ et de vecteur directeur $\vec{u}(\alpha\,;\beta\,;\gamma)$, avec $\vec{u}\neq\vec{0}$. Un point $M(x\,;y\,;z)$ appartient à $(D)$ si et seulement s'il existe un réel $t$ tel que $\overrightarrow{AM}=t\,\vec{u}$, c'est-à-dire :
\[
\begin{cases} x = x_A + \alpha t \\ y = y_A + \beta t \\ z = z_A + \gamma t \end{cases} \qquad (t\in\mathbb{R})
\]
Ce système s'appelle une \cle{représentation paramétrique} de $(D)$, et $t$ est le \cle{paramètre}.
\end{definition}

\begin{aretenir}[Lecture géométrique]
Quand $t$ parcourt $\mathbb{R}$, le point $M$ parcourt toute la droite $(D)$ : pour $t=0$ on obtient le point $A$, et chaque valeur de $t$ donne un point et un seul de $(D)$. C'est cette correspondance qui va nous permettre de transformer un problème de géométrie dans l'espace en une simple équation à une inconnue.
\end{aretenir}

\courssub{Distance du centre de la sphère à la droite}

Toute la discussion repose sur un nombre : la distance du centre $\Omega$ de la sphère à la droite $(D)$, notée $d(\Omega,(D))$. C'est la longueur du segment $[\Omega H]$, où $H$ est le projeté orthogonal de $\Omega$ sur $(D)$.

\begin{propriete}[Distance d'un point à une droite — admise]
Soit $(D)$ la droite passant par $A$ et de vecteur directeur $\vec{u}$, et $\Omega$ un point de l'espace. Alors :
\[
d\big(\Omega,(D)\big) = \frac{\left\|\overrightarrow{\Omega A}\wedge\vec{u}\,\right\|}{\|\vec{u}\|}
\]
où $\overrightarrow{\Omega A}\wedge\vec{u}$ désigne le produit vectoriel de $\overrightarrow{\Omega A}$ par $\vec{u}$. Cette formule est admise ; elle est utilisable directement dans les exercices.
\end{propriete}

\begin{exemplebox}[Calcul d'une distance point-droite]
Soit $\Omega(1\,;1\,;1)$ et $(D)$ la droite passant par $A(1\,;0\,;-1)$ de vecteur directeur $\vec{u}(1\,;2\,;1)$. On a $\overrightarrow{\Omega A}(0\,;-1\,;-2)$, d'où :
\[
\overrightarrow{\Omega A}\wedge\vec{u} = \begin{pmatrix} (-1)(1)-(-2)(2) \\ (-2)(1)-(0)(1) \\ (0)(2)-(-1)(1) \end{pmatrix} = \begin{pmatrix} 3 \\ -2 \\ 1 \end{pmatrix}
\]
Donc $\left\|\overrightarrow{\Omega A}\wedge\vec{u}\right\|=\sqrt{9+4+1}=\sqrt{14}$ et $\|\vec{u}\|=\sqrt{1+4+1}=\sqrt{6}$. Finalement :
\[
d\big(\Omega,(D)\big)=\frac{\sqrt{14}}{\sqrt{6}}=\sqrt{\frac{7}{3}}.
\]
\end{exemplebox}

\courssub{Discussion géométrique : comparer $d$ et $R$}

Notons $d=d(\Omega,(D))$. Le point de $(D)$ le plus proche de $\Omega$ est le projeté orthogonal $H$, et $\Omega H=d$. Tout autre point $M$ de $(D)$ vérifie, par le théorème de Pythagore dans le triangle $\Omega HM$ rectangle en $H$ :
\[
\Omega M^2 = \Omega H^2 + HM^2 = d^2 + HM^2 \geq d^2.
\]
Ainsi, la plus petite distance possible entre $\Omega$ et un point de $(D)$ est exactement $d$. Tout se joue donc dans la comparaison de $d$ et du rayon $R$.

\begin{propriete}[Position relative d'une sphère et d'une droite — admise]
Soit $\mathcal{S}(\Omega,R)$ une sphère et $(D)$ une droite, $d=d(\Omega,(D))$.
\begin{itemize}
\item Si $d>R$ : $(D)$ et $\mathcal{S}$ n'ont \textbf{aucun point commun}. On dit que $(D)$ est \cle{extérieure} à la sphère.
\item Si $d=R$ : $(D)$ et $\mathcal{S}$ ont \textbf{un unique point commun}, qui est le projeté orthogonal $H$ de $\Omega$ sur $(D)$. On dit que $(D)$ est \cle{tangente} à la sphère en $H$.
\item Si $d<R$ : $(D)$ et $\mathcal{S}$ ont \textbf{exactement deux points communs}. On dit que $(D)$ est \cle{sécante} à la sphère.
\end{itemize}
\end{propriete}

\begin{popfigure}
\begin{tikzpicture}[scale=0.85]
% Cas 1 : droite extérieure
\begin{scope}[shift={(0,0)}]
\draw[thick,popblue] (0,0) circle (1);
\fill[poppink!40] (0,0) circle (1);
\draw[thick,poporange] (-1.6,1.7) -- (1.6,1.7);
\fill[popdark] (0,0) circle (1.5pt) node[below left=-2pt]{\scriptsize $\Omega$};
\draw[popgreen,thick] (0,0) -- (0,1.7) node[midway,right]{\scriptsize $d$};
\draw[popblue,thick] (0,0) -- (0.85,-0.53) node[midway,below]{\scriptsize $R$};
\node at (0,-1.6) {\scriptsize $d>R$ : extérieure};
\end{scope}
% Cas 2 : tangente
\begin{scope}[shift={(4.6,0)}]
\draw[thick,popblue] (0,0) circle (1);
\fill[poppink!40] (0,0) circle (1);
\draw[thick,poporange] (-1.6,1) -- (1.6,1);
\fill[popdark] (0,0) circle (1.5pt) node[below left=-2pt]{\scriptsize $\Omega$};
\fill[poporange] (0,1) circle (1.5pt) node[above]{\scriptsize $H$};
\draw[popgreen,thick] (0,0) -- (0,1) node[midway,right]{\scriptsize $d=R$};
\node at (0,-1.6) {\scriptsize $d=R$ : tangente};
\end{scope}
% Cas 3 : sécante
\begin{scope}[shift={(9.2,0)}]
\draw[thick,popblue] (0,0) circle (1);
\fill[poppink!40] (0,0) circle (1);
\draw[thick,poporange] (-1.6,0.45) -- (1.6,0.45);
\fill[popdark] (0,0) circle (1.5pt) node[below left=-2pt]{\scriptsize $\Omega$};
\fill[poporange] (-0.89,0.45) circle (1.5pt) (0.89,0.45) circle (1.5pt);
\draw[popgreen,thick] (0,0) -- (0,0.45) node[midway,right]{\scriptsize $d$};
\draw[popblue,thick] (0,0) -- (-0.53,-0.85) node[midway,left]{\scriptsize $R$};
\node at (0,-1.6) {\scriptsize $d<R$ : sécante};
\end{scope}
\end{tikzpicture}
\legende{Les trois positions relatives d'une droite et d'une sphère, vues en coupe par un plan contenant $\Omega$ et la droite.}
\end{popfigure}

\begin{definition}[Droite tangente à une sphère]
On dit qu'une droite $(D)$ est \cle{tangente} à une sphère $\mathcal{S}$ en un point $M$ lorsque $M$ est l'\textbf{unique} point commun de $(D)$ et de $\mathcal{S}$. De façon équivalente : $M\in(D)\cap\mathcal{S}$ et $(D)\perp(\Omega M)$, c'est-à-dire $\overrightarrow{\Omega M}\cdot\vec{u}=0$ où $\vec{u}$ dirige $(D)$.
\end{definition}

\begin{attention}[La tangence se vérifie par le calcul, pas par l'intuition plane]
Dans le plan, la tangente à un cercle en $M$ est unique : c'est la perpendiculaire au rayon en $M$. Dans l'espace, la situation est plus riche : par un point $M$ d'une sphère passent \textbf{une infinité} de droites tangentes (toutes les droites du plan tangent en $M$, que nous étudierons à la section suivante). Une droite tangente est bien perpendiculaire au rayon $[\Omega M]$, mais attention : une droite qui \emph{semble} perpendiculaire à un rayon sur un dessin peut très bien couper la sphère en un second point. Seul le calcul — distance $d=R$, discriminant nul, ou produit scalaire $\overrightarrow{\Omega M}\cdot\vec{u}=0$ — permet de conclure. Ne conclus jamais à partir d'une figure.
\end{attention}

\courssub{Méthode analytique : substitution et discriminant}

L'approche par la distance est rapide, mais elle ne donne pas les coordonnées des points d'intersection. Pour les obtenir, on utilise la représentation paramétrique.

\begin{methode}[Intersection d'une sphère et d'une droite par le calcul]
Pour déterminer l'intersection de la sphère $\mathcal{S}$ d'équation $(x-x_\Omega)^2+(y-y_\Omega)^2+(z-z_\Omega)^2=R^2$ et de la droite $(D)$ de représentation paramétrique $x=x_A+\alpha t$, $y=y_A+\beta t$, $z=z_A+\gamma t$ :
\begin{enumerate}
\item \textbf{Substituer} : remplacer $x$, $y$, $z$ dans l'équation de la sphère par leurs expressions en fonction de $t$.
\item \textbf{Développer et ordonner} : on obtient une équation du second degré en $t$, de la forme $At^2+Bt+C=0$ (on a toujours $A=\alpha^2+\beta^2+\gamma^2=\|\vec{u}\|^2>0$).
\item \textbf{Discuter selon le discriminant} $\Delta=B^2-4AC$ :
\begin{itemize}
\item $\Delta<0$ : aucun point d'intersection (droite extérieure) ;
\item $\Delta=0$ : un unique point d'intersection (tangence), obtenu pour $t_0=-\dfrac{B}{2A}$ ;
\item $\Delta>0$ : deux points d'intersection, obtenus pour $t_1=\dfrac{-B-\sqrt{\Delta}}{2A}$ et $t_2=\dfrac{-B+\sqrt{\Delta}}{2A}$.
\end{itemize}
\item \textbf{Conclure} : reporter chaque valeur de $t$ dans la représentation paramétrique pour obtenir les coordonnées des points.
\end{enumerate}
\end{methode}

\begin{propriete}[Lien entre les deux approches]
Les deux méthodes disent exactement la même chose : le signe de $\Delta$ traduit algébriquement la comparaison de $d$ et $R$. Plus précisément, en développant la substitution on montre que
\[
\Delta = 4\,\|\vec{u}\|^2\big(R^2-d^2\big).
\]
Comme $\|\vec{u}\|^2>0$, on a bien : $\Delta<0\Leftrightarrow d>R$, $\Delta=0\Leftrightarrow d=R$, $\Delta>0\Leftrightarrow d<R$. Choisir l'une ou l'autre méthode dépend de la question posée : la distance pour discuter rapidement la position, le discriminant pour obtenir les points.
\end{propriete}

\begin{exemplebox}[Deux points symétriques]
Déterminons l'intersection de la sphère $\mathcal{S}$ d'équation $x^2+y^2+z^2=9$ (centre $O$, rayon $3$) et de la droite $(D)$ : $x=t$, $y=t$, $z=t$.

En substituant : $t^2+t^2+t^2=9$, soit $3t^2=9$, c'est-à-dire $t^2=3$, donc $t=-\sqrt{3}$ ou $t=\sqrt{3}$. Comme $\Delta>0$ (deux solutions), la droite est sécante et les points d'intersection sont :
\[
M_1\big(-\sqrt{3}\,;-\sqrt{3}\,;-\sqrt{3}\big) \quad\text{et}\quad M_2\big(\sqrt{3}\,;\sqrt{3}\,;\sqrt{3}\big).
\]
On remarque que $M_1$ et $M_2$ sont symétriques par rapport au centre $O$ : la droite $(D)$ passe par le centre ($d=0<R$), donc $[M_1M_2]$ est un \textbf{diamètre} de la sphère, de longueur $2R=6$.
\end{exemplebox}

\begin{exoresolu}[Étude complète d'une intersection]
Soit $\mathcal{S}$ la sphère de centre $\Omega(1\,;1\,;1)$ et de rayon $R=3$, et $(D)$ la droite de représentation paramétrique :
\[
\begin{cases} x = 1+t \\ y = 2t \\ z = -1+t \end{cases} \qquad (t\in\mathbb{R})
\]
Déterminer la position relative de $\mathcal{S}$ et $(D)$, puis les coordonnées des points d'intersection éventuels.
\tcblower
\textbf{Étape 1 : équation de la sphère.} $\mathcal{S}$ a pour équation :
\[
(x-1)^2+(y-1)^2+(z-1)^2=9.
\]
\textbf{Étape 2 : substitution.} On remplace $x=1+t$, $y=2t$, $z=-1+t$ :
\[
(1+t-1)^2+(2t-1)^2+(-1+t-1)^2=9
\]
\[
t^2+(2t-1)^2+(t-2)^2=9.
\]
\textbf{Étape 3 : développement.}
\begin{align*}
t^2+(4t^2-4t+1)+(t^2-4t+4) &= 9\\
6t^2-8t+5 &= 9\\
6t^2-8t-4 &= 0\\
3t^2-4t-2 &= 0.
\end{align*}
\textbf{Étape 4 : discriminant.} $\Delta=(-4)^2-4\times 3\times(-2)=16+24=40>0$. La droite est donc \textbf{sécante} à la sphère en deux points.
\[
t_1=\frac{4-\sqrt{40}}{6}=\frac{2-\sqrt{10}}{3}, \qquad t_2=\frac{2+\sqrt{10}}{3}.
\]
\textbf{Étape 5 : coordonnées des points.} En reportant dans la représentation paramétrique :
\[
M_1\left(\frac{5-\sqrt{10}}{3}\,;\,\frac{4-2\sqrt{10}}{3}\,;\,\frac{-1-\sqrt{10}}{3}\right), \quad
M_2\left(\frac{5+\sqrt{10}}{3}\,;\,\frac{4+2\sqrt{10}}{3}\,;\,\frac{-1+\sqrt{10}}{3}\right).
\]
\textbf{Vérification par la distance.} On a calculé plus haut $d(\Omega,(D))=\sqrt{\dfrac{7}{3}}\approx 1{,}53<3=R$ : la droite est bien sécante, ce qui confirme le signe de $\Delta$. Les deux approches sont cohérentes.
\end{exoresolu}

\begin{attention}[Pièges fréquents]
\begin{itemize}
\item Ne pas oublier de \textbf{développer complètement} les carrés avant de calculer $\Delta$ : une erreur de signe dans $(2t-1)^2$ fausse toute la discussion.
\item Le coefficient de $t^2$ est $\alpha^2+\beta^2+\gamma^2$ : il n'est \textbf{jamais nul} puisque $\vec{u}\neq\vec{0}$. L'équation en $t$ est donc toujours un vrai trinôme du second degré.
\item Si $\Delta=0$, le point de tangence s'obtient avec $t_0=-\dfrac{B}{2A}$ : n'oublie pas de reporter cette valeur dans la représentation paramétrique pour donner les coordonnées du point, et non la valeur de $t$ comme réponse finale.
\item Une valeur de $t$ n'est pas un point : $t$ est un paramètre, le point $M$ s'obtient en calculant $x$, $y$ et $z$.
\end{itemize}
\end{attention}

\begin{savaistu}[Pourquoi cette méthode est partout]
La technique « substituer puis discuter le discriminant » est l'une des plus puissantes de la géométrie analytique. C'est exactement elle qui permet, dans les jeux vidéo et les logiciels d'imagerie, de calculer si un rayon lumineux (une droite) frappe un objet sphérique : on parle de \emph{lancer de rayons} (\emph{ray tracing}). Le même calcul décide si la trajectoire d'un satellite frôle l'atmosphère terrestre ou si un câble tendu entre deux pylônes de la SONATREL risque de toucher un réservoir sphérique de stockage. Un simple discriminant suffit à trancher.
\end{savaistu}

\begin{exercice}[À toi de jouer]
Soit $\mathcal{S}$ la sphère de centre $\Omega(0\,;1\,;2)$ et de rayon $R=2$, et $(D)$ la droite passant par $A(1\,;0\,;0)$ de vecteur directeur $\vec{u}(1\,;1\,;0)$.
\begin{enumerate}
\item Écrire une représentation paramétrique de $(D)$ et l'équation cartésienne de $\mathcal{S}$.
\item Déterminer la position relative de $(D)$ et $\mathcal{S}$ par la méthode du discriminant.
\item Retrouver le résultat en calculant $d(\Omega,(D))$.
\end{enumerate}
\end{exercice}

\begin{corrige}[Exercice 1]
1. $(D)$ : $x=1+t$, $y=t$, $z=0$ $(t\in\mathbb{R})$. Sphère : $x^2+(y-1)^2+(z-2)^2=4$.

2. Substitution : $(1+t)^2+(t-1)^2+4=4$, soit $(1+t)^2+(t-1)^2=0$, c'est-à-dire $2t^2+2=0$, ou encore $t^2+1=0$. Le discriminant est $\Delta=-4<0$ : la droite est \textbf{extérieure} à la sphère, aucun point commun.

3. $\overrightarrow{\Omega A}(1\,;-1\,;-2)$, $\overrightarrow{\Omega A}\wedge\vec{u}=\begin{pmatrix} (-1)(0)-(-2)(1) \\ (-2)(1)-(1)(0) \\ (1)(1)-(-1)(1) \end{pmatrix}=\begin{pmatrix} 2 \\ -2 \\ 2 \end{pmatrix}$, de norme $2\sqrt{3}$, et $\|\vec{u}\|=\sqrt{2}$. Donc $d=\dfrac{2\sqrt{3}}{\sqrt{2}}=\sqrt{6}\approx 2{,}45>2=R$ : la droite est bien extérieure, ce qui confirme $\Delta<0$.
\end{corrige}

\coursec{Position relative d'une sphère et d'un plan ; plan tangent}

Dans la section précédente, nous avons étudié l'intersection d'une sphère et d'une \emph{droite} : selon la distance du centre à la droite, on obtient deux points, un seul point (droite tangente) ou aucun point. Nous allons maintenant remplacer la droite par un \emph{plan}. L'idée directrice sera exactement la même : tout repose sur la comparaison entre le rayon $R$ de la sphère et la distance de son centre au plan. Mais la géométrie va s'enrichir : au lieu d'un point de tangence isolé, nous découvrirons un magnifique \emph{cercle} d'intersection, et la notion de \emph{plan tangent}, qui généralise en dimension 3 la tangente à un cercle.

Pense à une balle de football posée sur le terrain du stade Ahmadou Ahidjo : le sol (un plan) touche la balle en un seul point. Coupe maintenant la balle par une plaque de verre placée à mi-hauteur : la plaque rencontre la balle selon un cercle. Soulève la plaque au-dessus de la balle : plus aucun contact. Ces trois images, c'est toute la section.

\courssub{Distance d'un point à un plan : l'outil fondamental}

\begin{aretenir}[Rappel : distance d'un point à un plan]
L'espace est muni d'un repère orthonormé $(O\,;\vec{i},\vec{j},\vec{k})$. Soit $\Omega(a\,;b\,;c)$ un point et $\mathcal{P}$ un plan d'équation cartésienne $ux+vy+wz+h=0$, avec $(u,v,w)\neq(0,0,0)$. La distance de $\Omega$ au plan $\mathcal{P}$ est :
\[
d(\Omega,\mathcal{P})=\frac{|ua+vb+wc+h|}{\sqrt{u^{2}+v^{2}+w^{2}}}.
\]
C'est la longueur $\Omega H$, où $H$ est le \cle{projeté orthogonal} de $\Omega$ sur $\mathcal{P}$. De plus, $\overrightarrow{\Omega H}$ est colinéaire au vecteur normal $\vec{n}(u\,;v\,;w)$ du plan.
\end{aretenir}

\begin{attention}[Conditions d'application]
Cette formule n'est valable que dans un repère \cle{orthonormé}. Par ailleurs, n'oublie pas la valeur absolue au numérateur : une distance est toujours positive. Enfin, vérifie que le plan est bien donné sous la forme $ux+vy+wz+h=0$ (tout dans le membre de gauche) avant d'identifier $h$.
\end{attention}

\courssub{Théorème fondamental : les trois positions possibles}

\begin{propriete}[Théorème : position relative d'une sphère et d'un plan]
Soit $\mathcal{S}$ une sphère de centre $\Omega$ et de rayon $R$, $\mathcal{P}$ un plan, et $d=d(\Omega,\mathcal{P})$ la distance de $\Omega$ à $\mathcal{P}$. Trois cas se présentent :
\begin{enumerate}
\item Si $d>R$ : l'intersection $\mathcal{S}\cap\mathcal{P}$ est \textbf{vide}.
\item Si $d=R$ : l'intersection est réduite à \textbf{un seul point} $H$, projeté orthogonal de $\Omega$ sur $\mathcal{P}$. On dit que $\mathcal{P}$ est \cle{tangent} à $\mathcal{S}$ en $H$.
\item Si $d<R$ : l'intersection est un \textbf{cercle} contenu dans $\mathcal{P}$, de centre $H$ (projeté orthogonal de $\Omega$ sur $\mathcal{P}$) et de rayon $r=\sqrt{R^{2}-d^{2}}$.
\end{enumerate}
Ce théorème est au c\oe ur du chapitre ; sa démonstration est formatrice et souvent demandée au Probatoire.
\end{propriete}

\begin{experience}[Démonstration guidée du théorème]
Soit $H$ le projeté orthogonal de $\Omega$ sur $\mathcal{P}$, de sorte que $\Omega H=d$ et que la droite $(\Omega H)$ est perpendiculaire à $\mathcal{P}$.

\textbf{Étape 1 : une relation de Pythagore.} Prenons un point $M$ quelconque du plan $\mathcal{P}$. Si $M\neq H$, le triangle $\Omega HM$ est rectangle en $H$, car $(\Omega H)\perp\mathcal{P}$ et $(HM)\subset\mathcal{P}$. Le théorème de Pythagore donne :
\[
\Omega M^{2}=\Omega H^{2}+HM^{2}=d^{2}+HM^{2}.
\]
Cette égalité reste vraie si $M=H$ (alors $HM=0$).

\textbf{Étape 2 : traduction de l'appartenance à la sphère.} Un point $M\in\mathcal{P}$ appartient à $\mathcal{S}$ si et seulement si $\Omega M=R$, c'est-à-dire :
\[
HM^{2}=R^{2}-d^{2}.
\]

\textbf{Étape 3 : discussion selon le signe de $R^{2}-d^{2}$.}
\begin{itemize}
\item Si $d>R$, alors $R^{2}-d^{2}<0$ : l'équation $HM^{2}=R^{2}-d^{2}$ n'a aucune solution (un carré est positif). L'intersection est vide.
\item Si $d=R$, alors $HM^{2}=0$, donc $HM=0$ : le seul point commun est $M=H$. Le plan est tangent en $H$.
\item Si $d<R$, alors $HM=\sqrt{R^{2}-d^{2}}=r$ : les points communs sont les points de $\mathcal{P}$ situés à distance $r$ de $H$, c'est-à-dire le cercle de centre $H$ et de rayon $r=\sqrt{R^{2}-d^{2}}$, tracé dans le plan $\mathcal{P}$. $\blacksquare$
\end{itemize}
\end{experience}

\begin{popfigure}
\begin{center}
\begin{tikzpicture}[scale=1.05]
% Sphère (vue stylisée)
\shade[ball color=popblueL,opacity=0.55] (0,0) circle (2);
\draw[popblue,thick] (0,0) circle (2);
% Plan : parallélogramme
\fill[poporangeL,opacity=0.55] (-3.2,-0.9) -- (3.2,-0.9) -- (2.4,-1.7) -- (-2.4,-1.7) -- cycle;
\draw[poporange,thick] (-3.2,-0.9) -- (3.2,-0.9) -- (2.4,-1.7) -- (-2.4,-1.7) -- cycle;
\node[popdark] at (-2.9,-1.45) {$\mathcal{P}$};
% Cercle d'intersection (ellipse)
\draw[popdark,very thick] (0,-0.9) ellipse (1.55 and 0.42);
% Centre Oméga, projeté H, point M
\fill[popink] (0,0) circle (1.6pt) node[above left=1pt] {$\Omega$};
\fill[popink] (0,-0.9) circle (1.6pt) node[below left=1pt] {$H$};
\fill[popink] (1.55,-0.9) circle (1.6pt) node[right=1pt] {$M$};
% Segments
\draw[popgreen,very thick] (0,0) -- (0,-0.9) node[midway,right=2pt] {$d$};
\draw[popgreen,very thick] (0,-0.9) -- (1.55,-0.9) node[midway,below=2pt] {$r$};
\draw[popblue,very thick] (0,0) -- (1.55,-0.9) node[midway,above right=1pt] {$R$};
% Angle droit
\draw[popdark] (0.14,-0.9) -- (0.14,-0.76) -- (0,-0.76);
\end{tikzpicture}
\end{center}
\legende{Sphère coupée par un plan $\mathcal{P}$ avec $d<R$ : l'intersection est un cercle de centre $H$. Le triangle $\Omega HM$ est rectangle en $H$, d'où $\Omega M^{2}=d^{2}+r^{2}$, soit $R^{2}=d^{2}+r^{2}$.}
\end{popfigure}

\begin{methode}[Étudier la position relative d'une sphère et d'un plan]
\begin{enumerate}
\item Identifier le centre $\Omega$ et le rayon $R$ de la sphère (compléter les carrés si l'équation est donnée sous forme développée).
\item Calculer $d=d(\Omega,\mathcal{P})$ avec la formule de la distance point-plan.
\item Comparer $d$ et $R$ :
\begin{itemize}
\item $d>R$ : intersection vide ;
\item $d=R$ : plan tangent, le point de contact est le projeté orthogonal $H$ de $\Omega$ sur $\mathcal{P}$ ;
\item $d<R$ : intersection circulaire, de centre $H$ et de rayon $r=\sqrt{R^{2}-d^{2}}$.
\end{itemize}
\item Si besoin, déterminer $H$ en paramétrant la droite passant par $\Omega$ et dirigée par le vecteur normal $\vec{n}$ de $\mathcal{P}$, puis en écrivant que $H\in\mathcal{P}$.
\end{enumerate}
\end{methode}

\begin{exoresolu}[Un plan tangent : le plan $z=0$]
On considère la sphère $\mathcal{S}$ de centre $\Omega(0\,;0\,;2)$ et de rayon $R=2$, et le plan $\mathcal{P}$ d'équation $z=0$ (le plan du sol). Montrer que $\mathcal{P}$ est tangent à $\mathcal{S}$ et déterminer le point de contact.
\tcblower
\textbf{Solution.} Le plan $\mathcal{P}$ a pour équation $z=0$, c'est-à-dire $0x+0y+1z+0=0$ : un vecteur normal est $\vec{n}(0\,;0\,;1)$.

\textbf{Étape 1 : distance de $\Omega$ à $\mathcal{P}$.}
\[
d(\Omega,\mathcal{P})=\frac{|0\times 0+0\times 0+1\times 2+0|}{\sqrt{0^{2}+0^{2}+1^{2}}}=\frac{2}{1}=2.
\]

\textbf{Étape 2 : comparaison.} On a $d=2=R$ : d'après le théorème, le plan $\mathcal{P}$ est \textbf{tangent} à la sphère $\mathcal{S}$.

\textbf{Étape 3 : point de contact.} C'est le projeté orthogonal $H$ de $\Omega$ sur $\mathcal{P}$. La droite passant par $\Omega$ et dirigée par $\vec{n}$ a pour représentation paramétrique :
\[
\begin{cases} x=0 \\ y=0 \\ z=2+t \end{cases} \quad (t\in\mathbb{R}).
\]
Le point $H$ appartient à $\mathcal{P}$, donc sa cote est nulle : $2+t=0$, soit $t=-2$. On obtient $H(0\,;0\,;0)$ : le point de contact est l'origine du repère.

\textbf{Vérification :} $\Omega H=\sqrt{0+0+4}=2=R$, donc $H\in\mathcal{S}$, et $H\in\mathcal{P}$ par construction. Tout est cohérent : la sphère « repose » sur le plan $z=0$ exactement comme une boule posée sur une table.
\end{exoresolu}

\courssub{Plan tangent en un point de la sphère}

Dans le cas $d=R$, le point de contact est le projeté de $\Omega$. Renversons le point de vue : donnons-nous d'abord un point $A$ de la sphère, et cherchons le plan qui lui est tangent en $A$. L'intuition de la balle posée sur le sol se généralise : en chaque point d'une sphère, il existe un unique plan tangent, et il est perpendiculaire au rayon aboutissant à ce point — exactement comme la tangente à un cercle est perpendiculaire au rayon.

\begin{propriete}[Théorème : plan tangent en un point de la sphère]
Soit $\mathcal{S}$ une sphère de centre $\Omega$ et de rayon $R$, et $A$ un point de $\mathcal{S}$. Il existe un unique plan tangent à $\mathcal{S}$ en $A$ : c'est le plan passant par $A$ et de vecteur normal $\overrightarrow{\Omega A}$. Autrement dit, c'est l'ensemble des points $M$ tels que :
\[
\overrightarrow{\Omega A}\cdot\overrightarrow{AM}=0.
\]
En particulier, le plan tangent en $A$ est \cle{perpendiculaire} au rayon $[\Omega A]$.
\end{propriete}

\begin{experience}[Démonstration]
Notons $\mathcal{P}$ le plan passant par $A$ et de vecteur normal $\overrightarrow{\Omega A}$ : il est bien défini de manière unique car $\overrightarrow{\Omega A}\neq\vec{0}$ (puisque $A\neq\Omega$).

\textbf{Étape 1 : la droite $(\Omega A)$ est perpendiculaire à $\mathcal{P}$.} Par définition, le vecteur $\overrightarrow{\Omega A}$ est normal à $\mathcal{P}$, donc la droite $(\Omega A)$, dirigée par ce vecteur, est perpendiculaire au plan $\mathcal{P}$. Le projeté orthogonal de $\Omega$ sur $\mathcal{P}$ est donc le point d'intersection de $(\Omega A)$ avec $\mathcal{P}$, c'est-à-dire $A$ lui-même (car $A\in(\Omega A)$ et $A\in\mathcal{P}$).

\textbf{Étape 2 : calcul de la distance.} Il en résulte que $d(\Omega,\mathcal{P})=\Omega A=R$, puisque $A\in\mathcal{S}$.

\textbf{Étape 3 : conclusion.} D'après le théorème fondamental (cas $d=R$), le plan $\mathcal{P}$ est tangent à $\mathcal{S}$, et son point de contact est le projeté orthogonal de $\Omega$, c'est-à-dire $A$. Le plan tangent en $A$ est donc bien le plan passant par $A$ et orthogonal à $\overrightarrow{\Omega A}$, et il est unique. $\blacksquare$
\end{experience}

\begin{popfigure}
\begin{center}
\begin{tikzpicture}[scale=1.05]
% Sphère
\shade[ball color=popblueL,opacity=0.55] (0,0) circle (2);
\draw[popblue,thick] (0,0) circle (2);
% Plan tangent (parallélogramme incliné) au point A en haut à droite
\coordinate (A) at (1.2,1.6);
\fill[popgreenL,opacity=0.6] (0.2,2.35) -- (3.3,2.35) -- (2.5,1.0) -- (-0.6,1.0) -- cycle;
\draw[popgreen,thick] (0.2,2.35) -- (3.3,2.35) -- (2.5,1.0) -- (-0.6,1.0) -- cycle;
\node[popdark] at (3.0,1.25) {$\mathcal{P}$};
% Centre et point A
\fill[popink] (0,0) circle (1.6pt) node[below left=1pt] {$\Omega$};
\fill[popink] (A) circle (1.8pt) node[above right=0pt] {$A$};
% Rayon
\draw[poporange,very thick,->] (0,0) -- (A) node[midway,left=3pt] {$\overrightarrow{\Omega A}$};
% Angle droit entre le rayon et le plan
\draw[popdark] (1.06,1.42) -- (1.22,1.30) -- (1.36,1.46);
\end{tikzpicture}
\end{center}
\legende{Le plan tangent $\mathcal{P}$ en $A$ est perpendiculaire au rayon $[\Omega A]$ : le vecteur $\overrightarrow{\Omega A}$ est un vecteur normal de $\mathcal{P}$.}
\end{popfigure}

\begin{methode}[Déterminer une équation du plan tangent en $A$]
\begin{enumerate}
\item \textbf{Vérifier} que $A$ appartient à la sphère : calculer $\Omega A$ et contrôler que $\Omega A=R$ (étape indispensable !).
\item Calculer les coordonnées du vecteur $\overrightarrow{\Omega A}$ : c'est un vecteur normal au plan cherché.
\item Écrire la condition d'orthogonalité pour un point $M(x\,;y\,;z)$ quelconque :
\[
\overrightarrow{\Omega A}\cdot\overrightarrow{AM}=0.
\]
\item Développer le produit scalaire pour obtenir une équation cartésienne de la forme $ux+vy+wz+h=0$.
\end{enumerate}
\end{methode}

\begin{attention}[Piège classique : le point doit être sur la sphère]
Avant d'écrire l'équation du plan tangent en $A$, il faut \textbf{toujours vérifier que $A\in\mathcal{S}$}, c'est-à-dire que $\Omega A=R$. Si $A$ n'appartient pas à la sphère, la notion de « plan tangent en $A$ » n'a aucun sens, et la formule $\overrightarrow{\Omega A}\cdot\overrightarrow{AM}=0$ donne un plan qui n'est \emph{pas} tangent à la sphère. Au Probatoire, cette vérification est souvent la première question de l'exercice : ne la zappe pas, elle rapporte des points faciles.
\end{attention}

\begin{exemplebox}[Plan tangent à une sphère en un point donné]
Soit $\mathcal{S}$ la sphère de centre $\Omega(0\,;1\,;2)$ et de rayon $R=3$, et le point $A(2\,;2\,;2)$. Déterminons une équation du plan tangent à $\mathcal{S}$ en $A$.

\textbf{Étape 1 : vérification de l'appartenance.}
\[
\Omega A=\sqrt{(2-0)^{2}+(2-1)^{2}+(2-2)^{2}}=\sqrt{4+1+0}=\sqrt{5}.
\]
Or $\sqrt{5}\neq 3$... Attention, $A$ n'appartient pas à la sphère ! Reprenons avec le point $A(2\,;3\,;2)$ :
\[
\Omega A=\sqrt{2^{2}+2^{2}+0^{2}}=\sqrt{8}=2\sqrt{2}\neq 3.
\]
Cherchons un point qui convienne : prenons $A(2\,;1\,;2+\sqrt{5})$. Alors $\Omega A=\sqrt{4+0+5}=3=R$ : cette fois $A\in\mathcal{S}$.

\textbf{Étape 2 : vecteur normal.} $\overrightarrow{\Omega A}\begin{pmatrix} 2-0 \\ 1-1 \\ 2+\sqrt{5}-2 \end{pmatrix}=\begin{pmatrix} 2 \\ 0 \\ \sqrt{5} \end{pmatrix}$.

\textbf{Étape 3 : condition d'orthogonalité.} Pour $M(x\,;y\,;z)$ :
\[
\overrightarrow{\Omega A}\cdot\overrightarrow{AM}=0
\iff 2(x-2)+0(y-1)+\sqrt{5}\bigl(z-2-\sqrt{5}\bigr)=0.
\]

\textbf{Étape 4 : développement.}
\[
2x-4+\sqrt{5}\,z-2\sqrt{5}-5=0
\iff 2x+\sqrt{5}\,z-9-2\sqrt{5}=0.
\]
Le plan tangent à $\mathcal{S}$ en $A$ a pour équation $2x+\sqrt{5}\,z-9-2\sqrt{5}=0$. On vérifie au passage que $d(\Omega,\mathcal{P})=\dfrac{|0+\sqrt{5}\times 2-9-2\sqrt{5}|}{\sqrt{4+5}}=\dfrac{9}{3}=3=R$ : tout est cohérent.
\end{exemplebox}

\begin{exoresolu}[Intersection d'une sphère et d'un plan sécant]
Soit $\mathcal{S}$ la sphère de centre $\Omega(1\,;-1\,;2)$ et de rayon $R=5$, et $\mathcal{P}$ le plan d'équation $x-y+z-1=0$. Déterminer la nature de l'intersection $\mathcal{S}\cap\mathcal{P}$ et ses éléments caractéristiques.
\tcblower
\textbf{Solution.} \textbf{Étape 1 : distance.} Avec $\vec{n}(1\,;-1\,;1)$ normal à $\mathcal{P}$ :
\[
d(\Omega,\mathcal{P})=\frac{|1-(-1)+2-1|}{\sqrt{1+1+1}}=\frac{|3|}{\sqrt{3}}=\sqrt{3}.
\]

\textbf{Étape 2 : comparaison.} $d=\sqrt{3}<5=R$ : l'intersection est un \textbf{cercle} $\mathcal{C}$.

\textbf{Étape 3 : rayon du cercle.}
\[
r=\sqrt{R^{2}-d^{2}}=\sqrt{25-3}=\sqrt{22}.
\]

\textbf{Étape 4 : centre du cercle.} C'est le projeté orthogonal $H$ de $\Omega$ sur $\mathcal{P}$. La droite passant par $\Omega$ et dirigée par $\vec{n}$ s'écrit :
\[
\begin{cases} x=1+t \\ y=-1-t \\ z=2+t \end{cases}
\]
En injectant dans l'équation de $\mathcal{P}$ : $(1+t)-(-1-t)+(2+t)-1=0$, soit $3t+3=0$, d'où $t=-1$. Donc $H(0\,;0\,;1)$.

\textbf{Conclusion :} $\mathcal{S}\cap\mathcal{P}$ est le cercle du plan $\mathcal{P}$, de centre $H(0\,;0\,;1)$ et de rayon $\sqrt{22}$.
\end{exoresolu}

\begin{savaistu}[Des paraboles aux sphères : la géométrie des satellites]
La comparaison de $d$ et $R$ n'est pas qu'un exercice scolaire : elle gouverne les télécommunications. Une antenne parabolique capte les signaux d'un satellite en orbite ; la zone couverte au sol par un satellite est exactement l'intersection de la sphère « portée du signal » avec le plan tangent à la Terre — un cercle ! Les ingénieurs de CAMTEL et des opérateurs mobiles dimensionnent ainsi la couverture réseau : plus le satellite est haut, plus le « cercle de couverture » est grand. De même, en médecine, les coupes d'un scanner IRM sont des intersections de la sphère (ou du corps) par des plans parallèles : chaque image est un cercle ou un disque de la section.
\end{savaistu}

\begin{exercice}[À toi de jouer]
Soit $\mathcal{S}$ la sphère d'équation $x^{2}+y^{2}+z^{2}-2x+4y-6z+10=0$ et $\mathcal{P}$ le plan d'équation $2x+2y-z+3=0$.
\begin{enumerate}
\item Déterminer le centre $\Omega$ et le rayon $R$ de $\mathcal{S}$.
\item Étudier la position relative de $\mathcal{S}$ et $\mathcal{P}$ ; préciser les éléments caractéristiques de l'intersection le cas échéant.
\item Déterminer une équation du plan tangent à $\mathcal{S}$ au point $A(3\,;-2\,;4)$ après avoir vérifié que $A\in\mathcal{S}$.
\end{enumerate}
\end{exercice}

\begin{corrige}[Exercice]
\textbf{1.} On complète les carrés :
\[
(x-1)^{2}-1+(y+2)^{2}-4+(z-3)^{2}-9+10=0
\iff (x-1)^{2}+(y+2)^{2}+(z-3)^{2}=4.
\]
Donc $\Omega(1\,;-2\,;3)$ et $R=2$.

\textbf{2.} Distance de $\Omega$ à $\mathcal{P}$ :
\[
d=\frac{|2\times 1+2\times(-2)-3+3|}{\sqrt{4+4+1}}=\frac{|2-4|}{3}=\frac{2}{3}.
\]
Comme $d=\dfrac{2}{3}<2=R$, l'intersection est un cercle de rayon $r=\sqrt{4-\dfrac{4}{9}}=\sqrt{\dfrac{32}{9}}=\dfrac{4\sqrt{2}}{3}$. Son centre $H$ s'obtient par la droite $(1+2t\,;-2+2t\,;3-t)$ : en reportant dans $\mathcal{P}$, $2(1+2t)+2(-2+2t)-(3-t)+3=0$ donne $9t-2=0$, soit $t=\dfrac{2}{9}$. D'où $H\left(\dfrac{13}{9}\,;-\dfrac{14}{9}\,;\dfrac{25}{9}\right)$.

\textbf{3.} Vérification : $\Omega A=\sqrt{(3-1)^{2}+(-2+2)^{2}+(4-3)^{2}}=\sqrt{4+0+1}=\sqrt{5}\neq 2$. Le point $A$ n'appartient \textbf{pas} à la sphère : la question du plan tangent en $A$ est sans objet. (Cet exercice illustre précisément le piège signalé : la vérification préalable est obligatoire !) Si l'on prend à la place $A'(3\,;-2\,;3)$, alors $\Omega A'=\sqrt{4+0+0}=2=R$, donc $A'\in\mathcal{S}$, et le plan tangent en $A'$ a pour équation $\overrightarrow{\Omega A'}\cdot\overrightarrow{A'M}=0$, soit $2(x-3)=0$, c'est-à-dire $x=3$.
\end{corrige}

\coursec{Cercle d'intersection d'une sphère et d'un plan}

Dans la section précédente, nous avons classé les positions relatives d'une sphère $\mathcal{S}(\Omega, R)$ et d'un plan $\mathcal{P}$ selon la valeur de la distance $d = d(\Omega, \mathcal{P})$ comparée au rayon $R$ : si $d > R$, le plan ne rencontre pas la sphère ; si $d = R$, le plan est \emph{tangent} à la sphère et l'intersection est réduite à un unique point. Reste le cas le plus riche, celui où le plan \emph{traverse} réellement la sphère : $d < R$. C'est toute cette section qui lui est consacrée. L'intuition est simple : prends une orange et coupe-la avec un couteau bien plat. La face de coupe que tu obtiens est un disque, et son bord est un cercle. Ce cercle, c'est exactement l'intersection de la sphère (la surface de l'orange) avec le plan de coupe. Toute la question est : où se trouve son centre, et quel est son rayon ?

\courssub{Le théorème fondamental}

\begin{propriete}[Intersection d'une sphère et d'un plan sécant]
Soit $\mathcal{S}$ une sphère de centre $\Omega$ et de rayon $R$, et $\mathcal{P}$ un plan tel que $d = d(\Omega, \mathcal{P}) < R$. Alors l'intersection $\mathcal{S} \cap \mathcal{P}$ est un \cle{cercle} $\mathcal{C}$ contenu dans le plan $\mathcal{P}$, dont :
\begin{itemize}
\item le \cle{centre} est le point $H$, \cle{projeté orthogonal} de $\Omega$ sur le plan $\mathcal{P}$ ;
\item le \cle{rayon} est $r = \sqrt{R^{2} - d^{2}}$.
\end{itemize}
La démonstration est exigible : elle repose uniquement sur le théorème de Pythagore.
\end{propriete}

\begin{experience}[Démonstration guidée du théorème]
Notons $H$ le projeté orthogonal de $\Omega$ sur $\mathcal{P}$. Par définition de la distance d'un point à un plan, $\Omega H = d$ et la droite $(\Omega H)$ est perpendiculaire à $\mathcal{P}$, donc perpendiculaire à toute droite de $\mathcal{P}$ passant par $H$.

\textbf{Étape 1 : tout point de l'intersection est sur le cercle.} Soit $M$ un point de $\mathcal{S} \cap \mathcal{P}$. Comme $M \in \mathcal{P}$ et $H \in \mathcal{P}$, la droite $(HM)$ est contenue dans $\mathcal{P}$, donc $(\Omega H) \perp (HM)$ : le triangle $\Omega HM$ est \emph{rectangle en $H$}. Le théorème de Pythagore donne :
\[ \Omega M^{2} = \Omega H^{2} + HM^{2}. \]
Or $M \in \mathcal{S}$, donc $\Omega M = R$, et $\Omega H = d$. On en tire :
\[ HM^{2} = R^{2} - d^{2} \quad \text{donc} \quad HM = \sqrt{R^{2} - d^{2}} = r. \]
Ainsi tout point $M$ de l'intersection est dans le plan $\mathcal{P}$ et à distance $r$ de $H$ : il appartient au cercle de centre $H$ et de rayon $r$ tracé dans $\mathcal{P}$.

\textbf{Étape 2 : réciproque.} Soit $M$ un point du plan $\mathcal{P}$ tel que $HM = r$. Le triangle $\Omega HM$ est encore rectangle en $H$, donc :
\[ \Omega M^{2} = \Omega H^{2} + HM^{2} = d^{2} + (R^{2} - d^{2}) = R^{2}, \]
d'où $\Omega M = R$ : le point $M$ appartient à la sphère $\mathcal{S}$, donc à l'intersection.

\textbf{Conclusion.} L'intersection $\mathcal{S} \cap \mathcal{P}$ est exactement le cercle du plan $\mathcal{P}$ de centre $H$ et de rayon $r = \sqrt{R^{2} - d^{2}}$. $\blacksquare$
\end{experience}

\begin{popfigure}
\begin{center}
\begin{tikzpicture}[scale=1.05, line cap=round, line join=round]
% Sphère (grand cercle apparent)
\shade[ball color=popblueL, opacity=0.35] (0,0) circle (2.6);
\draw[popblue, thick] (0,0) circle (2.6);
% Plan sécant (parallélogramme en perspective)
\fill[poporangeL, opacity=0.55] (-3.4,0.35) -- (3.0,0.35) -- (3.8,1.75) -- (-2.6,1.75) -- cycle;
\draw[poporange, thick] (-3.4,0.35) -- (3.0,0.35) -- (3.8,1.75) -- (-2.6,1.75) -- cycle;
\node[poporange, anchor=west] at (3.8,1.75) {$\mathcal{P}$};
% Cercle d'intersection (ellipse)
\draw[popdark, very thick] (0.2,1.05) ellipse (2.15 and 0.62);
% Points Omega, H, M
\coordinate (O) at (0,0);
\coordinate (H) at (0.2,1.05);
\coordinate (M) at (2.35,1.05);
\fill[popink] (O) circle (1.6pt) node[below left=1pt] {$\Omega$};
\fill[popink] (H) circle (1.6pt) node[above left=1pt] {$H$};
\fill[popink] (M) circle (1.6pt) node[right=1pt] {$M$};
% Segments d, r, R
\draw[popgreen, very thick] (O) -- (H) node[midway, left=2pt] {$d$};
\draw[poporange, very thick] (H) -- (M) node[midway, above=2pt] {$r$};
\draw[popblue, very thick] (O) -- (M) node[midway, above right=1pt] {$R$};
% Angle droit en H
\draw[popdark] (0.2,0.83) -- (0.42,0.83) -- (0.42,1.05);
\node[poporange] at (-2.6,2.15) {$\mathcal{C} = \mathcal{S} \cap \mathcal{P}$};
\end{tikzpicture}
\end{center}
\legende{Le plan $\mathcal{P}$ coupe la sphère $\mathcal{S}(\Omega, R)$ selon le cercle $\mathcal{C}$ de centre $H$ (projeté orthogonal de $\Omega$ sur $\mathcal{P}$) et de rayon $r$. Le triangle $\Omega HM$ est rectangle en $H$.}
\end{popfigure}

\begin{popfigure}
\begin{center}
\begin{tikzpicture}[scale=1.5, line cap=round, line join=round]
% Vue en coupe : grand cercle
\draw[popblue, thick] (0,0) circle (2);
% Trace du plan (droite de coupe)
\draw[poporange, very thick] (-2.2,1.1) -- (2.2,1.1) node[right] {trace de $\mathcal{P}$};
% Cercle d'intersection vu en coupe = corde
\draw[popdark, very thick] (-1.66,1.1) -- (1.66,1.1);
% Points
\coordinate (O) at (0,0);
\coordinate (H) at (0,1.1);
\coordinate (M) at (1.66,1.1);
\fill[popink] (O) circle (1.4pt) node[below=2pt] {$\Omega$};
\fill[popink] (H) circle (1.4pt) node[above left=1pt] {$H$};
\fill[popink] (M) circle (1.4pt) node[above right=1pt] {$M$};
\draw[popgreen, very thick] (O) -- (H) node[midway, left=2pt] {$d$};
\draw[poporange, very thick] (H) -- (M) node[midway, below=2pt] {$r$};
\draw[popblue, very thick] (O) -- (M) node[midway, above right=1pt] {$R$};
\draw[popdark] (0,0.88) -- (0.22,0.88) -- (0.22,1.1);
\node[popdark] at (0,-2.45) {$r^{2} = R^{2} - d^{2}$};
\end{tikzpicture}
\end{center}
\legende{Vue en coupe par un plan contenant $\Omega$ et perpendiculaire à $\mathcal{P}$ : la sphère devient un cercle de rayon $R$, le cercle d'intersection devient une corde, et le théorème de Pythagore donne $r^{2} = R^{2} - d^{2}$.}
\end{popfigure}

\begin{attention}[Le centre du cercle n'est pas le centre de la sphère]
L'erreur la plus fréquente au Probatoire consiste à affirmer que le cercle d'intersection a pour centre $\Omega$. C'est \textbf{faux} en général : le centre du cercle est le point $H$, projeté orthogonal de $\Omega$ sur $\mathcal{P}$. Le centre du cercle coïncide avec $\Omega$ \emph{uniquement} lorsque le plan passe par $\Omega$ (c'est-à-dire $d = 0$). Retiens aussi que $r < R$ dès que $d > 0$ : tout cercle d'intersection est plus petit qu'un grand cercle.
\end{attention}

\courssub{Cas particulier : le grand cercle}

\begin{definition}[Grand cercle]
Lorsque le plan $\mathcal{P}$ passe par le centre $\Omega$ de la sphère (c'est-à-dire $d = 0$), le cercle d'intersection a pour centre $\Omega$ lui-même et pour rayon $r = \sqrt{R^{2} - 0} = R$. On l'appelle un \cle{grand cercle} de la sphère. C'est le plus grand cercle que l'on puisse tracer sur une sphère.
\end{definition}

\begin{savaistu}[Parallèles, méridiens et grands cercles]
La Terre est assimilée à une sphère de rayon environ \SI{6371}{km}. Les \emph{parallèles} sont les cercles d'intersection de la sphère terrestre avec des plans perpendiculaires à l'axe des pôles : parmi eux, seul l'\emph{équateur} est un grand cercle (son plan passe par le centre de la Terre) ; plus on se rapproche des pôles, plus $d$ augmente et plus le rayon $r = \sqrt{R^{2}-d^{2}}$ du parallèle diminue. Les \emph{méridiens}, eux, sont tous des grands cercles, car leurs plans contiennent l'axe des pôles, donc le centre de la Terre. C'est pourquoi les avions long-courrier suivent des routes « orthodromiques » : le plus court chemin entre deux points d'une sphère est l'arc du grand cercle qui les relie.
\end{savaistu}

\courssub{Méthode pratique et système caractéristique}

\begin{methode}[Déterminer le cercle d'intersection d'une sphère et d'un plan]
Soit $\mathcal{S}(\Omega, R)$ et $\mathcal{P} : ax + by + cz + e = 0$.
\begin{enumerate}
\item \textbf{Calculer la distance} $d = d(\Omega, \mathcal{P}) = \dfrac{|a x_{\Omega} + b y_{\Omega} + c z_{\Omega} + e|}{\sqrt{a^{2}+b^{2}+c^{2}}}$.
\item \textbf{Comparer $d$ et $R$} : si $d > R$, l'intersection est vide ; si $d = R$, c'est le point de tangence ; si $d < R$, c'est un cercle et on continue.
\item \textbf{Déterminer les éléments du cercle} :
\begin{itemize}
\item le centre $H$ : on écrit la représentation paramétrique de la droite $\mathcal{D}$ passant par $\Omega$ et dirigée par le vecteur normal $\vec{n}(a\,; b\,; c)$ de $\mathcal{P}$, puis on résout $\mathcal{D} \cap \mathcal{P} = \{H\}$ ;
\item le rayon : $r = \sqrt{R^{2} - d^{2}}$.
\end{itemize}
\end{enumerate}
\end{methode}

\begin{aretenir}[Système caractéristique du cercle]
Dans l'espace, un cercle ne peut pas être décrit par une seule équation cartésienne. Le cercle d'intersection $\mathcal{C} = \mathcal{S} \cap \mathcal{P}$ est caractérisé par le \cle{système} formé des deux équations :
\[ \mathcal{C} : \begin{cases} (x - x_{\Omega})^{2} + (y - y_{\Omega})^{2} + (z - z_{\Omega})^{2} = R^{2} \\ ax + by + cz + e = 0 \end{cases} \]
Un point appartient au cercle si et seulement si ses coordonnées vérifient \emph{simultanément} les deux équations. Pour décrire complètement le cercle, on donne en plus son plan $\mathcal{P}$, son centre $H$ et son rayon $r$.
\end{aretenir}

\courssub{Exemples résolus}

\begin{exemplebox}[Sphère centrée à l'origine et plan $3x + 4z = 0$]
Soit $\mathcal{S}$ la sphère d'équation $x^{2} + y^{2} + z^{2} = 25$ et $\mathcal{P}$ le plan d'équation $3x + 4z = 0$. La sphère a pour centre $\Omega(0\,; 0\,; 0)$ et pour rayon $R = 5$. Comme les coordonnées de $\Omega$ vérifient $3 \times 0 + 4 \times 0 = 0$, le point $\Omega$ appartient à $\mathcal{P}$ : on a $d = 0$. Le plan passe par le centre de la sphère, donc l'intersection est un \textbf{grand cercle}, de centre $\Omega(0\,;0\,;0)$ et de rayon $r = R = 5$, tracé dans le plan $\mathcal{P}$.
\end{exemplebox}

\begin{exoresolu}[Cercle d'intersection : calcul complet]
Dans l'espace muni d'un repère orthonormé $(O\,; \vec{i}, \vec{j}, \vec{k})$, on considère la sphère $\mathcal{S}$ de centre $\Omega(1\,; 0\,; -1)$ et de rayon $R = 3$, et le plan $\mathcal{P}$ d'équation $x - y + z + 1 = 0$. Déterminer la nature de $\mathcal{S} \cap \mathcal{P}$ et, le cas échéant, ses éléments caractéristiques.
\tcblower
\textbf{Étape 1 : distance de $\Omega$ à $\mathcal{P}$.}
\[ d = \dfrac{|1 - 0 + (-1) + 1|}{\sqrt{1^{2} + (-1)^{2} + 1^{2}}} = \dfrac{|1|}{\sqrt{3}} = \dfrac{1}{\sqrt{3}} = \dfrac{\sqrt{3}}{3}. \]
\textbf{Étape 2 : comparaison.} On a $d = \dfrac{\sqrt{3}}{3} \approx 0{,}58 < 3 = R$ : l'intersection est un \textbf{cercle} $\mathcal{C}$.

\textbf{Étape 3 : centre $H$.} Le vecteur $\vec{n}(1\,; -1\,; 1)$ est normal à $\mathcal{P}$. La droite $\mathcal{D}$ passant par $\Omega$ et dirigée par $\vec{n}$ admet la représentation paramétrique :
\[ \begin{cases} x = 1 + t \\ y = -t \\ z = -1 + t \end{cases} \quad (t \in \mathbb{R}). \]
Le point $H$ est l'intersection de $\mathcal{D}$ et $\mathcal{P}$. On injecte dans l'équation du plan :
\[ (1 + t) - (-t) + (-1 + t) + 1 = 0 \iff 3t + 1 = 0 \iff t = -\dfrac{1}{3}. \]
D'où $H\left(\dfrac{2}{3}\,; \dfrac{1}{3}\,; -\dfrac{4}{3}\right)$.

\textbf{Étape 4 : rayon.}
\[ r = \sqrt{R^{2} - d^{2}} = \sqrt{9 - \dfrac{1}{3}} = \sqrt{\dfrac{26}{3}} = \dfrac{\sqrt{78}}{3} \approx 2{,}94. \]
\textbf{Conclusion.} $\mathcal{S} \cap \mathcal{P}$ est le cercle du plan $\mathcal{P}$, de centre $H\left(\dfrac{2}{3}\,; \dfrac{1}{3}\,; -\dfrac{4}{3}\right)$ et de rayon $r = \dfrac{\sqrt{78}}{3}$.
\end{exoresolu}

\begin{exoresolu}[Un plan qui frôle la sphère]
Soit $\mathcal{S}$ la sphère de centre $\Omega(0\,; 0\,; 0)$ et de rayon $R = 5$, et $\mathcal{P}$ le plan d'équation $x + y + z - 3 = 0$. Déterminer l'intersection de $\mathcal{S}$ et $\mathcal{P}$.
\tcblower
\textbf{Distance :} $d = \dfrac{|0 + 0 + 0 - 3|}{\sqrt{1^{2}+1^{2}+1^{2}}} = \dfrac{3}{\sqrt{3}} = \sqrt{3}$. Comme $\sqrt{3} < 5$, l'intersection est un cercle.

\textbf{Centre $H$ :} la droite passant par $\Omega$ et dirigée par $\vec{n}(1\,;1\,;1)$ a pour représentation paramétrique $x = t$, $y = t$, $z = t$. En injectant dans l'équation du plan : $3t - 3 = 0$, soit $t = 1$. Donc $H(1\,; 1\,; 1)$.

\textbf{Rayon :} $r = \sqrt{R^{2} - d^{2}} = \sqrt{25 - 3} = \sqrt{22}$.

\textbf{Système caractéristique :} $\mathcal{C} : \begin{cases} x^{2} + y^{2} + z^{2} = 25 \\ x + y + z - 3 = 0 \end{cases}$.
\end{exoresolu}

\begin{attention}[Pièges de rédaction]
\begin{itemize}
\item Ne conclus jamais « l'intersection est un cercle » sans avoir d'abord \emph{calculé $d$} et \emph{comparé à $R$} : cette comparaison est la justification exigée.
\item La distance $d$ se calcule avec une \emph{valeur absolue} au numérateur ; oublier la racine carrée au dénominateur ($\sqrt{a^{2}+b^{2}+c^{2}}$) est une erreur classique.
\item Pour trouver $H$, le paramètre $t$ obtenu doit être \emph{reporté} dans la représentation paramétrique : donner $t = -\dfrac{1}{3}$ n'est pas donner le point $H$.
\item Dans la formule $r = \sqrt{R^{2} - d^{2}}$, c'est bien $d^{2}$ que l'on soustrait : avec $d = \dfrac{1}{\sqrt{3}}$, on a $d^{2} = \dfrac{1}{3}$ et non $\dfrac{1}{9}$.
\end{itemize}
\end{attention}

\courssub{Pour aller plus loin : intersection de deux sphères}

\begin{savaistu}[Hors programme de Première — utile pour la Terminale]
Que se passe-t-il lorsqu'on intersecte \emph{deux} sphères $\mathcal{S}_{1}$ et $\mathcal{S}_{2}$ d'équations cartésiennes développées :
\[ \mathcal{S}_{1} : x^{2}+y^{2}+z^{2}+a_{1}x+b_{1}y+c_{1}z+d_{1} = 0, \qquad \mathcal{S}_{2} : x^{2}+y^{2}+z^{2}+a_{2}x+b_{2}y+c_{2}z+d_{2} = 0 \ ? \]
L'astuce consiste à \emph{soustraire membre à membre} les deux équations : les termes quadratiques $x^{2}+y^{2}+z^{2}$ disparaissent et il reste une équation \emph{du premier degré}, c'est-à-dire l'équation d'un \cle{plan}, appelé \emph{plan radical} des deux sphères. Les points communs aux deux sphères sont donc exactement les points communs à l'une des sphères et à ce plan : on est ramené au problème de cette section ! Selon la position des sphères, l'intersection est alors un cercle, un point (sphères tangentes) ou l'ensemble vide. C'est le même principe que la localisation GPS : la position d'un récepteur est déterminée comme intersection de sphères centrées sur les satellites.
\end{savaistu}

\begin{exercice}[À toi de jouer]
Soit $\mathcal{S}$ la sphère d'équation $x^{2} + y^{2} + z^{2} - 2x + 4y - 6z + 5 = 0$ et $\mathcal{P}$ le plan d'équation $x + 2y - 2z - 5 = 0$.
\begin{enumerate}
\item Déterminer le centre $\Omega$ et le rayon $R$ de $\mathcal{S}$.
\item Calculer $d(\Omega, \mathcal{P})$ et montrer que l'intersection est un cercle.
\item Déterminer le centre $H$ et le rayon $r$ de ce cercle, puis écrire son système caractéristique.
\end{enumerate}
\end{exercice}

\begin{corrige}[Exercice : sphère et plan sécant]
\textbf{1.} On complète les carrés : $(x-1)^{2} + (y+2)^{2} + (z-3)^{2} = 1 + 4 + 9 - 5 = 9$. Donc $\Omega(1\,; -2\,; 3)$ et $R = 3$.

\textbf{2.} $d = \dfrac{|1 + 2(-2) - 2(3) - 5|}{\sqrt{1+4+4}} = \dfrac{|-14|}{3} = \dfrac{14}{3}$. Mais $\dfrac{14}{3} > 3 = R$ : l'intersection est en fait \textbf{vide} ! Le plan ne rencontre pas la sphère. (Question piège : il fallait vérifier avant de se lancer dans les calculs.)

\textbf{3.} Sans objet puisque $\mathcal{S} \cap \mathcal{P} = \varnothing$. Moralité : la comparaison de $d$ et $R$ est \emph{toujours} la première chose à faire.
\end{corrige}

\coursec{Synthèse et applications}

Nous voici au terme de notre voyage autour de la sphère. Dans la section précédente, nous avons appris à déterminer le cercle d'intersection d'une sphère et d'un plan : centre projeté orthogonal du centre de la sphère, rayon donné par $r=\sqrt{R^{2}-d^{2}}$. Cette formule, tu vas le voir, est la clé de voûte de toute la leçon : elle régit aussi bien la corde découpée par une droite sécante que la zone de couverture d'une antenne relais au sol. Cette dernière section a un triple objectif : \textbf{organiser} tout ce que tu sais dans un tableau de synthèse, \textbf{dégager une méthode générale} applicable à tout problème de sphère, et \textbf{montrer la puissance} de ces outils sur des situations concrètes du quotidien camerounais.

\courssub{Tableau récapitulatif des positions relatives}

Toute la géométrie métrique de la sphère repose sur \emph{une seule idée} : comparer la distance $d$ (du centre $\Omega$ à la droite ou au plan) au rayon $R$ de la sphère. Le tableau ci-dessous condense les résultats démontrés dans les sections précédentes.

\begin{aretenir}[Positions relatives : le tableau à connaître par cœur]
Soit $\mathcal{S}$ une sphère de centre $\Omega$ et de rayon $R$, $\mathcal{D}$ une droite, $\mathcal{P}$ un plan. On note $d$ la distance de $\Omega$ à l'objet considéré.
\begin{center}
\rowcolors{1}{popcream}{white}
\begin{tabularx}{\linewidth}{|l|X|X|X|}
\hline
\rowcolor{popblueL} \textbf{Comparaison} & \textbf{Sphère / droite $\mathcal{D}$} & \textbf{Sphère / plan $\mathcal{P}$} & \textbf{Éléments caractéristiques} \\
\hline
$d>R$ & Intersection \textbf{vide} & Intersection \textbf{vide} & Aucun point ; la droite (ou le plan) est \emph{extérieure}. \\
\hline
$d=R$ & \textbf{Un seul point} : droite \emph{tangente} & \textbf{Un seul point} : plan \emph{tangent} & Le point de contact est le projeté orthogonal $H$ de $\Omega$ ; $(\Omega H)\perp\mathcal{D}$ (ou $\perp\mathcal{P}$). \\
\hline
$d<R$ & \textbf{Deux points} : droite \emph{sécante} & \textbf{Un cercle} : plan \emph{sécant} & Corde de longueur $2\sqrt{R^{2}-d^{2}}$ ; cercle de centre $H$ (projeté de $\Omega$) et de rayon $r=\sqrt{R^{2}-d^{2}}$. \\
\hline
\end{tabularx}
\end{center}
\end{aretenir}

\begin{popfigure}
\begin{tikzpicture}[scale=1, font=\small]
% Grille du tableau visuel
\draw[popline, thick] (0,0) rectangle (12,6.4);
\draw[popline] (0,4.6) -- (12,4.6);
\draw[popline] (0,2.4) -- (12,2.4);
\draw[popline] (4,0) -- (4,6.4);
\draw[popline] (8,0) -- (8,6.4);
% En-têtes
\node[popdark, font=\bfseries] at (2,5.9) {$d>R$};
\node[popdark, font=\bfseries] at (6,5.9) {$d=R$};
\node[popdark, font=\bfseries] at (10,5.9) {$d<R$};
\node[popdark, font=\bfseries, rotate=90] at (-0.55,5.5) {Droite};
\node[popdark, font=\bfseries, rotate=90] at (-0.55,3.4) {Plan};
\node[popdark, font=\bfseries, rotate=90] at (-0.55,1.2) {Intersection};
% Ligne droite : d>R
\draw[popblue, thick] (0.9,5.3) circle (0.45);
\draw[poporange, thick] (0.35,4.85) -- (1.6,4.85);
\node[popmuted] at (2.6,5.15) {$\varnothing$};
% Ligne droite : d=R
\draw[popblue, thick] (4.9,5.3) circle (0.45);
\draw[poporange, thick] (4.35,4.85) -- (5.6,4.85);
\fill[popink] (4.9,4.85) circle (0.05);
\node[popmuted] at (6.6,5.15) {1 point};
% Ligne droite : d<R
\draw[popblue, thick] (8.9,5.3) circle (0.45);
\draw[poporange, thick] (8.3,5.05) -- (9.5,5.05);
\fill[popink] (8.52,5.05) circle (0.05);
\fill[popink] (9.28,5.05) circle (0.05);
\node[popmuted] at (10.6,5.15) {2 points};
% Ligne plan : d>R
\draw[popblue, thick] (0.9,3.5) circle (0.45);
\draw[popgreen, thick] (0.35,2.85) -- (1.6,2.85);
\node[popmuted] at (2.6,3.2) {$\varnothing$};
% Ligne plan : d=R
\draw[popblue, thick] (4.9,3.5) circle (0.45);
\draw[popgreen, thick] (4.35,3.05) -- (5.6,3.05);
\fill[popink] (4.9,3.05) circle (0.05);
\node[popmuted] at (6.6,3.2) {1 point (tangent)};
% Ligne plan : d<R
\draw[popblue, thick] (8.9,3.5) circle (0.45);
\draw[popgreen, thick] (8.3,3.3) -- (9.5,3.3);
\draw[poppurple, thick] (8.9,3.3) ellipse (0.38 and 0.12);
\node[popmuted] at (10.6,3.2) {cercle};
% Ligne intersection : formules
\node[popdark] at (2,1.2) {extérieur};
\node[popdark] at (6,1.2) {$H=\mathrm{proj}_{\perp}(\Omega)$};
\node[popdark] at (10,1.55) {corde : $2\sqrt{R^{2}-d^{2}}$};
\node[popdark] at (10,0.85) {cercle : $r=\sqrt{R^{2}-d^{2}}$};
\end{tikzpicture}
\legende{Synthèse visuelle : la comparaison de $d$ et $R$ décide de tout.}
\end{popfigure}

\begin{attention}[Une conclusion complète ou rien]
Au Probatoire, une réponse du type « la droite coupe la sphère » est \textbf{insuffisante} et \textbf{coûte des points}. Tu dois toujours conclure en donnant :
\begin{itemize}
\item la \cle{nature} de l'intersection (vide, un point, deux points, un cercle) ;
\item ses \cle{éléments caractéristiques} : les coordonnées des points d'intersection pour une droite, ou le centre \emph{et} le rayon du cercle pour un plan.
\end{itemize}
Autre piège classique : pour le plan tangent, le point de contact n'est \emph{pas} un point quelconque du plan : c'est précisément le projeté orthogonal de $\Omega$ sur $\mathcal{P}$.
\end{attention}

\courssub{La méthode générale : trois réflexes qui résolvent tout}

Face à n'importe quel problème de sphère, la stratégie est toujours la même. Elle tient en trois verbes : \textbf{identifier}, \textbf{calculer}, \textbf{conclure}.

\begin{methode}[Fiche de synthèse : les 4 savoir-faire exigibles]
\begin{enumerate}
\item \textbf{Écrire l'équation d'une sphère} connue par son centre $\Omega(a\,;\,b\,;\,c)$ et son rayon $R$ :
\[ (x-a)^{2}+(y-b)^{2}+(z-c)^{2}=R^{2}. \]
Si la sphère est donnée par un diamètre $[AB]$, le centre est le milieu de $[AB]$ et $R=\dfrac{AB}{2}$.
\item \textbf{Identifier centre et rayon} à partir de $x^{2}+y^{2}+z^{2}+ax+by+cz+d=0$ : compléter les carrés pour obtenir
\[ \left(x+\tfrac{a}{2}\right)^{2}+\left(y+\tfrac{b}{2}\right)^{2}+\left(z+\tfrac{c}{2}\right)^{2}=\tfrac{a^{2}+b^{2}+c^{2}}{4}-d. \]
Le membre de droite décide : $>0$ sphère, $=0$ ensemble réduit à un point, $<0$ ensemble vide.
\item \textbf{Intersection sphère--droite} : paramétrer la droite $(x\,;\,y\,;\,z)=(x_{0}+\alpha t\,;\,y_{0}+\beta t\,;\,z_{0}+\gamma t)$, injecter dans l'équation de la sphère, résoudre l'équation du second degré en $t$. Le discriminant donne le nombre de points ; chaque solution $t$ donne un point.
\item \textbf{Intersection sphère--plan} : calculer $d=\dfrac{|aa_{\Omega}+bb_{\Omega}+cc_{\Omega}+e|}{\sqrt{a^{2}+b^{2}+c^{2}}}$ pour le plan $ax+by+cz+e=0$, comparer à $R$, puis, si $d<R$, trouver le centre $H$ du cercle comme intersection du plan avec la droite passant par $\Omega$ et dirigée par le vecteur normal, et conclure par $r=\sqrt{R^{2}-d^{2}}$.
\end{enumerate}
\end{methode}

\courssub{La corde découpée par une droite sécante}

Voici un résultat que le programme ne liste pas explicitement mais qui découle immédiatement de ce que tu sais, et qui tombe régulièrement aux examens.

\begin{propriete}[Longueur de la corde]
Soit $\mathcal{S}$ une sphère de centre $\Omega$ et de rayon $R$, et $\mathcal{D}$ une droite telle que $d=d(\Omega,\mathcal{D})<R$. La droite découpe sur la sphère deux points $A$ et $B$, et le segment $[AB]$ (la \cle{corde}) a pour longueur
\[ AB=2\sqrt{R^{2}-d^{2}}. \]
\end{propriete}

\begin{experience}[Pourquoi cette formule ?]
Soit $H$ le projeté orthogonal de $\Omega$ sur $\mathcal{D}$. Par symétrie de la sphère autour de la droite $(\Omega H)$, $H$ est le milieu de la corde $[AB]$. Dans le triangle $\Omega HA$, rectangle en $H$, le théorème de Pythagore donne :
\[ \Omega A^{2}=\Omega H^{2}+HA^{2}\quad\Longrightarrow\quad R^{2}=d^{2}+HA^{2}\quad\Longrightarrow\quad HA=\sqrt{R^{2}-d^{2}}. \]
Comme $AB=2\,HA$, on obtient $AB=2\sqrt{R^{2}-d^{2}}$. C'est exactement le même Pythagore que pour le cercle-section d'un plan : la géométrie de la sphère est un éternel retour du théorème de Pythagore !
\end{experience}

\begin{exemplebox}[Calcul de corde]
Soit $\mathcal{S}$ la sphère de centre $\Omega(1\,;\,2\,;\,0)$ et de rayon $R=5$, et $\mathcal{D}$ la droite passant par $A(0\,;\,0\,;\,4)$ et dirigée par $\vec{u}(1\,;\,0\,;\,0)$. La distance de $\Omega$ à $\mathcal{D}$ est
\[ d=\dfrac{\|\overrightarrow{A\Omega}\wedge\vec{u}\|}{\|\vec{u}\|}. \]
Ici $\overrightarrow{A\Omega}(1\,;\,2\,;\,-4)$ et $\vec{u}(1\,;\,0\,;\,0)$, donc $\overrightarrow{A\Omega}\wedge\vec{u}=(0\,;\,-4\,;\,-2)$ (à un signe près), de norme $\sqrt{20}=2\sqrt{5}$. Ainsi $d=2\sqrt{5}\approx 4{,}47<R=5$ : la droite est sécante et la corde mesure
\[ 2\sqrt{R^{2}-d^{2}}=2\sqrt{25-20}=2\sqrt{5}. \]
Remarque l'élégance : la corde a la même longueur que la distance $d$ elle-même dans ce cas particulier.
\end{exemplebox}

\courssub{Application 1 : zone de couverture d'un émetteur}

Une antenne relais MTN installée sur la colline de Nkolbisson émet de façon isotrope : son signal remplit la \textbf{boule fermée} de centre $\Omega$ (l'émetteur) et de rayon $R$ (la portée). Un abonné reçoit le signal si et seulement si son récepteur est \emph{à l'intérieur ou sur} la sphère frontière.

\begin{exoresolu}[Le récepteur est-il couvert ?]
Un émetteur est placé au point $\Omega(2\,;\,1\,;\,3)$ (unité : le kilomètre) et a une portée de $R=6$ km. Un village $V$ est situé au point $(8\,;\,5\,;\,6)$. Le village est-il couvert ?
\tcblower
L'équation de la sphère frontière est $(x-2)^{2}+(y-1)^{2}+(z-3)^{2}=36$. Calculons la distance $\Omega V$ :
\[ \Omega V=\sqrt{(8-2)^{2}+(5-1)^{2}+(6-3)^{2}}=\sqrt{36+16+9}=\sqrt{61}\approx 7{,}81. \]
Comme $\Omega V=\sqrt{61}>\sqrt{36}=6=R$, le point $V$ est \textbf{strictement à l'extérieur} de la boule de couverture : le village n'est \textbf{pas couvert}. Il faudrait rapprocher l'antenne d'environ $1{,}81$ km ou augmenter la puissance d'émission.
\end{exoresolu}

\begin{popfigure}
\begin{tikzpicture}[scale=0.9, font=\small]
% Sol (plan horizontal)
\fill[popgreenL, opacity=0.4] (-3.5,0) rectangle (5.5, -0.3);
\draw[popgreen, thick] (-3.5,0) -- (5.5,0);
\node[popgreen!60!black, right] at (5.5, -0.1) {sol (plan horizontal)};
% Sphère de couverture
\shade[ball color=popblue!30, opacity=0.5] (1,2.5) circle (3);
\draw[popblue, thick] (1,2.5) circle (3);
\draw[popblue, dashed] (1,2.5) ellipse (3 and 0.8);
% Centre antenne
\fill[popink] (1,2.5) circle (0.07);
\node[popink, above right] at (1,2.5) {$\Omega$ (antenne)};
% Hauteur d
\draw[poporange, thick, <->] (1,2.5) -- (1,0);
\node[poporange, right] at (1.1,1.25) {$d$};
\fill[poporange] (1,0) circle (0.05);
\node[poporange, below left] at (1,-0.1) {$H$};
% Cercle de couverture au sol
\draw[poppurple, very thick] (1,0) ellipse (1.66 and 0.45);
\draw[poppurple, thick, <->] (1,0) -- (2.66,0);
\node[poppurple, below] at (1.8,-0.15) {$r=\sqrt{R^{2}-d^{2}}$};
% Rayon R
\draw[popblue, thick, ->] (1,2.5) -- (2.8, 4.1);
\node[popblue] at (2.2,3.5) {$R$};
% Récepteurs
\fill[popgreen] (1.5,-0.15) circle (0.06);
\node[popgreen, below] at (1.5,-0.3) {couvert};
\fill[popink] (4.5,-0.15) circle (0.06);
\node[popink, below] at (4.5,-0.3) {hors couverture};
\end{tikzpicture}
\legende{La sphère de portée coupe le plan du sol selon le cercle de couverture, de rayon $r=\sqrt{R^{2}-d^{2}}$, où $d$ est la hauteur de l'antenne au-dessus du sol. Ici $d<R$, l'intersection est bien un cercle.}
\end{popfigure}

\begin{savaistu}[Des antennes relais aux satellites]
La formule $r=\sqrt{R^{2}-d^{2}}$ est utilisée quotidiennement par les ingénieurs télécoms. Pour une antenne relais, elle donne le rayon du cercle de couverture au sol en fonction de la hauteur du pylône. À plus grande échelle, c'est le même calcul qui détermine la « zone de visibilité » d'un satellite géostationnaire au-dessus du Cameroun : la Terre est la sphère, le satellite définit un plan tangent limite, et le cercle de couverture se calcule avec... le théorème de Pythagore, comme en classe de Première !
\end{savaistu}

\courssub{Application 2 : forage d'un puits à travers une colline sphérique}

Une colline est modélisée par une demi-sphère de centre $\Omega$ et de rayon $R$. On fore un puits \emph{rectiligne} (une droite $\mathcal{D}$) qui traverse la colline de part en part. La question pratique : quelle longueur de forage traverse réellement la colline ? C'est exactement la longueur de la corde.

\begin{exoresolu}[Longueur du tunnel de forage]
La colline est modélisée par la sphère de centre $\Omega(0\,;\,0\,;\,0)$ et de rayon $R=10$ (unité : la dizaine de mètres). Le forage suit la droite $\mathcal{D}$ passant par $A(6\,;\,0\,;\,-12)$ et dirigée par $\vec{u}(0\,;\,0\,;\,1)$. Déterminer la longueur de la partie du forage située à l'intérieur de la colline.
\tcblower
\textbf{Étape 1 : distance de $\Omega$ à $\mathcal{D}$.} La droite $\mathcal{D}$ est verticale et passe par le point $(6\,;\,0\,;\,z)$ : sa distance à $\Omega$ est la distance de $\Omega$ à son projeté $H(6\,;\,0\,;\,0)$, soit $d=6$.
\textbf{Étape 2 : comparaison.} $d=6<R=10$ : la droite est sécante, le forage traverse bien la colline.
\textbf{Étape 3 : longueur de la corde.}
\[ AB=2\sqrt{R^{2}-d^{2}}=2\sqrt{100-36}=2\sqrt{64}=16. \]
Le forage traverse la colline sur $16$ unités, soit $\SI{160}{m}$. (Vérification par le calcul : les points d'entrée et de sortie sont $(6\,;\,0\,;\,-8)$ et $(6\,;\,0\,;\,8)$, et l'on retrouve $AB=16$.)
\end{exoresolu}

\courssub{Application 3 : section d'une sphère par un plan (dôme, miroir sphérique)}

Un dôme de salle des fêtes, une calotte de miroir sphérique, un réservoir boule que l'on découpe : dans tous ces cas, on \emph{sectionne} une sphère par un plan et l'ouverture obtenue est un cercle dont il faut connaître le rayon.

\begin{exoresolu}[Rayon de l'ouverture d'un dôme]
Un dôme est une portion de la sphère de centre $\Omega(0\,;\,0\,;\,5)$ et de rayon $R=13$ (unité : le mètre). Le dôme repose sur le plan horizontal $\mathcal{P}$ d'équation $z=0$. Calculer le rayon de l'ouverture circulaire au sol.
\tcblower
Le plan $\mathcal{P}$ a pour équation $z=0$, c'est-à-dire $0x+0y+1z+0=0$. La distance de $\Omega$ à $\mathcal{P}$ est
\[ d=\dfrac{|0+0+5+0|}{\sqrt{0^{2}+0^{2}+1^{2}}}=5. \]
Comme $d=5<R=13$, l'intersection est un cercle. Son centre est le projeté orthogonal de $\Omega$ sur $\mathcal{P}$, soit $H(0\,;\,0\,;\,0)$, et son rayon est
\[ r=\sqrt{R^{2}-d^{2}}=\sqrt{169-25}=\sqrt{144}=12. \]
L'ouverture au sol est donc le cercle de centre $H(0\,;\,0\,;\,0)$ et de rayon $\SI{12}{m}$ : le dôme couvre une emprise au sol de $\pi r^{2}=144\pi\approx \SI{452}{m^{2}}$.
\end{exoresolu}

\begin{attention}[Le cas limite $d=R$ dans les applications]
Si le plan de coupe est tangent ($d=R$), le « cercle » dégénère en un point : $r=\sqrt{R^{2}-R^{2}}=0$. C'est cohérent : la formule $r=\sqrt{R^{2}-d^{2}}$ englobe le cas tangent. En revanche, si $d>R$, la quantité sous la racine est négative : c'est le signal algébrique que l'intersection est vide. La formule ne ment jamais, à condition de vérifier \emph{d'abord} le signe de $R^{2}-d^{2}$.
\end{attention}

\courssub{Bilan des savoir-faire exigibles}

Pour le Probatoire comme pour le Baccalauréat, voici ce que tu dois savoir faire \emph{sans hésitation et sans faute de rédaction} :

\begin{itemize}
\item \textbf{Écrire} l'équation cartésienne d'une sphère à partir de son centre et de son rayon, ou d'un diamètre ;
\item \textbf{Reconnaître} la nature de l'ensemble $x^{2}+y^{2}+z^{2}+ax+by+cz+d=0$ par complétion des carrés, et donner centre et rayon quand c'est une sphère ;
\item \textbf{Déterminer} l'intersection d'une sphère et d'une droite par résolution de l'équation paramétrique (discriminant, puis coordonnées des points) ;
\item \textbf{Déterminer} l'intersection d'une sphère et d'un plan : distance $d$, comparaison à $R$, centre $H$ projeté orthogonal, rayon $r=\sqrt{R^{2}-d^{2}}$ ;
\item \textbf{Rédiger} une conclusion complète : nature de l'intersection \emph{et} éléments caractéristiques ;
\item \textbf{Exploiter} ces résultats dans des situations concrètes : appartenance d'un point à une sphère, tangence, longueur de corde, cercle de section.
\end{itemize}

\begin{exercice}[Entraînement final]
Dans l'espace muni d'un repère orthonormé, on considère l'ensemble $(\mathcal{E})$ d'équation $x^{2}+y^{2}+z^{2}-2x+4y-6z+10=0$, le plan $\mathcal{P}$ d'équation $x+y+z-3=0$ et la droite $\mathcal{D}$ passant par $A(5\,;\,0\,;\,0)$, dirigée par $\vec{u}(0\,;\,1\,;\,1)$.
\begin{enumerate}
\item Montrer que $(\mathcal{E})$ est une sphère dont on précisera le centre $\Omega$ et le rayon $R$.
\item Déterminer l'intersection de $(\mathcal{E})$ et de $\mathcal{P}$.
\item Déterminer l'intersection de $(\mathcal{E})$ et de $\mathcal{D}$, puis la longueur de la corde correspondante.
\end{enumerate}
\end{exercice}

\begin{corrige}[Exercice 1]
\textbf{1.} Complétons les carrés :
\[ (x-1)^{2}-1+(y+2)^{2}-4+(z-3)^{2}-9+10=0\iff (x-1)^{2}+(y+2)^{2}+(z-3)^{2}=4. \]
$(\mathcal{E})$ est la sphère de centre $\Omega(1\,;\,-2\,;\,3)$ et de rayon $R=2$.

\textbf{2.} Distance de $\Omega$ à $\mathcal{P}$ :
\[ d=\dfrac{|1+(-2)+3-3|}{\sqrt{1^{2}+1^{2}+1^{2}}}=\dfrac{1}{\sqrt{3}}=\dfrac{\sqrt{3}}{3}. \]
Comme $d=\frac{\sqrt{3}}{3}<2=R$, l'intersection est un cercle de rayon $r=\sqrt{4-\frac{1}{3}}=\sqrt{\frac{11}{3}}=\frac{\sqrt{33}}{3}$. Le centre $H$ est le projeté orthogonal de $\Omega$ sur $\mathcal{P}$ : la droite $(\Omega H)$ est dirigée par $\vec{n}(1\,;\,1\,;\,1)$, donc $H(1+t\,;\,-2+t\,;\,3+t)$ avec $(1+t)+(-2+t)+(3+t)-3=0$, soit $3t-1=0$, $t=\frac{1}{3}$. D'où $H\left(\frac{4}{3}\,;\,-\frac{5}{3}\,;\,\frac{10}{3}\right)$.

\textbf{3.} Paramétrons $\mathcal{D}$ : $(x\,;\,y\,;\,z)=(5\,;\,t\,;\,t)$. Injectons :
\[ (5-1)^{2}+(t+2)^{2}+(t-3)^{2}=4\iff 16+t^{2}+4t+4+t^{2}-6t+9=4\iff 2t^{2}-2t+25=0. \]
Discriminant : $\Delta=4-200=-196<0$. La droite $\mathcal{D}$ est \textbf{extérieure} à la sphère : intersection vide, il n'y a pas de corde. (Cohérent avec $d(\Omega,\mathcal{D})>R$, ce que l'on peut vérifier.)
\end{corrige}

Tu disposes désormais de tous les outils de la géométrie analytique de la sphère. Retiens l'essentiel : \textbf{un centre, un rayon, une distance, une comparaison} — et le théorème de Pythagore fait le reste. Entraîne-toi sur les exercices de la plateforme : la maîtrise vient de la répétition de la méthode sur des situations variées.

\newpage

\coursec{Activité d'intégration}

\coursec{Sphères}

\courssub{Définition et équation cartésienne d'une sphère}

\begin{definition}[Sphère]
Dans l'espace muni d'un repère orthonormé, une \cle{sphère} $(\mathcal{S})$ de centre $\Omega(x_0\,;y_0\,;z_0)$ et de rayon $R>0$ est l'ensemble des points $M(x\,;y\,;z)$ tels que la distance $\Omega M$ est égale à $R$ :
\[ (\mathcal{S}) = \{\, M \in \mathbb{R}^3 \mid \Omega M = R \,\}. \]
La \cle{boule} fermée associée est l'ensemble des points $M$ tels que $\Omega M \le R$.
\end{definition}

\begin{propriete}[Équation cartésienne d'une sphère]
La sphère de centre $\Omega(x_0\,;y_0\,;z_0)$ et de rayon $R$ admet pour équation cartésienne :
\[ (x-x_0)^2 + (y-y_0)^2 + (z-z_0)^2 = R^2. \]
\end{propriete}

\begin{methode}[Écrire l'équation d'une sphère]
\begin{enumerate}
\item Identifier les coordonnées du centre $\Omega(x_0\,;y_0\,;z_0)$ et la valeur du rayon $R$.
\item Écrire l'équation sous la forme factorisée $(x-x_0)^2+(y-y_0)^2+(z-z_0)^2=R^2$.
\item Si demandé, développer pour obtenir la forme $x^2+y^2+z^2+ax+by+cz+d=0$ avec
$a=-2x_0,\; b=-2y_0,\; c=-2z_0,\; d=x_0^2+y_0^2+z_0^2-R^2$.
\end{enumerate}
\end{methode}

\begin{exemplebox}[Équation d'une sphère]
Soit la sphère de centre $\Omega(1\,;-2\,;3)$ et de rayon $R=4$.
\begin{itemize}
\item Forme factorisée : $(x-1)^2+(y+2)^2+(z-3)^2=16$.
\item Forme développée : $x^2-2x+1 + y^2+4y+4 + z^2-6z+9 = 16 \iff x^2+y^2+z^2-2x+4y-6z-2=0$.
\end{itemize}
\end{exemplebox}

\begin{attention}[Signes dans l'équation factorisée]
L'équation est $(x-x_0)^2+(y-y_0)^2+(z-z_0)^2=R^2$.
\begin{itemize}
\item Si $x_0=2$, le terme est $(x-2)^2$ (et non $(x+2)^2$).
\item Si $y_0=-3$, le terme est $(y-(-3))^2 = (y+3)^2$.
\end{itemize}
Relis toujours les coordonnées du centre sur la forme factorisée, pas sur la forme développée.
\end{attention}

\courssub{Ensemble des points d'équation $x^2+y^2+z^2+ax+by+cz+d=0$}

\begin{propriete}[Équation réduite et nature de l'ensemble]
Tout ensemble $E$ d'équation $x^2+y^2+z^2+ax+by+cz+d=0$ (avec $a,b,c,d \in \mathbb{R}$) peut s'écrire, en complétant les carrés :
\[ \left(x+\frac{a}{2}\right)^2 + \left(y+\frac{b}{2}\right)^2 + \left(z+\frac{c}{2}\right)^2 = \frac{a^2+b^2+c^2}{4} - d. \]
On note $\Delta' = \frac{a^2+b^2+c^2}{4} - d$ (le « discriminant réduit »).
\begin{itemize}
\item Si $\Delta' > 0$ : $E$ est une \cle{sphère} de centre $\Omega\left(-\frac{a}{2}\,;-\frac{b}{2}\,;-\frac{c}{2}\right)$ et de rayon $R=\sqrt{\Delta'}$.
\item Si $\Delta' = 0$ : $E$ est un \cle{point unique} (sphère de rayon nul), le centre $\Omega$.
\item Si $\Delta' < 0$ : $E = \varnothing$ (ensemble vide, aucune solution réelle).
\end{itemize}
\end{propriete}

\begin{methode}[Déterminer la nature et les éléments caractéristiques]
\begin{enumerate}
\item Calculer $\Delta' = \frac{a^2+b^2+c^2}{4} - d$.
\item Comparer $\Delta'$ à $0$.
\item Si $\Delta' \ge 0$, lire le centre $\Omega\left(-\frac{a}{2}\,;-\frac{b}{2}\,;-\frac{c}{2}\right)$ et, si $\Delta'>0$, le rayon $R=\sqrt{\Delta'}$.
\end{enumerate}
\end{methode}

\begin{exemplebox}[Nature d'un ensemble]
Déterminer la nature de $E : x^2+y^2+z^2-4x+6y-2z+14=0$.
$a=-4,\; b=6,\; c=-2,\; d=14$.
$\Delta' = \frac{16+36+4}{4} - 14 = \frac{56}{4} - 14 = 14-14=0$.
$E$ est le point unique $\Omega(2\,;-3\,;1)$.
\end{exemplebox}

\begin{attention}[Pièges classiques]
\begin{itemize}
\item Oublier le facteur $\frac{1}{4}$ dans $\Delta'$ : c'est $\frac{a^2}{4}$, pas $a^2$.
\item Confondre le centre : les coordonnées sont les \emph{opposées} des demi-coefficients ($-\frac{a}{2}$, etc.).
\item Ne pas conclure « ensemble vide » quand $\Delta'<0$ ; écrire « pas de solution » ne suffit pas, il faut préciser la nature géométrique.
\end{itemize}
\end{attention}

\courssub{Position relative d'une sphère et d'une droite}

\begin{propriete}[Intersection sphère -- droite]
Soit $(\mathcal{S})$ une sphère d'équation $(x-x_0)^2+(y-y_0)^2+(z-z_0)^2=R^2$ et $(d)$ une droite d'équations paramétriques
$\begin{cases} x = x_A + t\,u_x \\ y = y_A + t\,u_y \\ z = z_A + t\,u_z \end{cases}$ ($t\in\mathbb{R}$, $\vec{u}(u_x,u_y,u_z)\ne\vec{0}$).
L'intersection $(\mathcal{S})\cap(d)$ se trouve en remplaçant $x,y,z$ par leurs expressions en $t$ dans l'équation de la sphère.
On obtient une équation du second degré en $t$ : $At^2+Bt+C=0$ (avec $A=\|\vec{u}\|^2>0$).
Son discriminant $\Delta = B^2-4AC$ détermine le nombre de points communs :
\begin{itemize}
\item $\Delta > 0$ : deux points d'intersection distincts (la droite est \cle{sécante}).
\item $\Delta = 0$ : un point double (la droite est \cle{tangente} à la sphère).
\item $\Delta < 0$ : aucun point réel (la droite ne rencontre pas la sphère).
\end{itemize}
\end{propriete}

\begin{methode}[Intersection sphère -- droite]
\begin{enumerate}
\item Écrire les équations paramétriques de la droite (si elles ne sont pas données).
\item Substituer $x(t),y(t),z(t)$ dans l'équation de la sphère.
\item Simplifier pour obtenir l'équation $At^2+Bt+C=0$ et calculer $\Delta$.
\item Selon le signe de $\Delta$ :
  \begin{itemize}
  \item $\Delta>0$ : calculer $t_1,t_2$ puis les coordonnées des deux points $M_1,M_2$.
  \item $\Delta=0$ : calculer $t_0$ puis le point de tangence $M$.
  \item $\Delta<0$ : conclure à l'absence d'intersection.
  \end{itemize}
\end{enumerate}
\end{methode}

\begin{exemplebox}[Droite sécante]
Sphère $(\mathcal{S}) : x^2+y^2+z^2=9$ (centre $O$, $R=3$).
Droite $(d) : \begin{cases} x=1+t \\ y=2 \\ z=2-t \end{cases}$.
Substitution : $(1+t)^2 + 4 + (2-t)^2 = 9 \iff 1+2t+t^2+4+4-4t+t^2 = 9 \iff 2t^2-2t=0 \iff 2t(t-1)=0$.
$\Delta=4>0$. $t=0 \Rightarrow M_1(1,2,2)$ ; $t=1 \Rightarrow M_2(2,2,1)$.
La droite coupe la sphère en deux points.
\end{exemplebox}

\courssub{Position relative d'une sphère et d'un plan ; plan tangent}

\begin{propriete}[Intersection sphère -- plan]
Soit $(\mathcal{S})$ une sphère de centre $\Omega$ et de rayon $R$, et $(\mathcal{P})$ un plan.
Soit $d = d(\Omega,(\mathcal{P}))$ la distance du centre au plan.
\begin{itemize}
\item Si $d < R$ : l'intersection est un \cle{cercle} (le plan est \cle{sécant}).
\item Si $d = R$ : l'intersection est un \cle{point unique} ; le plan est \cle{tangent} à la sphère.
\item Si $d > R$ : l'intersection est \cle{vide}.
\end{itemize}
\end{propriete}

\begin{propriete}[Équation du plan tangent]
Soit $(\mathcal{S})$ de centre $\Omega(x_0,y_0,z_0)$ et de rayon $R$, et $M_0(x_M,y_M,z_M)$ un point de $(\mathcal{S})$.
Le plan tangent $(\mathcal{T})$ à $(\mathcal{S})$ en $M_0$ a pour vecteur normal $\overrightarrow{\Omega M_0}$.
Son équation cartésienne est :
\[ (x_M-x_0)(x-x_M) + (y_M-y_0)(y-y_M) + (z_M-z_0)(z-z_M) = 0 \]
ou, sous forme développée :
\[ (x_M-x_0)x + (y_M-y_0)y + (z_M-z_0)z = x_M^2+y_M^2+z_M^2 - (x_0x_M+y_0y_M+z_0z_M). \]
\end{propriete}

\begin{methode}[Position relative sphère -- plan]
\begin{enumerate}
\item Calculer la distance $d$ du centre $\Omega$ au plan $(\mathcal{P})$ (formule de distance point/plan).
\item Comparer $d$ et $R$.
\item Si $d=R$ (tangent) : déterminer le point de tangence $H$ (projeté orthogonal de $\Omega$ sur $(\mathcal{P})$) et, si besoin, l'équation du plan tangent.
\item Si $d<R$ (sécant) : le centre du cercle d'intersection est $H$ (projeté de $\Omega$), son rayon est $r=\sqrt{R^2-d^2}$ (voir section suivante).
\end{enumerate}
\end{methode}

\begin{exemplebox}[Plan tangent]
Sphère $(\mathcal{S}) : (x-1)^2+(y+2)^2+z^2=14$ ($\Omega(1,-2,0), R=\sqrt{14}$).
Plan $(\mathcal{P}) : 2x-y+3z-5=0$.
$d(\Omega,(\mathcal{P})) = \frac{|2(1)-(-2)+3(0)-5|}{\sqrt{4+1+9}} = \frac{|2+2-5|}{\sqrt{14}} = \frac{1}{\sqrt{14}} < \sqrt{14}$.
Le plan est sécant.
\end{exemplebox}

\courssub{Cercle d'intersection d'une sphère et d'un plan}

\begin{propriete}[Caractéristiques du cercle d'intersection]
Lorsque le plan $(\mathcal{P})$ est sécant à la sphère $(\mathcal{S})$ ($d<R$), l'intersection est un cercle $(\mathcal{C})$.
\begin{itemize}
\item Son \cle{centre} $H$ est le projeté orthogonal du centre $\Omega$ de la sphère sur le plan $(\mathcal{P})$.
\item Son \cle{rayon} $r$ est donné par le théorème de Pythagore dans le triangle rectangle $\Omega H M$ ($M \in (\mathcal{C})$) :
\[ r = \sqrt{R^2 - d^2}, \quad \text{où } d = \Omega H = d(\Omega,(\mathcal{P})). \]
\end{itemize}
\end{propriete}

\begin{methode}[Déterminer le cercle d'intersection]
\begin{enumerate}
\item Vérifier que $d<R$ (plan sécant).
\item Déterminer $H$, projeté orthogonal de $\Omega$ sur $(\mathcal{P})$ (résolution du système : $H \in (\mathcal{P})$ et $\overrightarrow{\Omega H}$ colinéaire à un vecteur normal de $(\mathcal{P})$).
\item Calculer $d = \Omega H$ (ou utiliser la distance point/plan déjà trouvée).
\item Calculer $r = \sqrt{R^2-d^2}$.
\item (Optionnel) Donner un système d'équations du cercle : $\begin{cases} \text{équation de }(\mathcal{S}) \\ \text{équation de }(\mathcal{P}) \end{cases}$.
\end{enumerate}
\end{methode}

\begin{exemplebox}[Cercle d'intersection]
$(\mathcal{S}) : x^2+y^2+z^2=25$ ($O,R=5$). $(\mathcal{P}) : z=3$.
$d = d(O,(\mathcal{P})) = 3 < 5$. Sécant.
$H$ projeté de $O$ sur $z=3 \Rightarrow H(0,0,3)$.
$r = \sqrt{5^2-3^2} = 4$.
Le cercle a pour centre $H(0,0,3)$ et rayon $4$. Ses équations : $\begin{cases} x^2+y^2+z^2=25 \\ z=3 \end{cases} \iff \begin{cases} x^2+y^2=16 \\ z=3 \end{cases}$.
\end{exemplebox}

\begin{savaistu}[GPS et trilatération]
Le GPS (Global Positioning System) utilise l'intersection de \textbf{sphères} : chaque satellite émet un signal qui donne la distance au récepteur. Le récepteur se situe sur une sphère centrée sur le satellite. Avec 4 satellites (4 sphères), l'intersection réduit à un point unique (la position), sous réserve de corriger l'horloge du récepteur. Au Cameroun, cette technologie guide les taxis à Douala, les camions sur le corridor Douala--Ndjamena, et les pirogues sur le fleuve Sanaga.
\end{savaistu}

\courssub{Synthèse et applications : Couverture téléphonique au Mont Fébé}

\courssub{Objectifs de l'activité}

À l'issue de cette activité d'intégration, tu devras être capable de :

\begin{itemize}
\item modéliser une situation concrète de la vie quotidienne (couverture d'un émetteur de télécommunication) par une sphère de l'espace ;
\item déterminer l'équation cartésienne d'une sphère connaissant son centre et son rayon ;
\item décider si un point appartient à une boule (intérieur d'une sphère) en comparant une distance au rayon ;
\item déterminer l'intersection d'une sphère et d'un plan, puis préciser la nature et les éléments caractéristiques du cercle obtenu (centre, rayon) ;
\item déterminer l'intersection d'une sphère et d'une droite, et interpréter géométriquement le résultat ;
\item confronter tes résultats à ceux de ton groupe, argumenter et valider collectivement une démarche.
\end{itemize}

\courssub{Situation}

L'entreprise de téléphonie mobile qui dessert la ville de Yaoundé souhaite vérifier la qualité de couverture d'un de ses émetteurs. Cet émetteur est installé au sommet du \cle{Mont Fébé}, qui domine la ville.

On modélise la situation dans un repère orthonormé $(O\,;\vec{i},\vec{j},\vec{k})$ de l'espace, où :
\begin{itemize}
\item l'unité sur chaque axe est le \cle{kilomètre} ;
\item le plan $(xOy)$, d'équation $z=0$, représente le \cle{sol} de la ville ;
\item l'axe $(Oz)$ est vertical, dirigé vers le ciel.
\end{itemize}

L'émetteur est assimilé au point $\Omega(0\,;0\,;4)$ : il est donc situé à \SI{4}{km} d'altitude au-dessus du point $O$ du sol. On admet que l'émetteur couvre tous les points situés à une distance inférieure ou égale à \SI{5}{km} de lui : sa \cle{portée} est de $5$ km. La zone de couverture est donc la boule de centre $\Omega$ et de rayon $5$, et sa frontière est la sphère $(\mathcal{S})$ de centre $\Omega$ et de rayon $5$.

On considère les objets suivants :
\begin{itemize}
\item le quartier d'habitation $A$, dont le centre est modélisé par le point $A(3\,;4\,;0)$ ;
\item une route rectiligne $(d)$, tracée au sol, modélisée par la droite d'équations paramétriques $x=t$, $y=6$, $z=0$ ($t\in\mathbb{R}$) ;
\item un avion de ligne qui survole la région en restant dans le plan horizontal $(\mathcal{P})$ d'équation $z=6$.
\end{itemize}

L'ingénieur de l'opérateur doit répondre aux questions suivantes : le quartier $A$ est-il couvert ? Quelle est l'étendue de la zone couverte au sol ? La route $(d)$ traverse-t-elle la zone couverte, et si oui sur quelle longueur ? L'avion capte-t-il le signal ?

\begin{popfigure}
\begin{center}
\begin{tikzpicture}[scale=0.85]
\draw[popline,fill=popgreenL] (-4.5,0) -- (4.5,0) -- (5.6,1.6) -- (-3.4,1.6) -- cycle;
\node[popdark] at (4.2,1.3) {\small sol $(z=0)$};
\draw[popblue,thick] (0,0.8) circle (2.1);
\draw[popblue,dashed] (0,0.8) ellipse (2.1 and 0.55);
\fill[popink] (0,0.8) circle (1.5pt) node[above left] {$\Omega$};
\draw[popink,thick,->] (0,0.8) -- (0,3.6) node[right] {$z$};
\draw[poporange,very thick] (0,0.8) -- (1.48,2.28) node[midway,above right] {$R=5$};
\fill[popink] (1.26,0.55) circle (1.5pt) node[below] {$A$};
\draw[popgold,very thick] (-3.6,1.05) -- (4.2,1.05) node[right] {$(d)$};
\end{tikzpicture}
\end{center}
\legende{Modélisation : l'émetteur $\Omega$, la sphère de couverture, le sol, le quartier $A$ et la route $(d)$.}
\end{popfigure}

\courssub{Consignes}

Travail en groupes de trois ou quatre élèves. Chaque groupe rédige sa solution complète et justifiée, puis la correction est menée collectivement au tableau, avec une figure en perspective.

\begin{enumerate}
\item Justifier que la sphère $(\mathcal{S})$ de couverture admet pour équation cartésienne :
\[ x^{2}+y^{2}+(z-4)^{2}=25, \]
puis la développer sous la forme $x^{2}+y^{2}+z^{2}+ax+by+cz+d=0$.
\item Le quartier $A(3\,;4\,;0)$ est-il couvert par l'émetteur ? On calculera la distance $\Omega A$ et on la comparera à la portée.
\item On s'intéresse à la zone couverte \emph{au sol}, c'est-à-dire à l'intersection de la sphère $(\mathcal{S})$ avec le plan $(xOy)$ d'équation $z=0$.
\begin{enumerate}
\item Calculer la distance du centre $\Omega$ au plan $(xOy)$.
\item En déduire la nature de l'intersection, puis déterminer son centre $H$ et son rayon $r$.
\item Interpréter concrètement le résultat pour les habitants de la ville.
\end{enumerate}
\item La route $(d)$ d'équations $x=t$, $y=6$, $z=0$ traverse-t-elle la zone couverte ?
\begin{enumerate}
\item Montrer que les points de $(d)$ appartenant à $(\mathcal{S})$ sont donnés par une équation du second degré en $t$.
\item Conclure sur l'intersection de $(d)$ et de la zone couverte, puis calculer la \emph{longueur} de la portion de route couverte.
\end{enumerate}
\item L'avion, volant dans le plan $(\mathcal{P})$ d'équation $z=6$, capte-t-il le signal de l'émetteur ? Justifier par un calcul de distance et une comparaison avec le rayon.
\item Bilan de groupe : rédiger une phrase de conclusion pour l'ingénieur, résumant la situation de la couverture (quartier, zone au sol, route, avion).
\end{enumerate}

\courssub{Ressources à mobiliser}

\begin{aretenir}[Boîte à outils de la leçon]
\begin{itemize}
\item \textbf{Sphère de centre $\Omega(x_0;y_0;z_0)$ et de rayon $R$} : équation cartésienne
\[ (x-x_0)^{2}+(y-y_0)^{2}+(z-z_0)^{2}=R^{2}. \]
Un point $M$ est à l'\emph{intérieur} de la sphère (zone couverte) si et seulement si $\Omega M < R$.
\item \textbf{Distance de deux points} : $\Omega M=\sqrt{(x_M-x_\Omega)^{2}+(y_M-y_\Omega)^{2}+(z_M-z_\Omega)^{2}}$.
\item \textbf{Sphère et plan} : si $d$ est la distance du centre $\Omega$ au plan $(\mathcal{P})$, alors :
\begin{itemize}
\item si $d<R$ : l'intersection est un \cle{cercle} de rayon $r=\sqrt{R^{2}-d^{2}}$, dont le centre est le projeté orthogonal de $\Omega$ sur $(\mathcal{P})$ ;
\item si $d=R$ : le plan est \cle{tangent} à la sphère (un seul point commun) ;
\item si $d>R$ : l'intersection est \cle{vide}.
\end{itemize}
\item \textbf{Sphère et droite} : on injecte les équations paramétriques de la droite dans l'équation de la sphère ; on obtient une équation du second degré dont le discriminant $\Delta$ donne le nombre de points d'intersection ($\Delta>0$ : deux points ; $\Delta=0$ : droite tangente ; $\Delta<0$ : aucun point).
\end{itemize}
\end{aretenir}

\courssub{Démarche et corrigé commenté}

\begin{exoresolu}[Zone de couverture de l'émetteur du Mont Fébé — corrigé complet]
Énoncé : émetteur en $\Omega(0\,;0\,;4)$, portée $R=5$ km ; quartier $A(3\,;4\,;0)$ ; route $(d):x=t,\ y=6,\ z=0$ ; avion dans le plan $z=6$. Répondre aux six consignes.

\tcblower

\textbf{1) Équation de la sphère de couverture.}

La sphère $(\mathcal{S})$ a pour centre $\Omega(0\,;0\,;4)$ et pour rayon $R=5$. Son équation cartésienne est donc :
\[ (x-0)^{2}+(y-0)^{2}+(z-4)^{2}=5^{2} \quad\text{c'est-à-dire}\quad x^{2}+y^{2}+(z-4)^{2}=25. \]
En développant $(z-4)^{2}=z^{2}-8z+16$ :
\[ x^{2}+y^{2}+z^{2}-8z+16=25 \iff x^{2}+y^{2}+z^{2}-8z-9=0. \]
C'est bien la forme $x^{2}+y^{2}+z^{2}+ax+by+cz+d=0$ avec $a=0$, $b=0$, $c=-8$, $d=-9$.

\textbf{2) Le quartier $A$ est-il couvert ?}

On calcule la distance de l'émetteur au quartier :
\[ \Omega A=\sqrt{(3-0)^{2}+(4-0)^{2}+(0-4)^{2}}=\sqrt{9+16+16}=\sqrt{41}. \]
Or $\sqrt{41}\approx 6{,}4 > 5$ : la distance dépasse la portée de l'émetteur.

\textbf{Conclusion :} le quartier $A$ n'est \emph{pas} couvert. L'opérateur devra envisager un relais ou augmenter la puissance de l'émetteur.

\textbf{3) Zone couverte au sol.}

\textbf{a)} Le plan du sol a pour équation $z=0$. La distance de $\Omega(0\,;0\,;4)$ à ce plan est :
\[ d=\lvert z_\Omega - 0\rvert = 4 \ \text{km}. \]

\textbf{b)} On compare : $d=4 < R=5$. L'intersection de $(\mathcal{S})$ et du plan du sol est donc un \cle{cercle} $(\mathcal{C})$.

Son centre est le projeté orthogonal de $\Omega$ sur le plan $z=0$, c'est-à-dire le point $H(0\,;0\,;0)=O$.

Son rayon est donné par le théorème de Pythagore dans le triangle $\Omega HM$ (rectangle en $H$, pour tout point $M$ du cercle) :
\[ r=\sqrt{R^{2}-d^{2}}=\sqrt{25-16}=\sqrt{9}=3 \ \text{km}. \]
Vérification par le calcul : l'intersection vérifie $z=0$ et $x^{2}+y^{2}+16=25$, soit $x^{2}+y^{2}=9$ dans le plan $z=0$ : c'est bien le cercle de centre $O$ et de rayon $3$.

\textbf{c)} Interprétation : au sol, la couverture est un \textbf{disque de \SI{3}{km} de rayon} centré au pied $O$ du Mont Fébé. Seuls les quartiers situés à moins de $3$ km de $O$ captent correctement le signal au niveau du sol.

\textbf{4) La route $(d)$ traverse-t-elle la zone couverte ?}

\textbf{a)} Un point $M(t\,;6\,;0)$ de $(d)$ appartient à $(\mathcal{S})$ si et seulement si :
\begin{align*}
t^{2}+6^{2}+(0-4)^{2}&=25\\
t^{2}+36+16&=25\\
t^{2}&=-27.
\end{align*}

\textbf{b)} Cette équation n'a \textbf{aucune solution réelle} (un carré est toujours positif ou nul). La droite $(d)$ ne rencontre donc \emph{pas} la sphère.

Autre point de vue : tous les points de $(d)$ ont pour ordonnée $y=6$, donc leur distance à $\Omega$ vaut au minimum $\sqrt{0^{2}+6^{2}+4^{2}}=\sqrt{52}\approx 7{,}2 > 5$ (minimum atteint en $t=0$).

\textbf{Conclusion :} la route $(d)$ ne traverse pas la zone couverte ; elle passe « trop loin » de l'émetteur (à plus de $7$ km au point le plus proche). La longueur de la portion couverte est nulle.

\textbf{5) L'avion dans le plan $z=6$ capte-t-il le signal ?}

La distance du centre $\Omega(0\,;0\,;4)$ au plan $(\mathcal{P})$ d'équation $z=6$ est :
\[ d'=\lvert 6-4\rvert = 2 \ \text{km}. \]
Comme $d'=2 < R=5$, le plan $(\mathcal{P})$ coupe la sphère selon un cercle de rayon :
\[ r'=\sqrt{R^{2}-d'^{2}}=\sqrt{25-4}=\sqrt{21}\approx 4{,}6 \ \text{km}, \]
centré au point $H'(0\,;0\,;6)$, projeté orthogonal de $\Omega$ sur $(\mathcal{P})$.

\textbf{Conclusion :} oui, l'avion capte le signal, tant qu'il reste à l'intérieur du disque de centre $H'(0\,;0\,;6)$ et de rayon $\sqrt{21}\approx 4{,}6$ km dans son plan de vol.

\textbf{6) Phrase de bilan pour l'ingénieur.}

« L'émetteur du Mont Fébé couvre au sol un disque de $3$ km de rayon autour du pied de la colline ; le quartier $A$ (à $6{,}4$ km de l'émetteur) et la route $(d)$ (à plus de $7$ km) ne sont pas couverts ; en revanche, les avions volant à $6$ km d'altitude captent le signal dans un rayon d'environ $4{,}6$ km autour de la verticale de l'émetteur. »
\end{exoresolu}

\begin{attention}[Pièges rencontrés pendant l'activité]
\begin{itemize}
\item Ne pas confondre la \emph{distance de $\Omega$ au plan} (ici $4$ km) avec le \emph{rayon du cercle d'intersection} (ici $3$ km) : c'est la relation $r=\sqrt{R^{2}-d^{2}}$ qui les relie.
\item Pour le quartier $A$, il ne suffit pas de regarder les coordonnées au sol : la distance se calcule en \emph{trois dimensions}, avec l'altitude $z=4$ de l'émetteur.
\item Obtenir $t^{2}=-27$ n'est pas une « erreur de calcul » : c'est la preuve algébrique que la droite ne rencontre pas la sphère. Il faut savoir conclure à partir d'une équation sans solution.
\item L'unité est le kilomètre : toutes les longueurs (rayons, distances) s'expriment en km, sans conversion.
\end{itemize}
\end{attention}

\courssub{Bilan de l'activité}

Cette activité a permis de mobiliser, dans un contexte authentique de télécommunication camerounaise, l'ensemble des savoir-faire de la leçon :

\begin{itemize}
\item l'\textbf{écriture de l'équation d'une sphère} à partir de son centre et de son rayon, sous forme factorisée puis développée ;
\item le \textbf{test d'appartenance d'un point à une boule} par comparaison d'une distance au rayon, qui a montré que le quartier $A$ n'est pas couvert ($\sqrt{41}>5$) ;
\item la \textbf{position relative d'une sphère et d'un plan} : la comparaison de $d$ et $R$ a donné un cercle de rayon $3$ km au sol, et un cercle de rayon $\sqrt{21}$ km à l'altitude de l'avion ;
\item la \textbf{position relative d'une sphère et d'une droite} : la résolution d'une équation du second degré (ici sans solution) a prouvé que la route ne traverse pas la zone couverte.
\end{itemize}

Sur le plan des savoir-être, le travail en groupe a développé l'esprit de collaboration, l'argumentation mathématique et l'esprit critique : chaque résultat numérique a dû être \emph{interprété} concrètement (couverture, longueur de route, signal en vol), ce qui est le cœur du métier d'ingénieur. La leçon « Sphères » apparaît ainsi comme un outil puissant de modélisation : dès qu'un phénomène se propage « à égale distance dans toutes les directions » (ondes radio, son, portée d'une alarme, zone de risque), c'est une sphère qui le décrit.

\newpage

\coursec{Méthodes et exercices résolus}

\coursec{Sphères}

\courssub{Sphère : définition et équation cartésienne}

\begin{definition}[Sphère]
Dans l'espace muni d'un repère orthonormé, une \cle{sphère} $(\mathcal{S})$ est l'ensemble des points $M$ tels que la distance $M\Omega$ à un point fixe $\Omega$ (le \cle{centre}) est constante et égale à un réel strictement positif $R$ (le \cle{rayon}).
\[ M \in (\mathcal{S}) \iff \Omega M = R \iff \Omega M^2 = R^2. \]
\end{definition}

\begin{propriete}[Équations cartésiennes d'une sphère]
Soit $\Omega(a\,;\,b\,;\,c)$ le centre et $R>0$ le rayon.
\begin{itemize}
\item \textbf{Forme réduite (canonique) :} $(x-a)^2 + (y-b)^2 + (z-c)^2 = R^2$.
\item \textbf{Forme développée :} $x^2+y^2+z^2 - 2ax - 2by - 2cz + (a^2+b^2+c^2-R^2) = 0$,
soit $x^2+y^2+z^2 + \alpha x + \beta y + \gamma z + \delta = 0$ avec $\alpha=-2a$, $\beta=-2b$, $\gamma=-2c$, $\delta=a^2+b^2+c^2-R^2$.
\end{itemize}
\textbf{Particuliers :}
\begin{itemize}
\item Sphère de diamètre $[AB]$ : centre = milieu de $[AB]$, rayon $=\dfrac{AB}{2}$ ; équation vectorielle $\overrightarrow{MA}\cdot\overrightarrow{MB}=0$.
\item Sphère passant par l'origine $O$ : le terme constant $\delta$ (ou $d$) est nul.
\end{itemize}
\end{propriete}

\begin{methode}[Déterminer l'équation d'une sphère de centre $\Omega$ et de rayon $R$]
\begin{enumerate}
\item Écrire la condition $\Omega M^2 = R^2$ pour $M(x\,;\,y\,;\,z)$.
\item Calculer $\Omega M^2 = (x-a)^2+(y-b)^2+(z-c)^2$.
\item Donner l'équation sous \textbf{forme réduite} : $(x-a)^2+(y-b)^2+(z-c)^2 = R^2$.
\item Si demandé, \textbf{développer} pour obtenir la forme $x^2+y^2+z^2+ax+by+cz+d=0$.
\end{enumerate}
\end{methode}

\begin{exoresolu}[Sphère de centre donné, sphère de diamètre donné]
L'espace est muni d'un repère orthonormé $(O\,;\vec{\imath},\vec{\jmath},\vec{k})$.
\begin{enumerate}
\item Déterminer une équation cartésienne de la sphère $(\mathcal{S})$ de centre $\Omega(2\,;\,-1\,;\,3)$ et de rayon $R=5$. La donner sous forme développée.
\item Déterminer une équation de la sphère $(\mathcal{S}')$ de diamètre $[AB]$ avec $A(1\,;\,0\,;\,2)$ et $B(3\,;\,4\,;\,-2)$.
\end{enumerate}
\tcblower
\textbf{1.} Soit $M(x\,;\,y\,;\,z)$. On a :
\[ M \in (\mathcal{S}) \iff \Omega M^2 = 25 \iff (x-2)^2+(y+1)^2+(z-3)^2=25. \]
Développons chaque carré :
\begin{align*}
(x-2)^2 &= x^2-4x+4,\\
(y+1)^2 &= y^2+2y+1,\\
(z-3)^2 &= z^2-6z+9.
\end{align*}
En sommant : $x^2+y^2+z^2-4x+2y-6z+14 = 25$, d'où :
\[ \boxed{(\mathcal{S}) : x^2+y^2+z^2-4x+2y-6z-11=0.} \]

\textbf{2.} Le centre de $(\mathcal{S}')$ est le milieu $I$ de $[AB]$ :
\[ I\left(\frac{1+3}{2}\,;\,\frac{0+4}{2}\,;\,\frac{2+(-2)}{2}\right) = I(2\,;\,2\,;\,0). \]
Le rayon est $R'=\dfrac{AB}{2}$. Calculons $AB$ :
\[ AB = \sqrt{(3-1)^2+(4-0)^2+(-2-2)^2}=\sqrt{4+16+16}=\sqrt{36}=6, \]
donc $R'=3$. Une équation de $(\mathcal{S}')$ sous forme réduite est :
\[ (x-2)^2+(y-2)^2+z^2=9. \]
\textbf{Vérification par le produit scalaire} : $\overrightarrow{MA}\cdot\overrightarrow{MB}=(1-x)(3-x)+(-y)(4-y)+(2-z)(-2-z)$.
En développant : $x^2-4x+3+y^2-4y+z^2-4=0$, soit $x^2+y^2+z^2-4x-4y-1=0$.
Le développement de $(x-2)^2+(y-2)^2+z^2=9$ donne $x^2+y^2+z^2-4x-4y+8=9$, c'est-à-dire $x^2+y^2+z^2-4x-4y-1=0$. Les deux méthodes concordent.

\textbf{Conclusion :} $(\mathcal{S}) : x^2+y^2+z^2-4x+2y-6z-11=0$ et $(\mathcal{S}') : (x-2)^2+(y-2)^2+z^2=9$.
\end{exoresolu}

\begin{popfigure}
\begin{tikzpicture}[scale=0.9, >=Stealth]
% Axes
\draw[->] (-1,0,0) -- (4,0,0) node[right] {$x$};
\draw[->] (0,-1,0) -- (0,3,0) node[above] {$y$};
\draw[->] (0,0,-1) -- (0,0,4) node[above] {$z$};
% Sphere (wireframe)
\pgfmathsetmacro{\R}{2.5}
\pgfmathsetmacro{\cx}{1.5}
\pgfmathsetmacro{\cy}{1}
\pgfmathsetmacro{\cz}{1.5}
% Circles for wireframe
\draw[popblue, thick] (\cx,\cy,\cz) ++({\R*cos(0)},{\R*sin(0)},0) 
  \foreach \ang in {10,20,...,360} { -- ++({\R*cos(\ang)-\R*cos(\ang-10)},{\R*sin(\ang)-\R*sin(\ang-10)},0) };
\draw[popblue, thick] (\cx,\cy,\cz) ++({\R*cos(0)},0,{\R*sin(0)}) 
  \foreach \ang in {10,20,...,360} { -- ++({\R*cos(\ang)-\R*cos(\ang-10)},0,{\R*sin(\ang)-\R*sin(\ang-10)}) };
\draw[popblue, thick] (\cx,\cy,\cz) ++(0,{\R*cos(0)},{\R*sin(0)}) 
  \foreach \ang in {10,20,...,360} { -- ++(0,{\R*cos(\ang)-\R*cos(\ang-10)},{\R*sin(\ang)-\R*sin(\ang-10)}) };
% Center
\fill[poporange] (\cx,\cy,\cz) circle (2pt) node[below left] {$\Omega(a,b,c)$};
% Point M
\pgfmathsetmacro{\tx}{30}
\pgfmathsetmacro{\ty}{60}
\pgfmathsetmacro{\mx}{\cx+\R*sin(\ty)*cos(\tx)}
\pgfmathsetmacro{\my}{\cy+\R*sin(\ty)*sin(\tx)}
\pgfmathsetmacro{\mz}{\cz+\R*cos(\ty)}
\fill[popgreen] (\mx,\my,\mz) circle (2pt) node[above right] {$M(x,y,z)$};
% Radius
\draw[popred, dashed, thick] (\cx,\cy,\cz) -- (\mx,\my,\mz) node[midway, above, sloped] {$R$};
% Projections for clarity (optional)
\draw[gray, thin, dashed] (\mx,\my,\mz) -- (\mx,\my,0);
\draw[gray, thin, dashed] (\mx,\my,0) -- (\mx,0,0);
\draw[gray, thin, dashed] (\mx,\my,0) -- (0,\my,0);
\end{tikzpicture}
\legende{Sphère de centre $\Omega$ et de rayon $R$ : $\Omega M = R$.}
\end{popfigure}

\courssub{Ensemble d'équation $x^2+y^2+z^2+ax+by+cz+d=0$}

\begin{methode}[Reconnaître une sphère, un point ou l'ensemble vide]
Soit $(E)$ l'ensemble des points $M(x\,;\,y\,;\,z)$ vérifiant $x^2+y^2+z^2+ax+by+cz+d=0$.
\begin{enumerate}
\item \textbf{Compléter les carrés} pour chaque variable (mise sous forme canonique) :
\[ x^2+ax = \left(x+\frac{a}{2}\right)^2-\frac{a^2}{4},\quad
   y^2+by = \left(y+\frac{b}{2}\right)^2-\frac{b^2}{4},\quad
   z^2+cz = \left(z+\frac{c}{2}\right)^2-\frac{c^2}{4}. \]
\item Regrouper les termes constants à droite :
\[ \left(x+\frac{a}{2}\right)^2+\left(y+\frac{b}{2}\right)^2+\left(z+\frac{c}{2}\right)^2 = \rho
\quad\text{avec}\quad \rho = \frac{a^2+b^2+c^2}{4}-d. \]
\item \textbf{Discuter selon le signe de $\rho$} :
\begin{itemize}
\item $\rho > 0$ : $(E)$ est la \cle{sphère} de centre $\Omega\left(-\dfrac{a}{2}\,;\,-\dfrac{b}{2}\,;\,-\dfrac{c}{2}\right)$ et de rayon $R=\sqrt{\rho}$.
\item $\rho = 0$ : $(E)$ est réduit au \cle{point} $\Omega$ (sphère de rayon nul).
\item $\rho < 0$ : $(E)$ est \cle{vide} (une somme de carrés de réels ne peut être strictement négative).
\end{itemize}
\end{enumerate}
\textbf{Réflexe Probatoire :} Ne pas appliquer la formule $R^2=\frac{a^2+b^2+c^2-4d}{4}$ de tête. Refaire la mise sous forme canonique : c'est plus sûr et c'est la démarche attendue.
\end{methode}

\begin{exoresolu}[Nature et éléments caractéristiques d'un ensemble]
Dans chaque cas, déterminer la nature de l'ensemble $(E)$ ; préciser, le cas échéant, le centre et le rayon.
\begin{enumerate}
\item $(E_1) : x^2+y^2+z^2-2x+4y-6z-2=0$ ;
\item $(E_2) : x^2+y^2+z^2+2x-4y+6z+14=0$ ;
\item $(E_3) : x^2+y^2+z^2+2x-4y+6z+15=0$.
\end{enumerate}
\tcblower
\textbf{1. Étude de $(E_1)$.} Mise sous forme canonique :
\begin{align*}
x^2-2x &= (x-1)^2-1,\\
y^2+4y &= (y+2)^2-4,\\
z^2-6z &= (z-3)^2-9.
\end{align*}
L'équation devient :
\[ (x-1)^2+(y+2)^2+(z-3)^2 - 14 - 2 = 0 \iff (x-1)^2+(y+2)^2+(z-3)^2 = 16. \]
Comme $16>0$, $(E_1)$ est la \textbf{sphère} de centre $\Omega_1(1\,;\,-2\,;\,3)$ et de rayon $R_1=\sqrt{16}=4$.

\textbf{2. Étude de $(E_2)$.} De même :
\[ x^2+2x=(x+1)^2-1,\quad y^2-4y=(y-2)^2-4,\quad z^2+6z=(z+3)^2-9. \]
L'équation devient :
\[ (x+1)^2+(y-2)^2+(z+3)^2 - 14 + 14 = 0 \iff (x+1)^2+(y-2)^2+(z+3)^2=0. \]
Une somme de carrés de réels est nulle ssi chaque carré est nul : $x=-1$, $y=2$, $z=-3$.
Donc $(E_2)$ est réduit au \textbf{point unique} $\Omega_2(-1\,;\,2\,;\,-3)$.

\textbf{3. Étude de $(E_3)$.} Avec les mêmes carrés :
\[ (x+1)^2+(y-2)^2+(z+3)^2 - 14 + 15 = 0 \iff (x+1)^2+(y-2)^2+(z+3)^2 = -1. \]
Le premier membre est $\ge 0$, le second est $<0$ : aucune solution. $(E_3)$ est \textbf{l'ensemble vide}.

\textbf{Conclusion :} $(E_1)$ sphère de centre $(1\,;\,-2\,;\,3)$, rayon $4$ ; $(E_2)$ singleton $\{(-1\,;\,2\,;\,-3)\}$ ; $(E_3)=\varnothing$.
\end{exoresolu}

\courssub{Position relative d'une sphère et d'une droite}

\begin{methode}[Intersection sphère-droite par paramétrage]
Pour déterminer l'intersection de la sphère $(\mathcal{S})$ et de la droite $(D)$ passant par $A(x_A\,;\,y_A\,;\,z_A)$ de vecteur directeur $\vec{u}(\alpha\,;\,\beta\,;\,\gamma)$ :
\begin{enumerate}
\item Écrire une \textbf{représentation paramétrique} de $(D)$ :
\[ M\in(D) \iff \begin{cases} x = x_A + \alpha t \\ y = y_A + \beta t \\ z = z_A + \gamma t \end{cases}, \quad t\in\mathbb{R}. \]
\item \textbf{Substituer} $x,y,z$ dans l'équation de $(\mathcal{S})$ (de préférence sous forme réduite). On obtient une équation du second degré en $t$ : $At^2+Bt+C=0$. Calculer son discriminant $\Delta = B^2-4AC$.
\item \textbf{Discussion} :
\begin{itemize}
\item $\Delta < 0$ : $(D)$ ne coupe pas $(\mathcal{S})$ (droite \emph{extérieure}).
\item $\Delta = 0$ : un unique point d'intersection, $(D)$ est \emph{tangente} à $(\mathcal{S})$.
\item $\Delta > 0$ : deux points d'intersection distincts, $(D)$ est \emph{sécante}.
\end{itemize}
\item Calculer les valeurs de $t$, puis les \textbf{coordonnées des points} en reportant dans le paramétrage.
\end{enumerate}
\textbf{Vérification :} Chaque point trouvé doit vérifier l'équation de la sphère.
\end{methode}

\begin{exoresolu}[Intersection sphère-droite]
Soit $(\mathcal{S})$ la sphère d'équation $x^2+y^2+z^2-2x-4y-4=0$ et $(D)$ la droite passant par $A(0\,;\,1\,;\,2)$ de vecteur directeur $\vec{u}(1\,;\,1\,;\,1)$.
\begin{enumerate}
\item Déterminer le centre et le rayon de $(\mathcal{S})$.
\item Déterminer les coordonnées des points d'intersection de $(\mathcal{S})$ et $(D)$.
\end{enumerate}
\tcblower
\textbf{1. Centre et rayon.} Forme canonique :
\[ (x-1)^2-1+(y-2)^2-4+z^2-4=0 \iff (x-1)^2+(y-2)^2+z^2=9. \]
Donc $\Omega(1\,;\,2\,;\,0)$ et $R=3$.

\textbf{2. Intersection.} Paramétrage de $(D)$ :
\[ \begin{cases} x = t \\ y = 1+t \\ z = 2+t \end{cases}, \quad t\in\mathbb{R}. \]
Substitution dans l'équation réduite de $(\mathcal{S})$ :
\[ (t-1)^2+(1+t-2)^2+(2+t)^2=9 \iff (t-1)^2+(t-1)^2+(t+2)^2=9. \]
Développement :
\[ 2(t^2-2t+1)+(t^2+4t+4)=9 \iff 3t^2+6=9 \iff 3t^2=3 \iff t^2=1. \]
$\Delta > 0$ : deux solutions, la droite est \textbf{sécante}.
\begin{itemize}
\item $t=1 \Rightarrow M_1(1\,;\,2\,;\,3)$.
\item $t=-1 \Rightarrow M_2(-1\,;\,0\,;\,1)$.
\end{itemize}
\textbf{Vérification :} $M_1$ : $1+4+9-2-8-4=0$ ; $M_2$ : $1+0+1+2-0-4=0$.

\textbf{Conclusion :} $(\mathcal{S})\cap(D)=\{M_1(1\,;\,2\,;\,3)\,;\,M_2(-1\,;\,0\,;\,1)\}$.
\end{exoresolu}

\begin{popfigure}
\begin{tikzpicture}[scale=0.85, >=Stealth]
% Three cases side by side
% Case 1: Secante
\begin{scope}[xshift=0cm]
\draw[popblue, thick] (0,0) circle (1.5);
\draw[poporange, thick] (-2.5,-0.5) -- (2.5,1.5);
\fill[poporange] (-1.1,-0.2) circle (1.5pt) node[below left] {$M_1$};
\fill[poporange] (1.1,0.6) circle (1.5pt) node[above right] {$M_2$};
\fill[popblue] (0,0) circle (1.5pt) node[below] {$\Omega$};
\node at (0,-2) {\textbf{Sécante} ($\Delta>0$)};
\end{scope}
% Case 2: Tangente
\begin{scope}[xshift=5.5cm]
\draw[popblue, thick] (0,0) circle (1.5);
\draw[poporange, thick] (-2.5,1.5) -- (2.5,1.5);
\fill[poporange] (0,1.5) circle (1.5pt) node[above] {$M$};
\fill[popblue] (0,0) circle (1.5pt) node[below] {$\Omega$};
\draw[popred, dashed] (0,0) -- (0,1.5) node[midway, right] {$R$};
\node at (0,-2) {\textbf{Tangente} ($\Delta=0$)};
\end{scope}
% Case 3: Exterieure
\begin{scope}[xshift=11cm]
\draw[popblue, thick] (0,0) circle (1.5);
\draw[poporange, thick] (-2.5,2) -- (2.5,2);
\fill[popblue] (0,0) circle (1.5pt) node[below] {$\Omega$};
\draw[popred, dashed] (0,0) -- (0,2) node[midway, right] {$d>R$};
\node at (0,-2) {\textbf{Extérieure} ($\Delta<0$)};
\end{scope}
\end{tikzpicture}
\legende{Les trois positions relatives d'une droite et d'une sphère.}
\end{popfigure}

\courssub{Position relative d'une sphère et d'un plan ; plan tangent}

\begin{methode}[Comparer la distance du centre au plan avec le rayon]
Soit $(\mathcal{S})$ de centre $\Omega$ et rayon $R$, et $(\mathcal{P})$ d'équation $ax+by+cz+d=0$.
\begin{enumerate}
\item Calculer la distance de $\Omega$ au plan $(\mathcal{P})$ :
\[ d\big(\Omega,(\mathcal{P})\big) = \frac{|a\,x_\Omega+b\,y_\Omega+c\,z_\Omega+d|}{\sqrt{a^2+b^2+c^2}}. \]
\item \textbf{Comparer} cette distance à $R$ :
\begin{itemize}
\item $d\big(\Omega,(\mathcal{P})\big) > R$ : $(\mathcal{P})$ ne coupe pas $(\mathcal{S})$ (intersection vide).
\item $d\big(\Omega,(\mathcal{P})\big) = R$ : $(\mathcal{P})$ est \cle{tangent} à $(\mathcal{S})$ en $H$, projeté orthogonal de $\Omega$ sur $(\mathcal{P})$.
\item $d\big(\Omega,(\mathcal{P})\big) < R$ : l'intersection est un \cle{cercle}.
\end{itemize}
\end{enumerate}
\textbf{Réflexes :} Valeur absolue au numérateur ; racine carrée au dénominateur ; coordonnées du \emph{centre} $\Omega$.
\end{methode}

\begin{exoresolu}[Discussion selon un paramètre]
Soit $(\mathcal{S})$ de centre $\Omega(1\,;\,-2\,;\,3)$ et rayon $R=3$, et $(\mathcal{P}_m) : 2x-y+2z+m=0$.
\begin{enumerate}
\item Calculer, en fonction de $m$, la distance de $\Omega$ au plan $(\mathcal{P}_m)$.
\item Discuter suivant $m$ la position relative de $(\mathcal{S})$ et $(\mathcal{P}_m)$.
\end{enumerate}
\tcblower
\textbf{1.} $a=2$, $b=-1$, $c=2$ :
\[ d\big(\Omega,(\mathcal{P}_m)\big) = \frac{|2(1)-(-2)+2(3)+m|}{\sqrt{4+1+4}} = \frac{|2+2+6+m|}{3} = \frac{|m+10|}{3}. \]

\textbf{2.} On compare $\dfrac{|m+10|}{3}$ à $R=3$, c'est-à-dire $|m+10|$ à $9$.
\begin{itemize}
\item $|m+10|>9 \iff m+10>9$ ou $m+10<-9 \iff m>-1$ ou $m<-19$ : \textbf{pas d'intersection}.
\item $|m+10|=9 \iff m=-1$ ou $m=-19$ : \textbf{plan tangent}.
\item $|m+10|<9 \iff -19<m<-1$ : \textbf{intersection = cercle}.
\end{itemize}

\textbf{Conclusion :} Pour $m\in\,]{-19}\,;\,{-1}[$, cercle ; pour $m=-19$ ou $m=-1$, tangent ; sinon, vide.
\end{exoresolu}

\begin{methode}[Équation du plan tangent en un point de la sphère]
Pour déterminer le plan $(\mathcal{T})$ tangent en $A$ à la sphère $(\mathcal{S})$ de centre $\Omega$ :
\begin{enumerate}
\item \textbf{Vérifier} que $A\in(\mathcal{S})$ (i.e. $\Omega A = R$).
\item Le plan tangent en $A$ est le plan passant par $A$ et de \textbf{vecteur normal} $\overrightarrow{\Omega A}$.
\item Calculer $\overrightarrow{\Omega A}$, puis écrire : $\overrightarrow{AM}\cdot\overrightarrow{\Omega A}=0$.
\item Développer pour obtenir l'équation cartésienne $ax+by+cz+d=0$.
\end{enumerate}
\textbf{Rappel :} Le point de contact $H$ d'un plan tangent $(\mathcal{P})$ est le projeté orthogonal de $\Omega$ sur $(\mathcal{P})$ (intersection de $(\mathcal{P})$ avec la droite $(\Omega\,;\vec{n})$).
\end{methode}

\begin{exoresolu}[Plan tangent et point de contact]
Soit $(\mathcal{S}) : x^2+y^2+z^2-2x+2y-4z-3=0$.
\begin{enumerate}
\item Centre $\Omega$ et rayon $R$ de $(\mathcal{S})$.
\item Vérifier $A(3\,;\,0\,;\,4)\in(\mathcal{S})$ ; équation du plan tangent $(\mathcal{T})$ en $A$.
\item Plan $(\mathcal{P}) : x+y+z-3=0$. Position relative ? Point de contact si tangent.
\end{enumerate}
\tcblower
\textbf{1.} Forme canonique :
\[ (x-1)^2-1+(y+1)^2-1+(z-2)^2-4-3=0 \iff (x-1)^2+(y+1)^2+(z-2)^2=9. \]
$\Omega(1\,;\,-1\,;\,2)$, $R=3$.

\textbf{2.} $\Omega A^2=(3-1)^2+(0+1)^2+(4-2)^2=4+1+4=9=R^2 \Rightarrow A\in(\mathcal{S})$.
Vecteur normal $\overrightarrow{\Omega A}(2\,;\,1\,;\,2)$.
$\overrightarrow{AM}\cdot\overrightarrow{\Omega A}=0 \iff 2(x-3)+1(y-0)+2(z-4)=0 \iff \boxed{2x+y+2z-14=0}$.

\textbf{3.} Distance $\Omega$ à $(\mathcal{P})$ :
\[ d=\frac{|1+(-1)+2-3|}{\sqrt{3}}=\frac{1}{\sqrt{3}}=\frac{\sqrt{3}}{3}\approx 0,58 < 3=R. \]
$(\mathcal{P})$ n'est \textbf{pas tangent} ; il coupe $(\mathcal{S})$ selon un cercle.
Déterminons le projeté $H$ de $\Omega$ sur $(\mathcal{P})$ (centre du cercle).
Droite $(\Delta)\perp(\mathcal{P})$ par $\Omega$ : $\vec{n}(1\,;\,1\,;\,1)$.
\[ \begin{cases} x=1+t \\ y=-1+t \\ z=2+t \end{cases} \xrightarrow{(\mathcal{P})} (1+t)+(-1+t)+(2+t)-3=0 \iff 3t-1=0 \iff t=\frac{1}{3}. \]
\[ H\left(\frac{4}{3}\,;\,-\frac{2}{3}\,;\,\frac{7}{3}\right). \]

\textbf{Conclusion :} $(\mathcal{T}) : 2x+y+2z-14=0$ ; $(\mathcal{P})$ sécante (cercle), centre du cercle $H\left(\frac{4}{3}\,;\,-\frac{2}{3}\,;\,\frac{7}{3}\right)$.
\end{exoresolu}

\begin{popfigure}
\begin{tikzpicture}[scale=0.9, >=Stealth]
% Sphere wireframe
\pgfmathsetmacro{\R}{2.5}
\pgfmathsetmacro{\cx}{0}
\pgfmathsetmacro{\cy}{0}
\pgfmathsetmacro{\cz}{0}
% Equator
\draw[popblue, thick] (\cx,\cy,\cz) ++({\R*cos(0)},{\R*sin(0)},0) 
  \foreach \ang in {10,20,...,360} { -- ++({\R*cos(\ang)-\R*cos(\ang-10)},{\R*sin(\ang)-\R*sin(\ang-10)},0) };
% Meridians
\draw[popblue, thick] (\cx,\cy,\cz) ++({\R*cos(0)},0,{\R*sin(0)}) 
  \foreach \ang in {10,20,...,360} { -- ++({\R*cos(\ang)-\R*cos(\ang-10)},0,{\R*sin(\ang)-\R*sin(\ang-10)}) };
\draw[popblue, thick] (\cx,\cy,\cz) ++(0,{\R*cos(0)},{\R*sin(0)}) 
  \foreach \ang in {10,20,...,360} { -- ++(0,{\R*cos(\ang)-\R*cos(\ang-10)},{\R*sin(\ang)-\R*sin(\ang-10)}) };
% Center
\fill[poporange] (\cx,\cy,\cz) circle (2pt) node[below left] {$\Omega$};
% Tangent Plane (visual representation as a square)
\pgfmathsetmacro{\px}{2.5} % Point of tangency on x-axis
\fill[poporange!80!black, opacity=0.3] 
  ({\px},-1.5,-1.5) -- ({\px},1.5,-1.5) -- ({\px},1.5,1.5) -- ({\px},-1.5,1.5) -- cycle;
\draw[poporange, thick] 
  ({\px},-1.5,-1.5) -- ({\px},1.5,-1.5) -- ({\px},1.5,1.5) -- ({\px},-1.5,1.5) -- cycle;
% Radius to tangency point
\fill[popgreen] (\px,0,0) circle (2pt) node[right] {$A$};
\draw[popred, thick, dashed] (\cx,\cy,\cz) -- (\px,0,0) node[midway, above] {$R$};
% Normal vector
\draw[->, poppurple, thick] (\px,0,0) -- ++(1,0,0) node[right] {$\vec{n}=\overrightarrow{\Omega A}$};
\end{tikzpicture}
\legende{Plan tangent $(\mathcal{T})$ en $A$ : $\vec{n}=\overrightarrow{\Omega A}$ est normal à $(\mathcal{T})$.}
\end{popfigure}

\courssub{Cercle d'intersection d'une sphère et d'un plan}

\begin{methode}[Centre et rayon du cercle d'intersection]
Lorsque $d\big(\Omega,(\mathcal{P})\big) < R$, l'intersection est un cercle $(\mathcal{C})$.
\begin{enumerate}
\item \textbf{Centre $H$} : projeté orthogonal de $\Omega$ sur $(\mathcal{P})$.
Paramétrer la droite $(\Delta)$ passant par $\Omega$ et dirigée par un vecteur normal $\vec{n}$ de $(\mathcal{P})$, puis résoudre $M(t)\in(\mathcal{P})$.
\item \textbf{Rayon $r$} : Dans le triangle rectangle $\Omega H M$ ($H$ pied de la perpendiculaire), Pythagore donne :
\[ r = \sqrt{R^2 - d\big(\Omega,(\mathcal{P})\big)^2}. \]
\item Conclure : $(\mathcal{C})$ est le cercle de centre $H$, de rayon $r$, \textbf{contenu dans le plan $(\mathcal{P})$}.
\end{enumerate}
\textbf{Cas particulier :} Si le plan passe par $\Omega$ ($d=0$), c'est un \emph{grand cercle} : $H=\Omega$, $r=R$.
\end{methode}

\begin{exoresolu}[Cercle d'intersection — type Probatoire]
$(\mathcal{S}) : x^2+y^2+z^2-2x+4y-6z+5=0$, $(\mathcal{P}) : x-2y+2z-1=0$.
\begin{enumerate}
\item Centre $\Omega$ et rayon $R$ de $(\mathcal{S})$.
\item Montrer que $(\mathcal{P})$ coupe $(\mathcal{S})$ selon un cercle $(\mathcal{C})$.
\item Centre $H$ et rayon $r$ de $(\mathcal{C})$.
\end{enumerate}
\tcblower
\textbf{1.} Forme canonique :
\[ (x-1)^2-1+(y+2)^2-4+(z-3)^2-9+5=0 \iff (x-1)^2+(y+2)^2+(z-3)^2=9. \]
$\Omega(1\,;\,-2\,;\,3)$, $R=3$.

\textbf{2.} Distance de $\Omega$ à $(\mathcal{P})$ ($a=1,b=-2,c=2,d=-1$) :
\[ d = \frac{|1-2(-2)+2(3)-1|}{\sqrt{1+4+4}} = \frac{|1+4+6-1|}{3} = \frac{10}{3} \approx 3,33. \]
$d > R$ : $(\mathcal{P})$ \textbf{ne coupe pas} $(\mathcal{S})$.
\textbf{Correction de l'énoncé :} Travaillons avec $(\mathcal{P}') : x-2y+2z-7=0$ (donnée corrigée pour avoir une intersection).
\[ d' = \frac{|1+4+6-7|}{3} = \frac{4}{3}. \]
$d'=\frac{4}{3}<3=R$ : $(\mathcal{P}')$ coupe $(\mathcal{S})$ selon un cercle $(\mathcal{C})$.

\textbf{3. Centre $H$.} Droite $\perp (\mathcal{P}')$ par $\Omega$, $\vec{n}(1\,;\,-2\,;\,2)$ :
\[ \begin{cases} x=1+t \\ y=-2-2t \\ z=3+2t \end{cases} \xrightarrow{(\mathcal{P}')} (1+t)-2(-2-2t)+2(3+2t)-7=0 \]
\[ \iff 1+t+4+4t+6+4t-7=0 \iff 9t+4=0 \iff t=-\frac{4}{9}. \]
\[ H\left(1-\frac{4}{9}\,;\,-2+\frac{8}{9}\,;\,3-\frac{8}{9}\right) = H\left(\frac{5}{9}\,;\,-\frac{10}{9}\,;\,\frac{19}{9}\right). \]

\textbf{Rayon $r$.} $r=\sqrt{R^2-d'^2}=\sqrt{9-\frac{16}{9}}=\sqrt{\frac{65}{9}}=\frac{\sqrt{65}}{3}$.

\textbf{Conclusion :} Intersection = cercle de centre $H\left(\frac{5}{9}\,;\,-\frac{10}{9}\,;\,\frac{19}{9}\right)$, rayon $r=\frac{\sqrt{65}}{3}$, dans le plan $(\mathcal{P}')$.
\end{exoresolu}

\begin{popfigure}
\begin{tikzpicture}[scale=0.9, >=Stealth]
% Sphere
\pgfmathsetmacro{\R}{2.5}
\pgfmathsetmacro{\cx}{0}
\pgfmathsetmacro{\cy}{0}
\pgfmathsetmacro{\cz}{0}
% Wireframe
\draw[popblue, thick] (\cx,\cy,\cz) ++({\R*cos(0)},{\R*sin(0)},0) 
  \foreach \ang in {10,20,...,360} { -- ++({\R*cos(\ang)-\R*cos(\ang-10)},{\R*sin(\ang)-\R*sin(\ang-10)},0) };
\draw[popblue, thick] (\cx,\cy,\cz) ++({\R*cos(0)},0,{\R*sin(0)}) 
  \foreach \ang in {10,20,...,360} { -- ++({\R*cos(\ang)-\R*cos(\ang-10)},0,{\R*sin(\ang)-\R*sin(\ang-10)}) };
\draw[popblue, thick] (\cx,\cy,\cz) ++(0,{\R*cos(0)},{\R*sin(0)}) 
  \foreach \ang in {10,20,...,360} { -- ++(0,{\R*cos(\ang)-\R*cos(\ang-10)},{\R*sin(\ang)-\R*sin(\ang-10)}) };
% Center
\fill[poporange] (\cx,\cy,\cz) circle (2pt) node[below left] {$\Omega$};
% Cutting Plane (ellipse)
\pgfmathsetmacro{\d}{1.0} % distance from center
\pgfmathsetmacro{\r}{sqrt(\R*\R-\d*\d)} % radius of circle
% Plane visual: tilted ellipse
\draw[popgreen, thick, fill=popgreen!20, opacity=0.5] 
  ({\d}, {-\r}, 0) 
  .. controls ({\d}, {-\r}, {\r}) and ({\d}, {\r}, {\r}) .. ({\d}, {\r}, 0)
  .. controls ({\d}, {\r}, {-\r}) and ({\d}, {-\r}, {-\r}) .. cycle;
% Center H of circle
\fill[popgreen] (\d,0,0) circle (2pt) node[right] {$H$};
% Radius r
\draw[popgreen, thick, dashed] (\d,0,0) -- (\d,\r,0) node[midway, above] {$r$};
% Distance d
\draw[popred, thick, dashed] (\cx,\cy,\cz) -- (\d,0,0) node[midway, below] {$d$};
% Radius R to point on circle
\pgfmathsetmacro{\mx}{\d}
\pgfmathsetmacro{\my}{\r}
\pgfmathsetmacro{\mz}{0}
\draw[poporange, thick, dashed] (\cx,\cy,\cz) -- (\mx,\my,\mz) node[midway, above left] {$R$};
\fill[poporange] (\mx,\my,\mz) circle (1.5pt);
\end{tikzpicture}
\legende{Section d'une sphère par un plan : cercle de centre $H$ (projeté de $\Omega$), rayon $r=\sqrt{R^2-d^2}$.}
\end{popfigure}

\courssub{Problèmes de synthèse — sphère définie par des conditions}

\begin{methode}[Déterminer une sphère à partir de conditions géométriques]
\begin{enumerate}
\item Choisir les \textbf{inconnues} : centre $\Omega(a\,;\,b\,;\,c)$ et rayon $R$ (ou coefficients $a,b,c,d$ de la forme développée).
\item \textbf{Traduire chaque condition} en équation(s) :
\begin{itemize}
\item $A\in(\mathcal{S}) \iff$ coordonnées de $A$ vérifient l'équation.
\item $(\mathcal{P})$ tangent à $(\mathcal{S}) \iff d\big(\Omega,(\mathcal{P})\big)=R$.
\item $\Omega$ sur une droite/plan $\iff$ coordonnées de $\Omega$ vérifient les équations.
\end{itemize}
\item \textbf{Résoudre le système}. Astuce : soustraire les équations « passage par des points » deux à deux élimine les carrés et donne des équations de plans (médiateurs).
\item Conclure : centre, rayon, équation (réduite et/ou développée).
\end{enumerate}
\end{methode}

\begin{exoresolu}[Sphère tangente à un plan, centre sur une droite]
Équation de la sphère $(\mathcal{S})$ de centre $\Omega$ sur $(D) : \begin{cases} x=t \\ y=t \\ z=t \end{cases}$, passant par $A(1\,;\,1\,;\,3)$ et tangente à $(\mathcal{P}) : x+y+z-3=0$.
\tcblower
$\Omega(t\,;\,t\,;\,t)$, rayon $R$.

\textbf{Condition 1 : $A\in(\mathcal{S})$.}
\[ R^2 = \Omega A^2 = (t-1)^2+(t-1)^2+(t-3)^2 = 2(t-1)^2+(t-3)^2. \]

\textbf{Condition 2 : tangence à $(\mathcal{P})$.}
\[ R = d\big(\Omega,(\mathcal{P})\big) = \frac{|t+t+t-3|}{\sqrt{3}} = \frac{|3t-3|}{\sqrt{3}} = \sqrt{3}\,|t-1|, \]
donc $R^2 = 3(t-1)^2$.

\textbf{Résolution.}
\[ 2(t-1)^2+(t-3)^2 = 3(t-1)^2 \iff (t-3)^2 = (t-1)^2. \]
$t-3=t-1$ (impossible) ou $t-3=-(t-1) \iff 2t=4 \iff t=2$.

$\Omega(2\,;\,2\,;\,2)$, $R^2=3(2-1)^2=3 \Rightarrow R=\sqrt{3}$.

\textbf{Vérification :} $\Omega A^2=1+1+1=3=R^2$ ; $d(\Omega,(\mathcal{P}))=\frac{|6-3|}{\sqrt{3}}=\sqrt{3}=R$.

\textbf{Conclusion :} $(\mathcal{S}) : (x-2)^2+(y-2)^2+(z-2)^2=3$, soit développée :
\[ \boxed{x^2+y^2+z^2-4x-4y-4z+9=0.} \]
\end{exoresolu}

\courssub{Synthèse et erreurs à éviter}

\begin{aretenir}[Formules clés à maîtriser]
\begin{itemize}
\item \textbf{Sphère centre $\Omega(a,b,c)$, rayon $R$ :} $(x-a)^2+(y-b)^2+(z-c)^2=R^2$.
\item \textbf{Forme développée :} $x^2+y^2+z^2+ax+by+cz+d=0 \iff$ centre $\left(-\frac{a}{2},-\frac{b}{2},-\frac{c}{2}\right)$, $\rho=\frac{a^2+b^2+c^2}{4}-d$.
  \begin{itemize}
  \item $\rho>0$ : sphère, $R=\sqrt{\rho}$.
  \item $\rho=0$ : point.
  \item $\rho<0$ : vide.
  \end{itemize}
\item \textbf{Distance point-plan :} $d(M,(\mathcal{P}))=\frac{|ax_M+by_M+cz_M+d|}{\sqrt{a^2+b^2+c^2}}$.
\item \textbf{Position sphère-plan :} comparer $d(\Omega,(\mathcal{P}))$ et $R$.
\item \textbf{Cercle d'intersection :} centre $H$ (projeté de $\Omega$), rayon $r=\sqrt{R^2-d^2}$.
\item \textbf{Plan tangent en $A$ :} vecteur normal $\overrightarrow{\Omega A}$, équation $\overrightarrow{AM}\cdot\overrightarrow{\Omega A}=0$.
\end{itemize}
\end{aretenir}

\begin{attention}[Les 5 erreurs qui coûtent des points au Probatoire]
\begin{enumerate}
\item \textbf{Oublier la mise sous forme canonique.} Face à $x^2+y^2+z^2+ax+by+cz+d=0$, annoncer centre/rayon sans écrire $\left(x+\frac{a}{2}\right)^2+\dots = \rho$ est une faute de méthode. Le signe de $\rho$ décide de la nature. Rédige la complétion des carrés explicitement.
\item \textbf{Confondre $R$ et $R^2$.} Si $(x-1)^2+(y+2)^2+z^2=16$, le rayon est $R=4$, pas $16$. Pour le cercle : $r=\sqrt{R^2-d^2}$ (bien mettre les carrés).
\item \textbf{Mal appliquer la formule de distance.} Erreurs classiques : oublier la valeur absolue au numérateur, oublier la racine carrée au dénominateur ($\sqrt{a^2+b^2+c^2}$), ou utiliser un point autre que le centre $\Omega$. Écris la formule en toutes lettres avant de substituer.
\item \textbf{Conclure sans comparer à $R$.} Dire « l'intersection est un cercle » sans justifier $d(\Omega,(\mathcal{P}))<R$ est une affirmation non démontrée. La comparaison distance/rayon \emph{est} la justification. Pour un plan tangent, vérifier $d=R$.
\item \textbf{Donner le centre du cercle égal à $\Omega$.} Le centre du cercle est le \emph{projeté orthogonal} $H$ de $\Omega$ sur le plan, pas $\Omega$ (sauf si le plan passe par $\Omega$). Il faut résoudre le paramétrage de la perpendiculaire ; préciser que le cercle est contenu dans le plan $(\mathcal{P})$.
\end{enumerate}
\end{attention}

\begin{savaistu}[Applications au Cameroun]
La modélisation par des sphères est omniprésente dans les télécommunications (couverture des antennes \textbf{MTN} / \textbf{Orange} : une antenne couvre une zone approximativement sphérique de rayon $R$ selon la puissance), la cartographie (GPS : intersection de sphères centrées sur les satellites pour localiser un récepteur), l'électrification rurale (zone de desserte d'un poste de transformation), ou encore l'agronomie (zone d'irrigation par pivot central vue en 3D pour les cultures en hauteur). La tangence sphère-plan modélise le contact d'un ballon avec le sol ou un mur.
\end{savaistu}

\begin{exercice}[Entraînement — Probatoire]
\begin{enumerate}
\item Soit $(\mathcal{S})$ la sphère d'équation $x^2+y^2+z^2-4x+2y-8z+11=0$. Déterminer son centre et son rayon.
\item La droite $(D)$ passe par $A(0\,;\,1\,;\,-1)$ et a pour vecteur directeur $\vec{u}(2\,;\,-1\,;\,2)$. Déterminer l'intersection de $(D)$ avec la sphère $(\mathcal{S})$ de centre $\Omega(1\,;\,-2\,;\,3)$ et de rayon $R=4$.
\item Soit $(\mathcal{S})$ la sphère de centre $\Omega(2\,;\,1\,;\,-1)$ et de rayon $R=5$. Le plan $(\mathcal{P})$ a pour équation $2x-y+2z-4=0$.
\begin{enumerate}
\item Déterminer la position relative de $(\mathcal{S})$ et $(\mathcal{P})$.
\item Si c'est un cercle, donner son centre $H$ et son rayon $r$.
\end{enumerate}
\item Déterminer l'équation du plan tangent à la sphère $x^2+y^2+z^2-2x+4y-6z-3=0$ au point $A(3\,;\,-1\,;\,4)$.
\item Trouver l'équation de la sphère passant par $A(1\,;\,0\,;\,0)$, $B(0\,;\,1\,;\,0)$, $C(0\,;\,0\,;\,1)$ et dont le centre appartient au plan $x+y+z=0$.
\end{enumerate}
\end{exercice}

\begin{corrige}[Exercice 1]
Forme canonique : $(x-2)^2-4+(y+1)^2-1+(z-4)^2-16+11=0 \iff (x-2)^2+(y+1)^2+(z-4)^2=10$.
Centre $\Omega(2\,;\,-1\,;\,4)$, rayon $R=\sqrt{10}$.
\end{corrige}

\begin{corrige}[Exercice 2]
Paramétrage $(D)$ : $\begin{cases} x=2t \\ y=1-t \\ z=-1+2t \end{cases}$.
$(\mathcal{S})$ : $(x-1)^2+(y+2)^2+(z-3)^2=16$.
Substitution : $(2t-1)^2+(3-t)^2+(2t-4)^2=16$.
$4t^2-4t+1 + t^2-6t+9 + 4t^2-16t+16 = 16 \iff 9t^2-26t+10=0$.
$\Delta = 676-360=316 > 0$ : deux points.
$t_{1,2} = \frac{26\pm\sqrt{316}}{18} = \frac{13\pm\sqrt{79}}{9}$. Coordonnées par report.
\end{corrige}

\begin{corrige}[Exercice 3]
$d(\Omega,(\mathcal{P})) = \frac{|4-1-2-4|}{\sqrt{4+1+4}} = \frac{3}{3}=1$.
$R=5$. $d=1 < 5$ : intersection = cercle.
Centre $H$ : droite $\Omega + t\vec{n}$, $\vec{n}(2,-1,2)$.
$\begin{cases} x=2+2t \\ y=1-t \\ z=-1+2t \end{cases} \xrightarrow{(\mathcal{P})} 2(2+2t)-(1-t)+2(-1+2t)-4=0 \iff 4+4t-1+t-2+4t-4=0 \iff 9t-3=0 \iff t=1/3$.
$H(8/3,\, 2/3,\, -1/3)$.
$r = \sqrt{25-1} = \sqrt{24} = 2\sqrt{6}$.
\end{corrige}

\begin{corrige}[Exercice 4]
$(\mathcal{S})$ : $(x-1)^2+(y+2)^2+(z-3)^2=23$. $\Omega(1,-2,3)$, $R=\sqrt{23}$.
$A(3,-1,4)$ : $\Omega A^2 = 4+1+1=6 \neq 23$. \textbf{Attention :} $A$ n'est pas sur la sphère ! (Vérifier l'énoncé).
Si l'énoncé impose $A$ sur la sphère, il y a une incohérence. En supposant une coquille dans l'énoncé et que $A$ est sur la sphère, la méthode reste : $\vec{n}=\overrightarrow{\Omega A}$, équation $\overrightarrow{AM}\cdot\vec{n}=0$.
\end{corrige}

\begin{corrige}[Exercice 5]
Soit $\Omega(a,b,c)$, $R$. Conditions :
$A\in\mathcal{S} \Rightarrow (a-1)^2+b^2+c^2=R^2$ \quad (1)
$B\in\mathcal{S} \Rightarrow a^2+(b-1)^2+c^2=R^2$ \quad (2)
$C\in\mathcal{S} \Rightarrow a^2+b^2+(c-1)^2=R^2$ \quad (3)
$\Omega \in (P) \Rightarrow a+b+c=0$ \quad (4)

(1)-(2) : $(a-1)^2 - a^2 + b^2 - (b-1)^2 = 0 \iff -2a+1+2b-1=0 \iff b=a$.
(1)-(3) : $(a-1)^2 - a^2 + c^2 - (c-1)^2 = 0 \iff -2a+1+2c-1=0 \iff c=a$.
Avec (4) : $a+a+a=0 \Rightarrow a=0$.
Donc $\Omega(0,0,0)$.
$R^2 = (0-1)^2+0+0 = 1 \Rightarrow R=1$.
Équation : $x^2+y^2+z^2=1$.
\end{corrige}

\newpage

\coursec{Exercices}

\courssub{Je vérifie mes connaissances}

\begin{exercice}[QCM : une seule réponse exacte]
Pour chaque question, choisir la bonne réponse (aucune justification n'est demandée).
\begin{enumerate}
\item La sphère de centre $\Omega(1\,;\,0\,;\,-2)$ et de rayon $3$ a pour équation :
\begin{enumerate}
\item[(a)] $(x-1)^2+y^2+(z+2)^2=3$
\item[(b)] $(x-1)^2+y^2+(z+2)^2=9$
\item[(c)] $(x+1)^2+y^2+(z-2)^2=9$
\end{enumerate}
\item L'ensemble des points $M(x\,;\,y\,;\,z)$ tels que $x^2+y^2+z^2-2x+4y+5=0$ est :
\begin{enumerate}
\item[(a)] une sphère de rayon $1$
\item[(b)] un point
\item[(c)] l'ensemble vide
\end{enumerate}
\item La distance du point $\Omega(0\,;\,0\,;\,0)$ au plan $(P) : 3x-4z+10=0$ vaut :
\begin{enumerate}
\item[(a)] $2$
\item[(b)] $5$
\item[(c)] $10$
\end{enumerate}
\item Si un plan $(P)$ est tangent à une sphère $\mathcal{S}(\Omega\,;\,R)$ en un point $H$, alors :
\begin{enumerate}
\item[(a)] $d(\Omega\,;\,P) > R$
\item[(b)] $(\Omega H)$ est perpendiculaire à $(P)$
\item[(c)] $H$ est le centre de la sphère
\end{enumerate}
\end{enumerate}
\end{exercice}

\begin{exercice}[Vrai ou faux, en justifiant]
L'espace est muni d'un repère orthonormé $(O\,;\,\vec{\imath},\vec{\jmath},\vec{k})$. Dire si chacune des affirmations suivantes est vraie ou fausse, en justifiant la réponse.
\begin{enumerate}
\item L'équation $x^2+y^2+z^2+2x-2y+2z+4=0$ définit une sphère.
\item Le point $A(1\,;\,2\,;\,3)$ appartient à la sphère d'équation $x^2+y^2+z^2=14$.
\item Si $d(\Omega\,;\,P) < R$, alors l'intersection de la sphère $\mathcal{S}(\Omega\,;\,R)$ et du plan $(P)$ est un cercle.
\item Deux points $A$ et $B$ étant donnés, l'ensemble des points $M$ tels que $\overrightarrow{MA}\cdot\overrightarrow{MB}=0$ est la sphère de diamètre $[AB]$.
\end{enumerate}
\end{exercice}

\begin{exercice}[Définitions et vocabulaire]
\begin{enumerate}
\item Donner la définition de la sphère de centre $\Omega$ et de rayon $R$.
\item Qu'appelle-t-on \cle{plan tangent} à une sphère en un point $H$ ? Citer la propriété caractéristique qui le définit.
\item Lorsqu'un plan $(P)$ coupe une sphère $\mathcal{S}(\Omega\,;\,R)$ selon un cercle $(\mathcal{C})$, comment obtient-on le centre et le rayon de $(\mathcal{C})$ ? Faire une figure à main levée.
\item Qu'appelle-t-on \cle{grand cercle} d'une sphère ? Donner une condition sur le plan sécant pour que le cercle d'intersection soit un grand cercle.
\end{enumerate}
\end{exercice}

\begin{exercice}[Calculs directs]
\begin{enumerate}
\item Déterminer une équation cartésienne de la sphère de centre $\Omega(2\,;\,-1\,;\,3)$ et de rayon $R=2$.
\item Déterminer le centre et le rayon de la sphère d'équation $x^2+y^2+z^2-6x+2y-4z+13=0$.
\item Calculer la distance du point $\Omega(1\,;\,2\,;\,-1)$ au plan $(P) : 2x-y+2z-6=0$.
\item Soit $\mathcal{S}$ la sphère de centre $O(0\,;\,0\,;\,0)$ et de rayon $5$. Le point $M(3\,;\,4\,;\,0)$ appartient-il à $\mathcal{S}$ ?
\end{enumerate}
\end{exercice}

\courssub{Je m'entraîne}

\begin{exercice}[Équation d'une sphère et test d'appartenance]
L'espace est muni d'un repère orthonormé $(O\,;\,\vec{\imath},\vec{\jmath},\vec{k})$. On considère les points $\Omega(2\,;\,-1\,;\,3)$, $A(0\,;\,1\,;\,3)$ et $B(4\,;\,-1\,;\,6)$.
\begin{enumerate}
\item Déterminer une équation cartésienne de la sphère $\mathcal{S}$ de centre $\Omega$ passant par $A$.
\item Le point $B$ appartient-il à $\mathcal{S}$ ? Justifier.
\item Déterminer les points d'intersection de $\mathcal{S}$ avec l'axe $(O\,;\,\vec{k})$.
\end{enumerate}
\end{exercice}

\begin{exercice}[Nature d'un ensemble : les trois cas]
Pour chacun des ensembles suivants, déterminer sa nature et, le cas échéant, ses éléments caractéristiques (centre et rayon) :
\begin{align*}
(E_1)&:\ x^2+y^2+z^2-4x+6y-2z+5=0\\
(E_2)&:\ x^2+y^2+z^2-4x+6y-2z+14=0\\
(E_3)&:\ x^2+y^2+z^2-4x+6y-2z+15=0
\end{align*}
Comparer les trois situations et commenter le rôle du terme constant.
\end{exercice}

\begin{exercice}[Intersection d'une sphère et d'une droite]
On considère la sphère $\mathcal{S}$ d'équation $x^2+y^2+z^2-2x+2y-4z-3=0$ et la droite $(D)$ de représentation paramétrique :
\[
\begin{cases}
x=1+2t\\
y=-t\\
z=2+t
\end{cases}\quad (t\in\mathbb{R})
\]
\begin{enumerate}
\item Déterminer le centre $\Omega$ et le rayon $R$ de la sphère $\mathcal{S}$.
\item Déterminer les points d'intersection de $\mathcal{S}$ et de $(D)$.
\item La droite $(D)$ est-elle tangente à $\mathcal{S}$ ? Justifier.
\end{enumerate}
\end{exercice}

\begin{exercice}[Cercle d'intersection et plan tangent]
\begin{enumerate}
\item Soit $\mathcal{S}$ la sphère de centre $\Omega(1\,;\,2\,;\,-1)$ et de rayon $R=4$, et $(P)$ le plan d'équation $2x-y+2z-6=0$. Montrer que l'intersection de $\mathcal{S}$ et $(P)$ est un cercle $(\mathcal{C})$ dont on déterminera le centre $H$ et le rayon $r$.
\item Déterminer une équation du plan tangent à la sphère d'équation $x^2+y^2+z^2=14$ au point $A(1\,;\,2\,;\,3)$, après avoir vérifié que $A$ appartient bien à cette sphère.
\end{enumerate}
\end{exercice}

\courssub{Je me prépare au Probatoire}

\begin{exercice}[La portée d'un radar]
Dans le cadre de la surveillance aérienne de l'aéroport international de Douala, on modélise la zone de détection d'un radar par une sphère $\mathcal{S}$ de centre $\Omega(0\,;\,0\,;\,3)$ et de rayon $R=5$, l'espace étant rapporté à un repère orthonormé $(O\,;\,\vec{\imath},\vec{\jmath},\vec{k})$ dont l'unité est le kilomètre. Le sol est assimilé au plan d'équation $z=0$.
\begin{enumerate}
\item Déterminer une équation cartésienne de la sphère $\mathcal{S}$.
\item 
\begin{enumerate}
\item Calculer la distance du centre $\Omega$ au plan du sol.
\item En déduire que la zone de détection au niveau du sol est un cercle $(\mathcal{C})$, appelé cercle de détection, dont on précisera le centre et le rayon.
\item Un village est situé au point $V(3\,;\,1\,;\,0)$. Ce village est-il couvert par le radar au niveau du sol ?
\end{enumerate}
\item Un avion vole en ligne droite selon la trajectoire $(D)$ de représentation paramétrique :
\[
\begin{cases}
x=t\\
y=t\\
z=3
\end{cases}\quad (t\in\mathbb{R},\ t\text{ en km})
\]
\begin{enumerate}
\item Montrer que la trajectoire de l'avion rencontre la sphère $\mathcal{S}$ en deux points $M_1$ et $M_2$ que l'on déterminera.
\item Calculer la distance $M_1M_2$, c'est-à-dire la longueur du trajet pendant lequel l'avion est détecté par le radar.
\end{enumerate}
\end{enumerate}
\end{exercice}

\begin{exercice}[Sphère de diamètre et plan médiateur]
L'espace est muni d'un repère orthonormé $(O\,;\,\vec{\imath},\vec{\jmath},\vec{k})$. On considère les points $A(1\,;\,-2\,;\,0)$ et $B(3\,;\,0\,;\,4)$.
\begin{enumerate}
\item 
\begin{enumerate}
\item Déterminer les coordonnées du milieu $I$ du segment $[AB]$ et calculer la distance $AB$.
\item En déduire une équation cartésienne de la sphère $\mathcal{S}$ de diamètre $[AB]$.
\item Vérifier que le point $C(3\,;\,-2\,;\,0)$ appartient à $\mathcal{S}$ en utilisant la propriété : $M\in\mathcal{S} \iff \overrightarrow{MA}\cdot\overrightarrow{MB}=0$.
\end{enumerate}
\item 
\begin{enumerate}
\item Déterminer une équation cartésienne du plan médiateur $(P)$ du segment $[AB]$.
\item Montrer que le centre $I$ de la sphère appartient au plan $(P)$.
\item En déduire la nature de l'intersection de $\mathcal{S}$ et de $(P)$, et préciser ses éléments caractéristiques.
\end{enumerate}
\item Déterminer une équation du plan $(Q)$ tangent à la sphère $\mathcal{S}$ au point $A$.
\end{enumerate}
\end{exercice}

\begin{exercice}[Antenne relais et zone de couverture]
Une entreprise de télécommunications installe une antenne relais au sommet d'un pylône dans la région de l'Ouest. On modélise la situation dans un repère orthonormé $(O\,;\,\vec{\imath},\vec{\jmath},\vec{k})$ de l'espace, l'unité étant l'hectomètre (\SI{100}{m}). L'émetteur est placé au point $\Omega(1\,;\,2\,;\,6)$ et sa portée est de \SI{500}{m}. Le sol est assimilé au plan d'équation $z=0$.
\begin{enumerate}
\item Montrer qu'une équation cartésienne de la sphère $\mathcal{S}$ modélisant la zone de couverture est :
\[x^2+y^2+z^2-2x-4y-12z+16=0.\]
\item 
\begin{enumerate}
\item Calculer la distance de $\Omega$ au plan du sol.
\item En déduire que la zone couverte au sol est un disque dont le bord est un cercle $(\mathcal{C})$ dont on déterminera le centre $H$ et le rayon $r$.
\item Un quartier est centré au point $Q(4\,;\,6\,;\,0)$. Son centre est-il couvert par l'antenne ?
\end{enumerate}
\item On considère le plan $(P)$ d'équation $z=6$, horizontal et passant par l'émetteur.
\begin{enumerate}
\item Déterminer la nature et les éléments caractéristiques de l'intersection de $\mathcal{S}$ et de $(P)$.
\item Comment appelle-t-on un tel cercle ?
\end{enumerate}
\item Un technicien se déplace le long de la droite $(\Delta)$ de représentation paramétrique :
\[
\begin{cases}
x=1+t\\
y=2+2t\\
z=6-2t
\end{cases}\quad (t\in\mathbb{R})
\]
Déterminer les points de $(\Delta)$ situés à la limite exacte de la zone de couverture.
\end{enumerate}
\end{exercice}

\newpage

\coursec{Corrigés des exercices}

\begin{corrige}[Exercice 1 — QCM : une seule réponse exacte]
\begin{enumerate}
\item \textbf{Réponse (b).} L'équation de la sphère de centre $\Omega(1\,;\,0\,;\,-2)$ et de rayon $3$ est $(x-1)^2+(y-0)^2+(z-(-2))^2=3^2$, c'est-à-dire $(x-1)^2+y^2+(z+2)^2=9$. Le rayon est \textbf{élevé au carré} : c'est ce qui élimine (a) ; les signes dans les parenthèses sont opposés aux coordonnées du centre : c'est ce qui élimine (c).
\item \textbf{Réponse (b).} On complète les carrés :
\[x^2-2x+y^2+4y+z^2+5=0 \iff (x-1)^2-1+(y+2)^2-4+z^2+5=0 \iff (x-1)^2+(y+2)^2+z^2=0.\]
Une somme de carrés nulle impose que chaque carré soit nul : l'ensemble est réduit au point $(1\,;\,-2\,;\,0)$.
\item \textbf{Réponse (a).} $d(\Omega\,;\,P)=\dfrac{|3\times 0-4\times 0+10|}{\sqrt{3^2+0^2+(-4)^2}}=\dfrac{10}{\sqrt{25}}=\dfrac{10}{5}=2$.
\item \textbf{Réponse (b).} Propriété caractéristique du plan tangent : $(P)$ est tangent à $\mathcal{S}$ en $H$ si et seulement si $(\Omega H)\perp(P)$. On a alors $d(\Omega\,;\,P)=R$ (et non $>R$), et $H$ est un point de la sphère, pas son centre.
\end{enumerate}
\end{corrige}

\begin{corrige}[Exercice 2 — Vrai ou faux, en justifiant]
\begin{enumerate}
\item \textbf{Faux.} Complétons les carrés :
\[x^2+2x+y^2-2y+z^2+2z+4=0 \iff (x+1)^2+(y-1)^2+(z+1)^2=1+1+1-4=-1.\]
Le « rayon au carré » serait $-1<0$ : l'équation définit \textbf{l'ensemble vide}, pas une sphère.
\item \textbf{Vrai.} $1^2+2^2+3^2=1+4+9=14$ : les coordonnées de $A$ vérifient l'équation, donc $A\in\mathcal{S}$.
\item \textbf{Vrai.} C'est le théorème du cours : si $d(\Omega\,;\,P)<R$, l'intersection est un cercle de centre le projeté orthogonal $H$ de $\Omega$ sur $(P)$ et de rayon $r=\sqrt{R^2-d^2}$.
\item \textbf{Vrai.} $\overrightarrow{MA}\cdot\overrightarrow{MB}=0$ signifie que le triangle $MAB$ est rectangle en $M$ (ou $M\in\{A,B\}$), donc $M$ appartient à la sphère de diamètre $[AB]$ (généralisation à l'espace du théorème du triangle rectangle inscrit dans un demi-cercle).
\end{enumerate}
\end{corrige}

\begin{corrige}[Exercice 3 — Définitions et vocabulaire]
\begin{enumerate}
\item La \cle{sphère} de centre $\Omega$ et de rayon $R$ ($R>0$) est l'ensemble des points $M$ de l'espace tels que $\Omega M=R$.
\item Le \cle{plan tangent} à une sphère $\mathcal{S}(\Omega\,;\,R)$ en un point $H$ de la sphère est le plan passant par $H$ et perpendiculaire à la droite $(\Omega H)$. Propriété caractéristique : $(P)$ est tangent à $\mathcal{S}$ en $H$ si et seulement si $(\Omega H)\perp(P)$ ; de manière équivalente, $d(\Omega\,;\,P)=R$.
\item Le centre de $(\mathcal{C})$ est le \textbf{projeté orthogonal} $H$ de $\Omega$ sur $(P)$. Le rayon est $r=\sqrt{R^2-d^2}$ où $d=d(\Omega\,;\,P)$ (théorème de Pythagore dans le triangle $\Omega HM$ rectangle en $H$, pour tout point $M$ du cercle). Figure à main levée : une sphère, un plan sécant, le segment $[\Omega H]$ perpendiculaire au plan, et le cercle de centre $H$.
\item Un \cle{grand cercle} est un cercle d'intersection dont le plan sécant \textbf{passe par le centre} $\Omega$ de la sphère (condition : $d(\Omega\,;\,P)=0$). Son centre est alors $\Omega$ et son rayon est $R$, le rayon de la sphère : c'est un cercle de rayon maximal.
\end{enumerate}
\end{corrige}

\begin{corrige}[Exercice 4 — Calculs directs]
\begin{enumerate}
\item $(x-2)^2+(y+1)^2+(z-3)^2=4$, soit en développant : $x^2+y^2+z^2-4x+2y-6z+10=0$.
\item On complète les carrés :
\[(x-3)^2-9+(y+1)^2-1+(z-2)^2-4+13=0 \iff (x-3)^2+(y+1)^2+(z-2)^2=1.\]
Centre $\Omega(3\,;\,-1\,;\,2)$, rayon $R=1$.
\item $d(\Omega\,;\,P)=\dfrac{|2\times 1-2+2\times(-1)-6|}{\sqrt{4+1+4}}=\dfrac{|2-2-2-6|}{3}=\dfrac{8}{3}$.
\item $OM=\sqrt{3^2+4^2+0^2}=\sqrt{25}=5=R$ : oui, $M\in\mathcal{S}$.
\end{enumerate}
\end{corrige}

\begin{corrige}[Exercice 5 — Équation d'une sphère et test d'appartenance]
\begin{enumerate}
\item Le rayon est $R=\Omega A=\sqrt{(0-2)^2+(1+1)^2+(3-3)^2}=\sqrt{4+4}=\sqrt{8}$. D'où
\[\mathcal{S}:\ (x-2)^2+(y+1)^2+(z-3)^2=8 \iff x^2+y^2+z^2-4x+2y-6z+6=0.\]
\item $\Omega B=\sqrt{(4-2)^2+(-1+1)^2+(6-3)^2}=\sqrt{4+0+9}=\sqrt{13}\neq\sqrt{8}$ : $B\notin\mathcal{S}$.
\item Un point de l'axe $(O\,;\,\vec{k})$ a des coordonnées $(0\,;\,0\,;\,z)$. On injecte dans l'équation :
\[(0-2)^2+(0+1)^2+(z-3)^2=8 \iff (z-3)^2=3 \iff z=3\pm\sqrt{3}.\]
Les points d'intersection sont $K_1(0\,;\,0\,;\,3+\sqrt{3})$ et $K_2(0\,;\,0\,;\,3-\sqrt{3})$.
\end{enumerate}
\end{corrige}

\begin{corrige}[Exercice 6 — Nature d'un ensemble : les trois cas]
Les trois équations ont les mêmes termes en $x$, $y$, $z$ : la mise sous forme canonique donne, pour le membre de gauche,
\[(x-2)^2-4+(y+3)^2-9+(z-1)^2-1 = (x-2)^2+(y+3)^2+(z-1)^2-14.\]
\begin{itemize}
\item $(E_1)$ : $(x-2)^2+(y+3)^2+(z-1)^2=14-5=9$. C'est la \textbf{sphère} de centre $\Omega(2\,;\,-3\,;\,1)$ et de rayon $R=3$.
\item $(E_2)$ : $(x-2)^2+(y+3)^2+(z-1)^2=14-14=0$. Somme de carrés nulle : l'ensemble est réduit au \textbf{point} $\Omega(2\,;\,-3\,;\,1)$ (sphère de rayon nul).
\item $(E_3)$ : $(x-2)^2+(y+3)^2+(z-1)^2=14-15=-1<0$ : \textbf{ensemble vide}.
\end{itemize}
\textbf{Commentaire.} Pour une équation $x^2+y^2+z^2+ax+by+cz+d=0$, la mise sous forme canonique donne $\left(x+\tfrac{a}{2}\right)^2+\left(y+\tfrac{b}{2}\right)^2+\left(z+\tfrac{c}{2}\right)^2=\rho$ avec $\rho=\dfrac{a^2+b^2+c^2}{4}-d$. Le terme constant $d$ détermine seul le signe de $\rho$ : $\rho>0$ (sphère), $\rho=0$ (point), $\rho<0$ (ensemble vide). Ici, faire passer $d$ de $5$ à $15$ fait décroître $\rho$ de $9$ à $-1$.
\end{corrige}

\begin{corrige}[Exercice 7 — Intersection d'une sphère et d'une droite]
\begin{enumerate}
\item Mise sous forme canonique :
\[(x-1)^2-1+(y+1)^2-1+(z-2)^2-4-3=0 \iff (x-1)^2+(y+1)^2+(z-2)^2=9.\]
Centre $\Omega(1\,;\,-1\,;\,2)$, rayon $R=3$.
\item On remplace $x$, $y$, $z$ par leurs expressions en $t$ dans l'équation développée :
\begin{align*}
&(1+2t)^2+(-t)^2+(2+t)^2-2(1+2t)+2(-t)-4(2+t)-3=0\\
\iff\ & (1+4t+4t^2)+t^2+(4+4t+t^2)-2-4t-2t-8-4t-3=0\\
\iff\ & 6t^2-2t-8=0 \iff 3t^2-t-4=0.
\end{align*}
$\Delta=1+48=49$, $\sqrt{\Delta}=7$, d'où $t_1=\dfrac{1-7}{6}=-1$ et $t_2=\dfrac{1+7}{6}=\dfrac{4}{3}$.
\begin{itemize}
\item Pour $t=-1$ : $M_1(-1\,;\,1\,;\,1)$.
\item Pour $t=\dfrac{4}{3}$ : $M_2\left(\dfrac{11}{3}\,;\,-\dfrac{4}{3}\,;\,\dfrac{10}{3}\right)$.
\end{itemize}
Vérification pour $M_1$ : $1+1+1+2+2-4-3=0$. \checkmark
\item \textbf{Non.} L'équation en $t$ admet \textbf{deux} solutions distinctes ($\Delta>0$) : la droite coupe la sphère en deux points, elle est \textbf{sécante}. Elle serait tangente si et seulement si $\Delta=0$ (point double).
\end{enumerate}
\end{corrige}

\begin{corrige}[Exercice 8 — Cercle d'intersection et plan tangent]
\begin{enumerate}
\item Distance du centre au plan :
\[d(\Omega\,;\,P)=\dfrac{|2\times 1-2+2\times(-1)-6|}{\sqrt{2^2+(-1)^2+2^2}}=\dfrac{|2-2-2-6|}{3}=\dfrac{8}{3}.\]
Comme $d=\dfrac{8}{3}<4=R$, l'intersection est un \textbf{cercle} $(\mathcal{C})$ de rayon
\[r=\sqrt{R^2-d^2}=\sqrt{16-\dfrac{64}{9}}=\sqrt{\dfrac{144-64}{9}}=\sqrt{\dfrac{80}{9}}=\dfrac{4\sqrt{5}}{3}.\]
Le centre $H$ est le projeté orthogonal de $\Omega$ sur $(P)$. Un vecteur normal est $\vec{n}(2\,;\,-1\,;\,2)$ ; on écrit $H=\Omega+t\vec{n}$ avec $H\in(P)$ :
\[2(1+2t)-(2-t)+2(-1+2t)-6=0 \iff 9t-8=0 \iff t=\dfrac{8}{9}.\]
D'où $H\left(1+\dfrac{16}{9}\,;\,2-\dfrac{8}{9}\,;\,-1+\dfrac{16}{9}\right)$, soit $H\left(\dfrac{25}{9}\,;\,\dfrac{10}{9}\,;\,\dfrac{7}{9}\right)$.
\item Vérification : $1^2+2^2+3^2=14$, donc $A$ appartient bien à la sphère de centre $O$ et de rayon $\sqrt{14}$. Le plan tangent en $A$ est perpendiculaire à $(OA)$, donc admet $\overrightarrow{OA}(1\,;\,2\,;\,3)$ comme vecteur normal. Pour $M(x\,;\,y\,;\,z)$ :
\[\overrightarrow{AM}\cdot\overrightarrow{OA}=0 \iff (x-1)+2(y-2)+3(z-3)=0 \iff x+2y+3z-14=0.\]
\end{enumerate}
\end{corrige}

\begin{corrige}[Exercice 9 — La portée d'un radar (type Probatoire)]
\begin{enumerate}
\item $\mathcal{S}$ : $(x-0)^2+(y-0)^2+(z-3)^2=25$, soit $x^2+y^2+z^2-6z+9-25=0$, c'est-à-dire
\[x^2+y^2+z^2-6z-16=0.\]
\item 
\begin{enumerate}
\item Le sol a pour équation $z=0$ : $d(\Omega\,;\,\text{sol})=|z_\Omega|=3$ km.
\item $d=3<5=R$ : l'intersection est un cercle $(\mathcal{C})$. Son centre est le projeté orthogonal de $\Omega$ sur le sol : $H(0\,;\,0\,;\,0)=O$. Son rayon : $r=\sqrt{R^2-d^2}=\sqrt{25-9}=\sqrt{16}=4$ km.
\item $OV=\sqrt{3^2+1^2+0^2}=\sqrt{10}\approx 3{,}16<4$ : $V$ est à l'intérieur du cercle de détection, donc le village \textbf{est couvert} par le radar au niveau du sol.
\end{enumerate}
\item 
\begin{enumerate}
\item Un point de $(D)$ est $M(t\,;\,t\,;\,3)$. Il appartient à $\mathcal{S}$ si et seulement si
\[t^2+t^2+9-18-16=0 \iff 2t^2=25 \iff t=\pm\dfrac{5\sqrt{2}}{2}.\]
Deux solutions : la trajectoire rencontre la sphère en
\[M_1\left(\dfrac{5\sqrt{2}}{2}\,;\,\dfrac{5\sqrt{2}}{2}\,;\,3\right) \quad\text{et}\quad M_2\left(-\dfrac{5\sqrt{2}}{2}\,;\,-\dfrac{5\sqrt{2}}{2}\,;\,3\right).\]
\item $M_1M_2=\sqrt{\left(5\sqrt{2}\right)^2+\left(5\sqrt{2}\right)^2+0^2}=\sqrt{50+50}=\sqrt{100}=10$ km. L'avion est détecté pendant un trajet de \SI{10}{km}.
\end{enumerate}
\end{enumerate}
\textbf{Barème indicatif (sur 10)} : 1 pt (équation) ; 0,5 + 1,5 + 1,5 pts (question 2) ; 2,5 + 2 pts (question 3) ; 1 pt (qualité de la rédaction).

\textbf{Points de vigilance.} Citer explicitement la condition $d<R$ avant de conclure à un cercle ; pour le village, comparer $OV$ au rayon $r=4$ du cercle (et non au rayon $R=5$ de la sphère) ; ne pas oublier l'unité (km) dans la réponse finale.
\end{corrige}

\begin{corrige}[Exercice 10 — Sphère de diamètre et plan médiateur (type Probatoire)]
\begin{enumerate}
\item 
\begin{enumerate}
\item $I\left(\dfrac{1+3}{2}\,;\,\dfrac{-2+0}{2}\,;\,\dfrac{0+4}{2}\right)=(2\,;\,-1\,;\,2)$. $AB=\sqrt{(3-1)^2+(0+2)^2+(4-0)^2}=\sqrt{4+4+16}=\sqrt{24}=2\sqrt{6}$.
\item Le rayon est $R=\dfrac{AB}{2}=\sqrt{6}$, donc $R^2=6$ :
\[\mathcal{S}:\ (x-2)^2+(y+1)^2+(z-2)^2=6 \iff x^2+y^2+z^2-4x+2y-4z+3=0.\]
\item $\overrightarrow{CA}=(1-3\,;\,-2+2\,;\,0-0)=(-2\,;\,0\,;\,0)$ et $\overrightarrow{CB}=(3-3\,;\,0+2\,;\,4-0)=(0\,;\,2\,;\,4)$.
\[\overrightarrow{CA}\cdot\overrightarrow{CB}=(-2)\times 0+0\times 2+0\times 4=0,\]
donc $C\in\mathcal{S}$ d'après la caractérisation de la sphère de diamètre $[AB]$.
\end{enumerate}
\item 
\begin{enumerate}
\item Le plan médiateur passe par $I$ et admet $\overrightarrow{AB}(2\,;\,2\,;\,4)$ comme vecteur normal :
\[2(x-2)+2(y+1)+4(z-2)=0 \iff 2x+2y+4z-10=0 \iff x+y+2z-5=0.\]
\item $2+(-1)+2\times 2-5=0$ : les coordonnées de $I$ vérifient l'équation, donc $I\in(P)$.
\item Puisque $I\in(P)$, on a $d(I\,;\,P)=0<R$ : l'intersection est un \textbf{grand cercle} de la sphère, de centre $I(2\,;\,-1\,;\,2)$ et de rayon $r=R=\sqrt{6}$.
\end{enumerate}
\item Le plan tangent en $A$ est perpendiculaire à $(IA)$, avec $\overrightarrow{IA}=(1-2\,;\,-2+1\,;\,0-2)=(-1\,;\,-1\,;\,-2)$ :
\[-(x-1)-(y+2)-2(z-0)=0 \iff x+y+2z+1=0.\]
Vérification : $A(1\,;\,-2\,;\,0)$ donne $1-2+0+1=0$. \checkmark
\end{enumerate}
\textbf{Barème indicatif (sur 10)} : 1 + 1,5 + 1 pts (question 1) ; 1,5 + 0,5 + 1,5 pts (question 2) ; 2 pts (question 3) ; 1 pt (rédaction).

\textbf{Points de vigilance.} Le rayon de la sphère de diamètre $[AB]$ est $\dfrac{AB}{2}$, pas $AB$ ; pour le plan médiateur, justifier par « ensemble des points équidistants de $A$ et $B$ » ou par « plan passant par $I$, de vecteur normal $\overrightarrow{AB}$ » ; pour le plan tangent, citer la propriété $(IA)\perp(Q)$.
\end{corrige}

\begin{corrige}[Exercice 11 — Antenne relais et zone de couverture (type Probatoire)]
La portée est \SI{500}{m} $=5$ hm, donc $R=5$.
\begin{enumerate}
\item $\mathcal{S}$ a pour équation $(x-1)^2+(y-2)^2+(z-6)^2=25$. En développant :
\[x^2-2x+1+y^2-4y+4+z^2-12z+36-25=0 \iff x^2+y^2+z^2-2x-4y-12z+16=0.\]
\item 
\begin{enumerate}
\item $d(\Omega\,;\,\text{sol})=|z_\Omega|=6$ hm.
\item \textbf{Point de vigilance important.} On constate que $d=6>5=R$ : avec les données numériques de l'énoncé (émetteur à $6$ hm $=$ \SI{600}{m} d'altitude, portée \SI{500}{m}), la sphère \textbf{n'atteint pas le sol} : l'intersection avec le plan $z=0$ est \textbf{vide}, et aucun point du sol n'est couvert. L'énoncé visait manifestement une hauteur inférieure à la portée (par exemple $z_\Omega=3$ hm, qui donnerait $H(1\,;\,2\,;\,0)$ et $r=\sqrt{25-9}=4$ hm). La démarche attendue reste : comparer $d$ et $R$, puis, si $d<R$, conclure à un disque de centre le projeté orthogonal $H$ de $\Omega$ sur le sol et de rayon $r=\sqrt{R^2-d^2}$.
\item $\Omega Q=\sqrt{(4-1)^2+(6-2)^2+(0-6)^2}=\sqrt{9+16+36}=\sqrt{61}\approx 7{,}8$ hm $>5$ hm : le centre du quartier n'est \textbf{pas couvert} (cohérent avec le constat du (b)).
\end{enumerate}
\item 
\begin{enumerate}
\item $(P):z=6$ passe par $\Omega(1\,;\,2\,;\,6)$, donc $d(\Omega\,;\,P)=0<5$ : l'intersection est le cercle de centre $\Omega(1\,;\,2\,;\,6)$ et de rayon $r=R=5$ hm, contenu dans le plan $z=6$.
\item C'est un \cle{grand cercle} de la sphère (cercle de rayon maximal, dont le plan passe par le centre).
\end{enumerate}
\item Un point de $(\Delta)$ est $M(1+t\,;\,2+2t\,;\,6-2t)$. Il est à la limite de la zone de couverture s'il appartient à $\mathcal{S}$, c'est-à-dire si $\Omega M=5$ :
\[(t)^2+(2t)^2+(-2t)^2=25 \iff 9t^2=25 \iff t=\pm\dfrac{5}{3}.\]
\begin{itemize}
\item Pour $t=\dfrac{5}{3}$ : $T_1\left(\dfrac{8}{3}\,;\,\dfrac{16}{3}\,;\,\dfrac{8}{3}\right)$.
\item Pour $t=-\dfrac{5}{3}$ : $T_2\left(-\dfrac{2}{3}\,;\,-\dfrac{4}{3}\,;\,\dfrac{28}{3}\right)$.
\end{itemize}
Ce sont les deux points où le technicien entre dans la zone de couverture et en sort.
\end{enumerate}
\textbf{Barème indicatif (sur 10)} : 1,5 pt (question 1) ; 0,5 + 2 + 1 pts (question 2) ; 1 + 0,5 pts (question 3) ; 2,5 pts (question 4) ; 1 pt (rédaction et unités).

\textbf{Points de vigilance.} Convertir les unités avant tout calcul (\SI{500}{m} $=5$ hm) ; toujours comparer $d(\Omega\,;\,P)$ à $R$ \emph{avant} d'annoncer la nature de l'intersection — ici c'est précisément ce réflexe qui permet de détecter le cas $d>R$ ; pour la question 4, utiliser la forme $(x-1)^2+(y-2)^2+(z-6)^2=25$ qui évite tout développement superflu.
\end{corrige}

\newpage

\coursec{Fiche bilan}

\coursec{Sphères : Synthèse et applications}

\begin{popfigure}
\begin{tikzpicture}[mindmap, grow cyclic, every node/.style={concept}, text width=2.6cm, align=center,
root concept/.append style={concept color=popblue, font=\bfseries\large, text=white},
level 1 concept/.append style={level distance=4.6cm, sibling angle=60, font=\bfseries\small},
level 2 concept/.append style={level distance=3.1cm, font=\footnotesize, text width=2.4cm}]
\node[root concept]{Sphères}
  child[concept color=poporange]{ node {Définition et équation}
    child[concept color=poporangeL]{ node {Centre $\Omega$, rayon $R$ : $\Omega M = R$} }
    child[concept color=poporangeL]{ node {$(x-x_\Omega)^2+(y-y_\Omega)^2+(z-z_\Omega)^2=R^2$} }
  }
  child[concept color=popgreen]{ node {Forme développée}
    child[concept color=popgreenL]{ node {$x^2+y^2+z^2+Ax+By+Cz+D=0$} }
    child[concept color=popgreenL]{ node {$\Delta'=\frac{A^2+B^2+C^2}{4}-D$} }
  }
  child[concept color=poppurple]{ node {Sphère et droite}
    child[concept color=poppurpleL]{ node {Distance $d(\Omega,\mathcal{D})$ vs $R$} }
    child[concept color=poppurpleL]{ node {0, 1 ou 2 points d'intersection} }
  }
  child[concept color=popgold]{ node {Sphère et plan}
    child[concept color=popgoldL]{ node {$d(\Omega,\mathcal{P})$ vs $R$} }
    child[concept color=popgoldL]{ node {Plan tangent : $d=R$} }
  }
  child[concept color=poppink]{ node {Cercle d'intersection}
    child[concept color=poppinkL]{ node {Rayon $r=\sqrt{R^2-d^2}$} }
    child[concept color=poppinkL]{ node {Centre $H$ : projeté orthogonal de $\Omega$} }
  }
  child[concept color=popblue]{ node {Méthodes clés}
    child[concept color=popblueL]{ node {Compléter les carrés} }
    child[concept color=popblueL]{ node {Représentation paramétrique d'une droite} }
  };
\end{tikzpicture}
\legende{Carte mentale de la leçon : les six idées forces sur les sphères.}
\end{popfigure}

\begin{definition}[Définitions fondamentales]
\begin{itemize}
  \item \cle{Sphère} de centre $\Omega$ et de rayon $R>0$ : ensemble des points $M$ de l'espace tels que $\Omega M = R$.
  \item \cle{Boule} : ensemble des points $M$ tels que $\Omega M \le R$.
  \item \cle{Plan tangent} à une sphère en un point $A$ : plan passant par $A$ et perpendiculaire à la droite $(\Omega A)$.
  \item \cle{Cercle d'intersection} : lorsque $d(\Omega,\mathcal{P}) < R$, l'intersection de la sphère et du plan $\mathcal{P}$ est un cercle de ce plan.
\end{itemize}
\end{definition}

\begin{propriete}[Formules et théorèmes à connaître par cœur]
\begin{center}
\begin{tabularx}{\linewidth}{|l|X|}\hline
\rowcolor{popblueL} \textbf{Situation} & \textbf{Résultat} \\ \hline
Équation réduite de la sphère $\mathcal{S}(\Omega,R)$ avec $\Omega(x_\Omega;y_\Omega;z_\Omega)$ & $(x-x_\Omega)^2+(y-y_\Omega)^2+(z-z_\Omega)^2=R^2$ \\ \hline
Forme développée & $x^2+y^2+z^2+Ax+By+Cz+D=0$ \\ \hline
Centre et rayon (forme développée) & $\Omega\left(-\dfrac{A}{2};-\dfrac{B}{2};-\dfrac{C}{2}\right)$, \quad $R^2=\dfrac{A^2+B^2+C^2}{4}-D$ \\ \hline
Nature de l'ensemble (forme développée) & $R^2>0$ : sphère ; \quad $R^2=0$ : point $\Omega$ ; \quad $R^2<0$ : ensemble vide \\ \hline
Distance d'un point $\Omega$ à un plan $\mathcal{P}: ux+vy+wz+t=0$ & $d(\Omega,\mathcal{P})=\dfrac{|u x_\Omega+v y_\Omega+w z_\Omega+t|}{\sqrt{u^2+v^2+w^2}}$ \\ \hline
Position relative sphère $\mathcal{S}(\Omega,R)$ et plan $\mathcal{P}$ ($d=d(\Omega,\mathcal{P})$) & $d>R$ : $\varnothing$ ; \quad $d=R$ : un point (plan tangent) ; \quad $d<R$ : cercle \\ \hline
Cercle d'intersection ($\mathcal{S} \cap \mathcal{P}$, $d<R$) & Centre $H$ : projeté orthogonal de $\Omega$ sur $\mathcal{P}$ ; \quad rayon $r=\sqrt{R^2-d^2}$ \\ \hline
Position relative sphère $\mathcal{S}(\Omega,R)$ et droite $\mathcal{D}$ ($\delta=d(\Omega,\mathcal{D})$) & $\delta>R$ : $\varnothing$ ; \quad $\delta=R$ : un point (droite tangente) ; \quad $\delta<R$ : deux points \\ \hline
\end{tabularx}
\end{center}
\end{propriete}

\begin{methode}[Démarches types pour le Probatoire]
\begin{enumerate}
  \item \textbf{Équation d'une sphère} : écrire $\Omega M^2 = R^2$ avec $M(x;y;z)$, puis développer pour obtenir la forme $x^2+y^2+z^2+Ax+By+Cz+D=0$.
  \item \textbf{Nature de $x^2+y^2+z^2+Ax+By+Cz+D=0$} : compléter les carrés $\left(x+\frac{A}{2}\right)^2$, $\left(y+\frac{B}{2}\right)^2$, $\left(z+\frac{C}{2}\right)^2$, puis comparer le membre de droite à $0$.
  \item \textbf{Intersection sphère $\cap$ droite $\mathcal{D}$} :
        \begin{enumerate}
          \item Écrire une représentation paramétrique de $\mathcal{D}$ : $\begin{cases} x = x_0 + t u_x \\ y = y_0 + t u_y \\ z = z_0 + t u_z \end{cases}$.
          \item Substituer dans l'équation de la sphère $\rightarrow$ équation du second degré en $t$ : $at^2+bt+c=0$.
          \item Calculer le discriminant $\Delta = b^2-4ac$ :
                $\Delta > 0$ : deux points ; $\Delta = 0$ : un point (tangente) ; $\Delta < 0$ : aucun point.
          \item Si solutions, calculer les coordonnées des points d'intersection.
        \end{enumerate}
  \item \textbf{Intersection sphère $\cap$ plan $\mathcal{P}$} :
        \begin{enumerate}
          \item Calculer $d = d(\Omega,\mathcal{P})$.
          \item Comparer $d$ et $R$ (conclure : vide, tangent, cercle).
          \item Si $d < R$ (cercle) :
                \begin{itemize}
                  \item Déterminer $H$, projeté orthogonal de $\Omega$ sur $\mathcal{P}$ : $\overrightarrow{\Omega H} = k\,\vec{n}$ ($\vec{n}$ normal à $\mathcal{P}$) et $H \in \mathcal{P}$.
                  \item Calculer le rayon $r = \sqrt{R^2 - d^2}$.
                \end{itemize}
        \end{enumerate}
  \item \textbf{Plan tangent en $A$ à la sphère $\mathcal{S}(\Omega,R)$} :
        \begin{enumerate}
          \item Vérifier qu'$A$ appartient à la sphère ($\Omega A = R$).
          \item $\overrightarrow{\Omega A}$ est un vecteur normal du plan tangent.
          \item Écrire l'équation cartésienne : $(x-x_A)(x_\Omega-x_A) + (y-y_A)(y_\Omega-y_A) + (z-z_A)(z_\Omega-z_A) = 0$ (ou forme développée).
        \end{enumerate}
\end{enumerate}
\end{methode}

\begin{exemplebox}[Exemple type : forme développée $\rightarrow$ centre et rayon]
Soit l'ensemble $E$ d'équation $x^2+y^2+z^2-4x+2y-6z-11=0$.
\begin{enumerate}
  \item Compléter les carrés :
  \[
  \begin{aligned}
  (x^2-4x) + (y^2+2y) + (z^2-6z) - 11 &= 0 \\
  (x-2)^2 - 4 + (y+1)^2 - 1 + (z-3)^2 - 9 - 11 &= 0 \\
  (x-2)^2 + (y+1)^2 + (z-3)^2 &= 25
  \end{aligned}
  \]
  \item Identifier : $\Omega(2;-1;3)$ et $R = 5$.
  \item Conclure : $E$ est la sphère de centre $\Omega(2;-1;3)$ et de rayon $5$.
\end{enumerate}
\end{exemplebox}

\begin{savaistu}[Applications au Cameroun et dans le monde]
\begin{itemize}
  \item \textbf{Géolocalisation (GPS)} : Un récepteur GPS (téléphone, tracker de camion sur le corridor Douala-Bamenda) calcule sa position à l'intersection de plusieurs sphères dont les centres sont les satellites (MTN/Orange/CAMTEL utilisent cette technologie pour la synchronisation réseau).
  \item \textbf{Couverture réseau} : La zone de couverture d'une antenne relais (MTN, Orange, Camtel, Nexttel) se modélise souvent par une sphère (ou un secteur sphérique) de centre l'antenne et de rayon la portée du signal.
  \item \textbf{Réservoirs d'eau} : Les châteaux d'eau sphériques (ex. à Yaoundé, Douala, Bafoussam) stockent l'eau ; leur volume $V=\frac{4}{3}\pi R^3$ et leur surface $S=4\pi R^2$ conditionnent le dimensionnement et le coût de peinture.
  \item \textbf{Modélisation de la Terre} : Pour les grands travaux (barrage de Lom Pangar, pipeline Tchad-Cameroun), on approxime la Terre par une sphère de rayon $\approx 6371$ km pour calculer des distances orthodromiques.
\end{itemize}
\end{savaistu}

\begin{aretenir}[Checklist avant le Probatoire]
\begin{itemize}
  \item[$\square$] Je sais écrire l'équation réduite d'une sphère connaissant son centre $\Omega(x_\Omega;y_\Omega;z_\Omega)$ et son rayon $R$.
  \item[$\square$] Je sais passer de la forme réduite à la forme développée $x^2+y^2+z^2+Ax+By+Cz+D=0$ (et inversement par complétion des carrés).
  \item[$\square$] Je sais déterminer le centre $\Omega\left(-\frac{A}{2};-\frac{B}{2};-\frac{C}{2}\right)$ et le rayon $R=\sqrt{\frac{A^2+B^2+C^2}{4}-D}$.
  \item[$\square$] Je sais conclure sur la nature de l'ensemble : sphère ($R^2>0$), point réduit à $\Omega$ ($R^2=0$), ou ensemble vide ($R^2<0$).
  \item[$\square$] Je sais calculer la distance d'un point à un plan avec la formule $\dfrac{|u x_\Omega+v y_\Omega+w z_\Omega+t|}{\sqrt{u^2+v^2+w^2}}$.
  \item[$\square$] Je sais discuter l'intersection sphère--plan en comparant $d(\Omega,\mathcal{P})$ et $R$ (trois cas distincts).
  \item[$\square$] Je sais que le plan tangent vérifie $d(\Omega,\mathcal{P})=R$ et que $\overrightarrow{\Omega A}$ en est un vecteur normal.
  \item[$\square$] Je sais déterminer le centre $H$ (projeté orthogonal de $\Omega$ sur $\mathcal{P}$) et le rayon $r=\sqrt{R^2-d^2}$ du cercle d'intersection.
  \item[$\square$] Je sais résoudre l'intersection sphère--droite par substitution paramétrique et discussion sur le discriminant $\Delta$.
  \item[$\square$] Je pense à vérifier qu'un point appartient à la sphère ($\Omega A = R$) avant de parler de tangence.
  \item[$\square$] Je rédige proprement : je cite la condition ($d<R$, $d=R$ ou $d>R$ ; $\Delta>0$, $\Delta=0$, $\Delta<0$) avant de conclure.
\end{itemize}
\end{aretenir}

\findecours

\end{document}
